% Formulation des idées essentielles d’un texte à résumer — Français Tle Tech — Cours signé OnBuch+
% Programme officiel MINESEC. Compiler avec : tectonic N21-formulation-idees-essentielles-texte-resumer.tex
\newcommand{\DOCMATIERE}{Français}
\newcommand{\DOCNIVEAU}{Tle Tech}
\newcommand{\DOCMODULE}{Français – classe de Terminale technique}
\newcommand{\DOCLECON}{N21}
\newcommand{\DOCTITRE}{Formulation des idées essentielles d’un texte à résumer}
% OnBuch+ — préambule « cours signés » ENRICHI (compilé avec tectonic / XeLaTeX)
% Système visuel « Pop » d'OnBuch+ : fond crème (jamais blanc), bordures encre
% épaisses, ombres « dures » décalées, titres Archivo Black en capitales.
% Version MATHÉMATIQUES : graphes (pgfplots/tikz), boîtes pédagogiques
% (définition, propriété, méthode, attention, le savais-tu, exercice résolu,
% exercice, TP, bilan), page de garde et signature OnBuch+ ancrée en bas.
%
% Macros attendues dans main.tex AVANT \input{preamble} :
%   \DOCMATIERE  \DOCNIVEAU  \DOCMODULE  \DOCLECON (numéro)  \DOCTITRE
\documentclass[11pt]{article}

\usepackage{fontspec}
\usepackage[french]{babel}
\usepackage[a4paper,margin=2cm,top=3.4cm,bottom=2.8cm,headheight=1.4cm,footskip=1.3cm]{geometry}
\usepackage{amsmath,amssymb,amsfonts}
\usepackage{mathtools} % \xmapsto, \coloneqq... souvent utilisés par les modèles
\usepackage{cancel} % simplifications barrées (bilans de forces)
% \abs{z} (module, valeur absolue) et \norm{u} (norme), utilisés par les modèles
\providecommand{\abs}{}\let\abs\relax\DeclarePairedDelimiter{\abs}{\lvert}{\rvert}
\providecommand{\norm}{}\let\norm\relax\DeclarePairedDelimiter{\norm}{\lVert}{\rVert}
\usepackage[table]{xcolor} % [table] : fournit aussi \rowcolors (alternance de lignes)
\usepackage{pagecolor}
\usepackage{tikz}
\usetikzlibrary{arrows.meta,positioning,calc,shapes.geometric,shapes.misc,shapes.arrows,decorations.pathmorphing,decorations.pathreplacing,decorations.markings,patterns,shadows,mindmap,trees,fit,backgrounds,matrix,intersections,angles,quotes}
\usepackage{pgfplots}
\usepackage[version=4]{mhchem} % équations nucléaires \ce{^{235}_{92}U}
\usepackage[european,siunitx]{circuitikz} % schémas électriques
\pgfplotsset{compat=1.18}
\usepgfplotslibrary{fillbetween,groupplots} % aires entre courbes (calcul intégral)
\usepackage{enumitem}
\usepackage{fancyhdr}
% [most] charge toutes les bibliothèques tcolorbox (listings, minted, raster,
% documentation...) et gonfle inutilement la mémoire TeX sur un document long
% (~300 boîtes) : on ne charge que ce qui est réellement utilisé.
\usepackage[skins,breakable]{tcolorbox}
\usepackage{graphicx}
\usepackage{colortbl}
\usepackage{booktabs}
\usepackage{tabularx}
\usepackage{multirow}
\usepackage{diagbox} % case d'en-tête diagonale des tableaux à double entrée
% Colonnes extensibles centrée / gauche / droite (C, L, R), souvent utilisées par les modèles
\newcolumntype{C}{>{\centering\arraybackslash}X}
\newcolumntype{L}{>{\raggedright\arraybackslash}X}
\newcolumntype{R}{>{\raggedleft\arraybackslash}X}
\usepackage{array}
\usepackage{multicol}
\usepackage{siunitx}
% parse-numbers=false : accepte aussi \SI{2,5 \times 10^{-3}}{...}
\sisetup{locale=FR,output-decimal-marker={,},group-minimum-digits=4,parse-numbers=false}
\DeclareSIUnit\ohm{\ensuremath{\symup{\Omega}}} % U+2126 absent de la police
\DeclareSIUnit\division{div} % réglages d'oscilloscope (ms/div, V/div)
\DeclareSIUnit\tr{tr} % tours (vitesse de rotation, ex. \SI{1500}{\tr\per\minute})
\DeclareSIUnit\molar{M} % concentration molaire (ex. \SI{3}{\molar}, \SI{1}{\micro\molar})
\DeclareSIUnit\an{an} % année (le modèle confond parfois avec \year, un compteur TeX) — ex. \SI{5}{\kilo\meter\per\an}
\DeclareSIUnit\FCFA{FCFA} % franc CFA (ex. \SI{500}{\FCFA\per\kilo\gram})
\DeclareSIUnit\degreeBrix{\textdegree Brix} % teneur en sucre d'un sirop/jus (réfractométrie)
\DeclareSIUnit\month{mois} % le modèle utilise parfois \month, un compteur TeX, comme unité de durée
\DeclareSIUnit\microliter{\micro\liter} % le modèle écrit parfois \microliter en un seul token
\DeclareSIUnit\million{millions} % dénombrement (ex. \SI{15}{\million\per\mL} spermatozoïdes)
\usepackage{qrcode}
% \degree : commande fréquemment utilisée par les modèles pour un angle ou une
% température en dehors de \SI{}{} (non fournie par les paquets chargés).
\providecommand{\degree}{^{\circ}}
% \qed (fin de démonstration), utilisé par les modèles sans amsthm
\providecommand{\qed}{\hfill$\square$}
% \barre : double barre des valeurs interdites dans un tableau de variations (tkz-tab)
\providecommand{\barre}{\ensuremath{\|}}
% environnement proof (amsthm non chargé) utilisé par les modèles
\ifdefined\proof\else\newenvironment{proof}[1][Démonstration]{\par\noindent\textit{#1.}\ }{\qed\par}\fi
\usepackage{lastpage}
\usepackage{hyperref}
\hypersetup{hidelinks}

% ============================================================
% Palette « Pop » OnBuch+ — mêmes hex que la classe P (plus_ui.dart)
% ============================================================
\definecolor{poporange}{HTML}{F59321}
\definecolor{popdark}{HTML}{E07A0C}
\definecolor{popcream}{HTML}{FFF6E8}
\definecolor{popink}{HTML}{1C1714}
\definecolor{poppurple}{HTML}{7A5AE0}
\definecolor{popgreen}{HTML}{2E9E6A}
\definecolor{poppink}{HTML}{E0457B}
\definecolor{popblue}{HTML}{2F7DE1}
\definecolor{popgold}{HTML}{C98A12}
\definecolor{popmuted}{HTML}{8A7F76}
\definecolor{popline}{HTML}{EADFD2}
\definecolor{popgray}{HTML}{8A7F76}
\colorlet{popgrey}{popgray}
% Alias tolérants (le modèle nomme parfois les couleurs différemment)
\colorlet{popwhite}{white}
\colorlet{popblack}{popink}
\colorlet{popred}{poppink}
\colorlet{popyellow}{popgold}
% Teintes claires (fonds de tableaux/encadrés) — le modèle nomme parfois
% les couleurs avec un suffixe L (ex. popblueL) sans les avoir définies.
\colorlet{poporangeL}{poporange!20}
\colorlet{popdarkL}{popdark!20}
\colorlet{poppurpleL}{poppurple!20}
\colorlet{popgreenL}{popgreen!20}
\colorlet{poppinkL}{poppink!20}
\colorlet{popblueL}{popblue!20}
\colorlet{popgoldL}{popgold!20}
\colorlet{popredL}{popred!20}
\colorlet{popgrayL}{popgray!20}
% Teintes claires pour fonds de tableaux / zones
\colorlet{popblueL}{popblue!10!white}
\colorlet{poporangeL}{poporange!14!white}
\colorlet{popgreenL}{popgreen!12!white}
\colorlet{poppurpleL}{poppurple!12!white}
\colorlet{poppinkL}{poppink!12!white}
\colorlet{popgoldL}{popgold!14!white}

\pagecolor{popcream}

% ============================================================
% Typographie
% ============================================================
\defaultfontfeatures{Ligatures=TeX}
\setmainfont{PlusJakartaSans-Medium.ttf}[Path=../../../fonts/,BoldFont=PlusJakartaSans-Bold.ttf]
% Maths en sans-serif (Fira Math) avec chiffres et lettres droites de Plus
% Jakarta, pour que les formules se fondent dans le texte.
\usepackage[math-style=ISO,bold-style=ISO]{unicode-math}
\setmathfont{FiraMath-Regular.otf}[Path=../../../fonts/]
\setmathfont{PlusJakartaSans-Medium.ttf}[Path=../../../fonts/,range={up/num,up/{Latin,latin},\mathup}]
\usepackage{icomma} % virgule décimale sans espace en mode math
% Caractères mathématiques Unicode absents des polices de texte (Plus Jakarta,
% Archivo Black) : rendus par leur équivalent mathématique, en texte comme en
% formule — sinon ils s'affichent en carré vide (ex. « Arithmétique dans ℤ »).
\usepackage{newunicodechar}
\newunicodechar{ℝ}{\ensuremath{\mathbb{R}}}
\newunicodechar{ℤ}{\ensuremath{\mathbb{Z}}}
\newunicodechar{ℂ}{\ensuremath{\mathbb{C}}}
\newunicodechar{ℕ}{\ensuremath{\mathbb{N}}}
\newunicodechar{ℚ}{\ensuremath{\mathbb{Q}}}
\newunicodechar{𝔻}{\ensuremath{\mathbb{D}}}
\newunicodechar{α}{\ensuremath{\alpha}}
\newunicodechar{β}{\ensuremath{\beta}}
\newunicodechar{γ}{\ensuremath{\gamma}}
\newunicodechar{δ}{\ensuremath{\delta}}
\newunicodechar{ε}{\ensuremath{\varepsilon}}
\newunicodechar{θ}{\ensuremath{\theta}}
\newunicodechar{λ}{\ensuremath{\lambda}}
\newunicodechar{μ}{\ensuremath{\mu}}
\newunicodechar{σ}{\ensuremath{\sigma}}
\newunicodechar{τ}{\ensuremath{\tau}}
\newunicodechar{φ}{\ensuremath{\varphi}}
\newunicodechar{ψ}{\ensuremath{\psi}}
\newunicodechar{ω}{\ensuremath{\omega}}
\newunicodechar{Σ}{\ensuremath{\Sigma}}
% majuscules produites par \MakeUppercase dans les titres (α-aminés -> Α)
\newunicodechar{Α}{\ensuremath{\alpha}}
\newunicodechar{Β}{\ensuremath{\beta}}
\newunicodechar{Γ}{\ensuremath{\Gamma}}
\newunicodechar{Δ}{\ensuremath{\Delta}}
\newunicodechar{Ε}{\ensuremath{\varepsilon}}
\newunicodechar{Θ}{\ensuremath{\Theta}}
\newunicodechar{Λ}{\ensuremath{\Lambda}}
\newunicodechar{Μ}{\ensuremath{\mu}}
\newunicodechar{Τ}{\ensuremath{\tau}}
\newunicodechar{Φ}{\ensuremath{\Phi}}
\newunicodechar{Ψ}{\ensuremath{\Psi}}
\newunicodechar{Ω}{\ensuremath{\Omega}}
\newunicodechar{↦}{\ensuremath{\mapsto}}
\newunicodechar{∈}{\ensuremath{\in}}
\newunicodechar{∉}{\ensuremath{\notin}}
\newunicodechar{∧}{\ensuremath{\wedge}}
\newunicodechar{⇔}{\ensuremath{\Leftrightarrow}}
\newunicodechar{⟺}{\ensuremath{\Longleftrightarrow}}
\newunicodechar{⇒}{\ensuremath{\Rightarrow}}
\newunicodechar{⟹}{\ensuremath{\Longrightarrow}}
\newunicodechar{∘}{\ensuremath{\circ}}
\newunicodechar{∪}{\ensuremath{\cup}}
\newunicodechar{∩}{\ensuremath{\cap}}
\newunicodechar{≡}{\ensuremath{\equiv}}
\newunicodechar{⊂}{\ensuremath{\subset}}
\newunicodechar{⊥}{\ensuremath{\perp}}
\newunicodechar{∃}{\ensuremath{\exists}}
\newunicodechar{∀}{\ensuremath{\forall}}
\newunicodechar{∛}{\ensuremath{\sqrt[3]{\,}}}
\newunicodechar{∜}{\ensuremath{\sqrt[4]{\,}}}
\newunicodechar{⃗}{\ensuremath{{}^{\to}}}
\newunicodechar{ⁿ}{\ifmmode^{n}\else\textsuperscript{\ensuremath{n}}\fi}
\newunicodechar{ˣ}{\ifmmode^{x}\else\textsuperscript{\ensuremath{x}}\fi}
\newunicodechar{ᵇ}{\ifmmode^{b}\else\textsuperscript{\ensuremath{b}}\fi}
\newunicodechar{ʳ}{\ifmmode^{r}\else\textsuperscript{\ensuremath{r}}\fi}
\newunicodechar{ᵘ}{\ifmmode^{u}\else\textsuperscript{\ensuremath{u}}\fi}
\newunicodechar{ʸ}{\ifmmode^{y}\else\textsuperscript{\ensuremath{y}}\fi}
\newunicodechar{ᵃ}{\ifmmode^{a}\else\textsuperscript{\ensuremath{a}}\fi}
\newunicodechar{ᵉ}{\ifmmode^{e}\else\textsuperscript{\ensuremath{e}}\fi}
\newunicodechar{ᵏ}{\ifmmode^{k}\else\textsuperscript{\ensuremath{k}}\fi}
\newunicodechar{ᵖ}{\ifmmode^{p}\else\textsuperscript{\ensuremath{p}}\fi}
\newunicodechar{ᴺ}{\ifmmode^{N}\else\textsuperscript{\ensuremath{N}}\fi}
\newunicodechar{ᶜ}{\ifmmode^{c}\else\textsuperscript{\ensuremath{c}}\fi}
\newunicodechar{ʰ}{\ifmmode^{h}\else\textsuperscript{\ensuremath{h}}\fi}
\newunicodechar{ᵠ}{\ifmmode^{q}\else\textsuperscript{\ensuremath{q}}\fi}
\newunicodechar{ₙ}{\ifmmode_{n}\else\textsubscript{\ensuremath{n}}\fi}
\newunicodechar{ₐ}{\ifmmode_{a}\else\textsubscript{\ensuremath{a}}\fi}
\newunicodechar{ᵢ}{\ifmmode_{i}\else\textsubscript{\ensuremath{i}}\fi}
\newunicodechar{ᵣ}{\ifmmode_{r}\else\textsubscript{\ensuremath{r}}\fi}
\newunicodechar{ₖ}{\ifmmode_{k}\else\textsubscript{\ensuremath{k}}\fi}
\newunicodechar{ᵦ}{\ifmmode_{\beta}\else\textsubscript{\ensuremath{\beta}}\fi}
\newunicodechar{ₓ}{\ifmmode_{x}\else\textsubscript{\ensuremath{x}}\fi}
\newfontfamily\popdisplay{ArchivoBlack-Regular.ttf}[Path=../../../fonts/]
\newfontfamily\poplogo{SpaceGrotesk-Bold.ttf}[Path=../../../fonts/]
\newfontfamily\popsemi{PlusJakartaSans-SemiBold.ttf}[Path=../../../fonts/]
\newfontfamily\popxbold{PlusJakartaSans-ExtraBold.ttf}[Path=../../../fonts/]
\newfontfamily\popxboldls{PlusJakartaSans-ExtraBold.ttf}[Path=../../../fonts/,LetterSpace=3.0]

\setlist[enumerate]{leftmargin=1.7em,itemsep=5pt,topsep=4pt}
\setlist[itemize]{leftmargin=1.7em,itemsep=4pt,topsep=4pt,label={\color{poporange}\textbullet}}
\setlength{\parindent}{0pt}
\setlength{\parskip}{5pt}
\renewcommand{\arraystretch}{1.25}

% pgfplots : style Pop par défaut
\pgfplotsset{
  every axis/.append style={axis line style={popink,line width=1pt},tick style={popink},
    grid style={popline,line width=0.5pt},label style={font=\small},tick label style={font=\footnotesize}},
  popcurve/.style={poporange,line width=1.8pt,no markers,smooth},
}
% Même style disponible dans un \draw TikZ simple (hors pgfplots)
\tikzset{popcurve/.style={poporange,line width=1.8pt}}
\tikzset{popgrid/.style={popline,very thin}}
\colorlet{popcurve}{poporange} % le nom du style est parfois utilisé comme couleur

% ============================================================
% Pill / squiggle
% ============================================================
\newcommand{\poppill}[3]{%
  \tikz[baseline=-0.55ex]\node[draw=popink,line width=1.2pt,rounded corners=2.6pt,%
    fill=#2,inner xsep=6.5pt,inner ysep=3pt]{%
    \popxboldls\fontsize{7}{7}\selectfont\color{#3}\MakeUppercase{#1}};%
}
\newcommand{\popsquiggle}[1]{%
  \tikz[baseline=-0.4ex]{
    \draw[#1,line width=2.6pt,line cap=round]
      plot [smooth] coordinates {(0,0.09) (0.34,0.01) (0.68,0.21) (1.02,0.01) (1.36,0.10)};
  }%
}
% Mot-clé surligné (marqueur orange sous le texte)
\newcommand{\cle}[1]{{\popxbold\color{popdark}#1}}

% ============================================================
% En-tête / pied
% ============================================================
\pagestyle{fancy}
\fancyhf{}
\renewcommand{\headrulewidth}{0pt}
\renewcommand{\footrulewidth}{0pt}
\lhead{\raisebox{-0.15ex}{\poplogo\fontsize{13}{13}\selectfont\color{popink}OnBuch\color{poporange}+}}
\rhead{\poppill{\DOCMATIERE\ \textbullet\ \DOCNIVEAU\ \textbullet\ Leçon \DOCLECON}{white}{popink}}
\lfoot{\popxbold\fontsize{7.6}{7.6}\selectfont\color{popmuted}\MakeUppercase{Cours signé OnBuch+}\ \textbullet\ {\popsemi\DOCTITRE}}
\rfoot{\tikz[baseline=-0.6ex]\node[rounded corners=8.5pt,fill=popink,minimum height=17pt,minimum width=17pt,inner xsep=5pt,inner sep=0]{\popxbold\fontsize{8}{8}\selectfont\color{popcream}\thepage\,/\,\pageref*{LastPage}};}

% ============================================================
% Titres
% ============================================================
\newcommand{\chaptitre}[2]{%
  \begin{tcolorbox}[enhanced,colback=poporange,colframe=popink,boxrule=2.6pt,arc=11pt,
      drop shadow=popink,left=15pt,right=15pt,top=12pt,bottom=11pt,before skip=3pt,after skip=18pt]
    \poppill{#2}{popink}{popcream}\par\vspace{9pt}
    {\raggedright\hyphenpenalty=10000\exhyphenpenalty=10000\popdisplay\fontsize{22}{24}\selectfont\color{popcream}\MakeUppercase{#1}\par}
  \end{tcolorbox}
}
% Section numérotée : pastille orange avec le numéro + titre Archivo
\newcounter{popsec}
\newcommand{\coursec}[1]{%
  \stepcounter{popsec}\setcounter{popsub}{0}%
  \par\vspace{16pt}\noindent
  \begin{minipage}{\linewidth}
  \tikz[baseline=-0.7ex]\node[draw=popink,line width=1.6pt,fill=poporange,rounded corners=4pt,minimum size=22pt,inner sep=0]{\popdisplay\fontsize{12}{12}\selectfont\color{popcream}\thepopsec};%
  \hspace{8pt}{\popdisplay\fontsize{15}{17}\selectfont\color{popink}\MakeUppercase{#1}}\par
  \vspace{4pt}\noindent{\color{poporange}\rule{36pt}{3.6pt}}
  \end{minipage}\par\nopagebreak\vspace{8pt}
}
\newcounter{popsub}
\newcommand{\courssub}[1]{%
  \stepcounter{popsub}%
  \par\vspace{9pt}\noindent{\popxbold\fontsize{11.5}{13}\selectfont\color{popink}{\color{poporange}\thepopsec.\thepopsub}\hspace{6pt}#1}\par\nopagebreak\vspace{4pt}
}

% ============================================================
% Boîtes pédagogiques — même recette : fond blanc, bordure épaisse colorée,
% ombre dure encre, badge plein. Titre optionnel [#1].
% ============================================================
\tcbset{popbase/.style={breakable,enhanced,colback=white,boxrule=2.3pt,arc=9pt,
  shadow={3pt}{-3pt}{0pt}{popink},left=14pt,right=14pt,top=11pt,bottom=11pt,before skip=11pt,after skip=15pt,
  fontupper=\popsemi\fontsize{10.6}{14.5}\selectfont\color{popink},
  fontlower=\popsemi\fontsize{10.6}{14.5}\selectfont\color{popink}}}
% Coche dessinée (le glyphe ✓ n'existe pas dans les polices du document)
\newcommand{\popcheck}[1]{\tikz[baseline=-0.5ex]\draw[#1,line width=1.6pt,line cap=round,line join=round](-0.13,0.0)--(-0.04,-0.09)--(0.13,0.1);}
\providecommand{\coche}{\popcheck{popgreen}}
\providecommand{\resultat}[1]{\par\smallskip\noindent\fcolorbox{popgreen}{popgreenL}{\parbox{\dimexpr\linewidth-2\fboxsep-2\fboxrule\relax}{\textbf{Résultat.} #1}}\par\smallskip}
\newcommand{\popboxhead}[3]{\poppill{#1}{#2}{white}\ifx\relax#3\relax\else\hspace{6pt}{\popxbold\fontsize{10.6}{12}\selectfont\color{popink}#3}\fi\par\vspace{7pt}}

\newtcolorbox{aretenir}[1][]{popbase,colframe=popgreen,before upper={\popboxhead{À retenir}{popgreen}{#1}}}
\newtcolorbox{exemplebox}[1][]{popbase,colframe=poporange,before upper={\popboxhead{Exemple}{poporange}{#1}}}
\newtcolorbox{definition}[1][]{popbase,colframe=popblue,before upper={\popboxhead{Définition}{popblue}{#1}}}
\newtcolorbox{propriete}[1][]{popbase,colframe=poppurple,before upper={\popboxhead{Propriété}{poppurple}{#1}}}
\newtcolorbox{methode}[1][]{popbase,colframe=popgold,before upper={\popboxhead{Méthode}{popgold}{#1}}}
\newtcolorbox{attention}[1][]{popbase,colframe=poppink,before upper={\popboxhead{Attention}{poppink}{#1}}}
\newtcolorbox{experience}[1][]{popbase,colframe=popblue,colback=popblueL,before upper={\popboxhead{Expérience}{popblue}{#1}}}
% « Le savais-tu ? » : carte violette pleine (contexte, culture, Cameroun)
\newtcolorbox{savaistu}[1][]{popbase,colback=poppurple,colframe=popink,
  fontupper=\popsemi\fontsize{10.4}{14.2}\selectfont\color{white},
  before upper={\poppill{Le savais-tu\,?}{white}{poppurple}\ifx\relax#1\relax\else\hspace{6pt}{\popxbold\color{white}#1}\fi\par\vspace{7pt}}}
% Exercice résolu : énoncé en haut, \tcblower puis la solution sur fond crème
\newcounter{popexo}\newcounter{popexr}
\newtcolorbox{exoresolu}[1][]{popbase,colframe=poporange,colbacklower=popcream,
  segmentation style={popink,line width=1.2pt,dashed},
  before upper={\stepcounter{popexr}\popboxhead{Exercice résolu \thepopexr}{poporange}{#1}},
  before lower={\poppill{Solution}{popink}{popcream}\par\vspace{6pt}}}
% Exercice (non résolu en place) — la correction est dans la partie Corrigés
\newtcolorbox{exercice}[1][]{popbase,colframe=popink,
  before upper={\stepcounter{popexo}\popboxhead{Exercice \thepopexo}{popink}{#1}}}
% Corrigé
\newtcolorbox{corrige}[1][]{popbase,colframe=popgreen,colback=popgreenL,
  before upper={\popboxhead{Corrigé}{popgreen}{#1}}}
% Objectifs / prérequis / bilan
\newtcolorbox{objectifs}{popbase,colframe=popink,colback=poporangeL,
  before upper={\popboxhead{Objectifs de la leçon}{popink}{}}}
\newtcolorbox{prerequis}{popbase,colframe=popink,colback=popblueL,
  before upper={\popboxhead{Prérequis}{popblue}{}}}
\newtcolorbox{bilan}{popbase,colframe=popink,colback=white,boxrule=2.8pt,
  before upper={\popboxhead{Fiche bilan}{poporange}{L'essentiel à connaître pour l'examen}}}
% Figure encadrée avec légende
\newtcolorbox{popfigure}[1][]{enhanced,breakable=false,colback=white,colframe=popline,boxrule=1.4pt,arc=8pt,
  left=10pt,right=10pt,top=10pt,bottom=8pt,before skip=10pt,after skip=12pt,halign=center,
  lower separated=false,halign lower=center,
  fontlower=\popsemi\fontsize{9}{11.5}\selectfont\color{popmuted},
  #1}
\newcommand{\legende}[1]{\par\vspace{6pt}{\popsemi\fontsize{9}{11.5}\selectfont\color{popmuted}#1}}

% ============================================================
% Signature OnBuch+
% ============================================================
\newcommand{\poplogoinline}[1]{{\poplogo\fontsize{#1}{#1}\selectfont\color{popink}OnBuch\color{poporange}+}}
\newcommand{\signatureonbuch}{%
  \begin{tcolorbox}[enhanced,colback=popink,colframe=popink,boxrule=2.2pt,arc=10pt,
      drop shadow=poporange,left=14pt,right=14pt,top=10pt,bottom=10pt,before skip=0pt,after skip=0pt]
    \tikz[baseline=-0.7ex]\node[circle,fill=popgreen,draw=popcream,line width=1.3pt,minimum size=22pt,inner sep=0]{\popcheck{white}};%
    \hspace{9pt}\begin{minipage}[c]{0.62\linewidth}
      {\popxbold\fontsize{12}{13}\selectfont\color{popcream}Signé\ \textbullet\ {\poplogo OnBuch\color{poporange}+}}\par\vspace{2pt}
      {\raggedright\popsemi\fontsize{8}{10.5}\selectfont\color{popline}\MakeUppercase{Équipe pédagogique OnBuch+}\\\MakeUppercase{Conforme au programme officiel MINESEC}\par}
    \end{minipage}\hfill
    {\poplogo\fontsize{22}{22}\selectfont\color{popcream}OnBuch\color{poporange}+}
  \end{tcolorbox}%
}

% ============================================================
% Page de garde
% #1 titre · #2 matière · #3 niveau · #4 module · #5 n° leçon · #6 durée ·
% #7 description / objectifs (texte ou itemize)
% ============================================================
\newcommand{\pagegarde}[7]{%
  \thispagestyle{empty}
  \newgeometry{a4paper,margin=1.7cm,top=1.6cm,bottom=1.4cm}%
  \begin{tikzpicture}[remember picture,overlay]
    \fill[poporange!18] ([xshift=-3.2cm,yshift=-3.6cm]current page.north east) circle (4.6cm);
    \fill[poppurple!14] ([xshift=2.4cm,yshift=7.6cm]current page.south west) circle (3.8cm);
  \end{tikzpicture}%
  {\poplogo\fontsize{30}{30}\selectfont\color{popink}OnBuch\color{poporange}+}\hfill
  \raisebox{0.6ex}{\poppill{Cours signé \textbullet\ Premium}{popink}{popcream}}\par
  \vspace{1pt}\popsquiggle{poppurple}\par
  \vspace{8pt}
  \begin{tcolorbox}[enhanced,colback=poporange,colframe=popink,boxrule=2.8pt,arc=13pt,
      drop shadow=popink,left=18pt,right=18pt,top=12pt,bottom=13pt,after skip=10pt]
    \poppill{#2 \textbullet\ #3}{popink}{popcream}\hspace{5pt}\poppill{#4}{white}{popink}\par\vspace{8pt}
    {\popdisplay\fontsize{14}{14}\selectfont\color{popink}LEÇON #5}\par\vspace{4pt}
    {\raggedright\hyphenpenalty=10000\exhyphenpenalty=10000\popdisplay\fontsize{25}{27}\selectfont\color{popcream}\MakeUppercase{#1}\par}
    \vspace{8pt}
    {\popxbold\fontsize{9.5}{9.5}\selectfont\color{popink}\MakeUppercase{Durée conseillée : #6}}
  \end{tcolorbox}
  \begin{tcolorbox}[enhanced,colback=white,colframe=popink,boxrule=2.2pt,arc=10pt,
      drop shadow=popink,left=16pt,right=16pt,top=10pt,bottom=10pt,after skip=8pt]
    \poppill{Dans cette leçon}{poppurple}{white}\par\vspace{5pt}
    \popsemi\fontsize{9.3}{12.6}\selectfont\color{popink}#7
  \end{tcolorbox}
  \noindent\begin{minipage}[t]{0.487\linewidth}
    \begin{tcolorbox}[enhanced,colback=popgreen,colframe=popink,boxrule=2.2pt,arc=10pt,
        drop shadow=popink,left=11pt,right=11pt,top=10pt,bottom=10pt,height=3.1cm,valign=center]
      \tcbox[colback=white,colframe=popink,boxrule=1.3pt,arc=5pt,boxsep=0pt,nobeforeafter,
        left=3pt,right=3pt,top=3pt,bottom=3pt,valign=center]{\qrcode[height=1.9cm]{https://play.google.com/store/apps/details?id=com.luvvix.onbuch&hl=fr}}%
      \hspace{8pt}\begin{minipage}[c]{0.5\linewidth}
        {\popdisplay\fontsize{11}{12.5}\selectfont\color{white}\MakeUppercase{Télécharge OnBuch}}\par\vspace{4pt}
        {\raggedright\popsemi\fontsize{8.4}{11}\selectfont\color{white}Gratuit \textbullet\ Google Play\\Tuteur IA, annales, concours, résultats\par}
      \end{minipage}
    \end{tcolorbox}
  \end{minipage}\hfill
  \begin{minipage}[t]{0.487\linewidth}
    \begin{tcolorbox}[enhanced,colback=poppurple,colframe=popink,boxrule=2.2pt,arc=10pt,
        drop shadow=popink,left=11pt,right=11pt,top=10pt,bottom=10pt,height=3.1cm,valign=center]
      \tikz[baseline=-0.7ex]\node[circle,fill=white,draw=popink,line width=1.3pt,minimum size=1.6cm,inner sep=0]{\popdisplay\fontsize{24}{24}\selectfont\color{poppurple}+};%
      \hspace{8pt}\begin{minipage}[c]{0.6\linewidth}
        {\popdisplay\fontsize{11}{12.5}\selectfont\color{white}\MakeUppercase{Passe à OnBuch+}}\par\vspace{4pt}
        {\raggedright\popsemi\fontsize{8.4}{11}\selectfont\color{white}Tous les cours signés, Tuteur IA illimité, fiches PDF — onglet OnBuch+ de l'app\par}
      \end{minipage}
    \end{tcolorbox}
  \end{minipage}
  \vspace{8pt}
  \popcontenu
  \vfill
  \signatureonbuch
  \restoregeometry
  \newpage
}

% Rangée « Ce cours contient » (page de garde)
\newcommand{\poptile}[3]{% #1 couleur #2 gros texte #3 légende
  \begin{tcolorbox}[enhanced,colback=#1,colframe=popink,boxrule=2pt,arc=9pt,shadow={2.5pt}{-2.5pt}{0pt}{popink},
      left=6pt,right=6pt,top=9pt,bottom=9pt,halign=center,valign=center,height=1.75cm,width=\linewidth,nobeforeafter]
    {\popdisplay\fontsize{12.5}{13}\selectfont\color{white}\MakeUppercase{#2}}\par\vspace{3pt}
    {\centering\popsemi\fontsize{7.8}{9.5}\selectfont\color{white}#3\par}
  \end{tcolorbox}}
\newcommand{\popcontenu}{%
  \par\noindent{\popxboldls\fontsize{7.5}{7.5}\selectfont\color{popmuted}CE COURS SIGNÉ CONTIENT}\par\vspace{7pt}
  \noindent\begin{minipage}[t]{0.235\linewidth}\poptile{popblue}{Cours}{complet, illustré, pas à pas}\end{minipage}\hfill
  \begin{minipage}[t]{0.235\linewidth}\poptile{poporange}{Méthodes}{et exercices résolus}\end{minipage}\hfill
  \begin{minipage}[t]{0.235\linewidth}\poptile{poppink}{Exercices}{type examen + corrigés}\end{minipage}\hfill
  \begin{minipage}[t]{0.235\linewidth}\poptile{popgreen}{Fiche bilan}{l'essentiel à retenir}\end{minipage}\par}

% Sommaire
\newtcolorbox{sommaire}{enhanced,colback=white,colframe=popink,boxrule=2.2pt,arc=10pt,
  drop shadow=popink,left=18pt,right=18pt,top=13pt,bottom=13pt,before skip=6pt,after skip=20pt,
  before upper={\poppill{Sommaire}{poppurple}{white}\par\vspace{9pt}\popsemi\fontsize{10.6}{14.5}\selectfont\color{popink}}}

% Signature de fin de cours — poussée tout en bas de la dernière page
\newcommand{\findecours}{%
  \par\vfill\nopagebreak
  \begin{center}\popsquiggle{poporange}\end{center}\vspace{4pt}
  \signatureonbuch
}
\AtBeginDocument{\def\llbracket{\mathopen{[\mkern-3.5mu[}}\def\rrbracket{\mathclose{]\mkern-3.5mu]}}} % intervalles d'entiers (glyphe ⟦ absent de la police)
% Glyphes absents de Fira Math : redéfinis à partir de glyphes disponibles
\AtBeginDocument{%
  \def\setminus{\mathbin{\backslash}}\let\smallsetminus\setminus
  \def\vdots{\mathord{\vbox{\baselineskip3.2pt\lineskiplimit0pt\kern4pt\hbox{.}\hbox{.}\hbox{.}}}}%
  \def\ddots{\mathinner{\mkern1mu\raise6.5pt\hbox{.}\mkern2mu\raise3.6pt\hbox{.}\mkern2mu\raise0.7pt\hbox{.}\mkern1mu}}%
  \def\triangle{\mathord{\scalebox{0.9}{$\symup{\Delta}$}}}%
  \def\nabla{\mathord{\raisebox{\depth}{\scalebox{1}[-1]{$\symup{\Delta}$}}}}%
  \def\checkmark{\popcheck{popdark}}%
  \def\star{\mathbin{\ast}}\def\bigstar{\mathord{\ast}}%
  \def\diamond{\mathbin{\scalebox{0.6}{$\Diamond$}}}%
  \def\bigcup{\mathop{\mathchoice{\raisebox{-0.3em}{\scalebox{1.7}{$\cup$}}}{\scalebox{1.25}{$\cup$}}{\cup}{\cup}}\displaylimits}%
  \def\bigcap{\mathop{\mathchoice{\raisebox{-0.3em}{\scalebox{1.7}{$\cap$}}}{\scalebox{1.25}{$\cap$}}{\cap}{\cap}}\displaylimits}%
  \def\lesseqqgtr{\mathrel{\lessgtr}}\def\bigoplus{\mathop{\oplus}\displaylimits}%
}
\providecommand{\ssi}{si et seulement si} % abréviation fréquente
\AtBeginDocument{\providecommand{\wideparen}{\overparen}} % arc de cercle

%% FIGFIX-BEGIN v4 — correctifs globaux des figures (docs/onbuch-generation/tools/figfix)
\usepackage{adjustbox}\usepackage{etoolbox}
% toute figure TikZ/pgfplots plus large que la ligne est réduite à la largeur disponible
\BeforeBeginEnvironment{tikzpicture}{\begin{adjustbox}{max width=\linewidth}}
\AfterEndEnvironment{tikzpicture}{\end{adjustbox}}
% légendes pgfplots lisibles par-dessus les courbes
\pgfplotsset{every axis/.append style={legend style={fill=white,fill opacity=0.92,text opacity=1,draw=black!15,inner sep=3pt}}}
% étiquettes d'arêtes (syntaxe edge["..."]) avec fond blanc
\tikzset{every edge quotes/.append style={fill=white,inner sep=1.2pt}}
%% FIGFIX-END
\begin{document}

\pagegarde{Formulation des idées essentielles d’un texte à résumer}{Français}{Tle Tech}{Français – classe de Terminale technique}{N21}{10 séances (fiche de progression harmonisée)}{Cette leçon outille l'élève de Terminale technique pour identifier la structure argumentative d'un texte, en extraire les idées essentielles et les reformuler dans une contraction de texte réussie. Elle mobilise la lecture méthodique du Vieux nègre et la médaille, les outils de liaison de la phrase et la distinction entre contenus manifestes et latents.
\vspace{4pt}
\begin{multicols}{2}\small
\begin{itemize}
\item À la fin de cette leçon, je sais identifier la thèse, les arguments et les exemples dans un texte argumentatif.
\item Je sais repérer les procédés stylistiques et rhétoriques qui soutiennent l'argumentation.
\item Je sais distinguer les idées essentielles des idées secondaires (exemples, répétitions, digressions).
\item Je sais appliquer les techniques de réduction : suppression, substitution, généralisation et fusion.
\item Je sais reformuler les idées essentielles avec mes propres mots sans trahir le sens du texte.
\item Je sais utiliser correctement les conjonctions de coordination et de subordination pour lier les idées dans mon résumé.
\end{itemize}\end{multicols}}

\begin{sommaire}
\begin{multicols}{2}
\begin{enumerate}
\item Le texte argumentatif : structure et procédés
\item Lecture méthodique n°4 et n°5 : Le Vieux nègre et la médaille
\item Les techniques de réduction du texte
\item Les outils de liaison : conjonctions de coordination et de subordination
\item Contenus manifestes et contenus latents
\item Reformulation, exposé et activités d'intégration
\item Activité d'intégration
\item Méthodes et exercices résolus
\item Exercices
\item Corrigés des exercices
\item Fiche bilan
\end{enumerate}
\end{multicols}
\end{sommaire}

\begin{objectifs}
\begin{itemize}
\item À la fin de cette leçon, je sais identifier la thèse, les arguments et les exemples dans un texte argumentatif.
\item Je sais repérer les procédés stylistiques et rhétoriques qui soutiennent l'argumentation.
\item Je sais distinguer les idées essentielles des idées secondaires (exemples, répétitions, digressions).
\item Je sais appliquer les techniques de réduction : suppression, substitution, généralisation et fusion.
\item Je sais reformuler les idées essentielles avec mes propres mots sans trahir le sens du texte.
\item Je sais utiliser correctement les conjonctions de coordination et de subordination pour lier les idées dans mon résumé.
\item Je sais distinguer les contenus manifestes et les contenus latents d'un message.
\item Je sais préparer et présenter un exposé oral structuré.
\item Je sais rédiger et justifier une démarche de contraction de texte respectant le nombre de mots imposé.
\end{itemize}
\end{objectifs}

\begin{prerequis}
\begin{itemize}
\item La notion de texte argumentatif et ses types de thèse (classe de Première)
\item Les connecteurs logiques simples (d'abord, ensuite, donc, cependant)
\item La lecture méthodique : repérage du thème et de la problématique
\item Les classes grammaticales et la phrase simple (niveau 4e)
\item Notions de base sur le résumé vues en Première
\end{itemize}
\end{prerequis}

\newpage

\coursec{Le texte argumentatif : structure et procédés}

\textbf{Situation d'accroche.} Lors des épreuves du Baccalauréat technique, de nombreux candidats de Yaoundé ou de Douala recopient de longs passages du texte au lieu de le résumer, et perdent des points précieux. Pourtant, savoir dégager et reformuler les idées essentielles d'un texte est une compétence utile partout : résumer un discours entendu à la radio CRTV, rédiger un compte rendu de réunion dans une entreprise, ou présenter un exposé devant la classe. Comment passer d'un texte de 500 mots à un résumé fidèle de 125 mots sans trahir la pensée de l'auteur ? La première étape de la méthode consiste à bien comprendre le texte à résumer. Or, la plupart des textes proposés à l'examen sont des \cle{textes argumentatifs}. C'est donc par eux que nous commençons.

\textbf{Problématique.} Qu'est-ce qu'un texte argumentatif ? Comment est-il construit ? Quels procédés l'auteur utilise-t-il pour nous convaincre ou nous persuader ?

\courssub{Qu'est-ce qu'un texte argumentatif ?}

\textbf{Intuition.} Imaginez un débat à la maison : votre grand frère veut convaincre votre père d'acheter une moto pour les livraisons de la boutique familiale. Il ne se contente pas de dire « il nous faut une moto ». Il explique que les clients se plaignent des retards, que le taxi-moto coûte cher à la longue, qu'un voisin a doublé son chiffre d'affaires grâce à une moto. Il défend une \emph{idée} (sa thèse) à l'aide de \emph{raisons} (ses arguments) appuyées par des \emph{faits concrets} (ses exemples). Un texte argumentatif fonctionne exactement de la même manière, par écrit.

\begin{definition}[Texte argumentatif, thèse, argument, exemple]
Un \cle{texte argumentatif} est un texte dans lequel un auteur défend une opinion (une \cle{thèse}) afin d'amener le lecteur à l'adopter. On distingue :
\begin{itemize}
\item la \cle{thèse} : l'idée principale que l'auteur veut faire accepter ;
\item l'\cle{argument} : une raison avancée pour justifier la thèse ;
\item l'\cle{exemple} : un cas concret, un fait précis qui illustre et renforce un argument.
\end{itemize}
L'argumentation poursuit deux buts possibles : \textbf{convaincre} (faire adhérer l'intelligence du lecteur par le raisonnement) ou \textbf{persuader} (toucher ses émotions, ses valeurs, pour le faire adhérer).
\end{definition}

\begin{exemplebox}[Convaincre ou persuader ?]
Pour défendre la scolarisation des filles, un journaliste peut :
\begin{itemize}
\item \emph{convaincre} : « Selon le MINESEC, une fille scolarisée jusqu'au secondaire fait reculer la pauvreté de son ménage ; c'est un investissement rentable pour le pays. » (raisonnement, chiffres) ;
\item \emph{persuader} : « Faut-il vraiment priver nos sœurs, nos filles, de la lumière de l'école et les condamner à l'ombre de l'ignorance ? » (émotions, valeurs, images).
\end{itemize}
\end{exemplebox}

\courssub{Les trois types de thèse}

La thèse n'est pas toujours écrite noir sur blanc dans le texte. Il faut savoir la débusquer. On distingue trois cas.

\begin{aretenir}[Les trois types de thèse]
\begin{itemize}
\item \textbf{Thèse explicite} : elle est clairement formulée dans le texte, souvent au début ou à la fin. Exemple : « Le téléphone portable doit être interdit en classe. »
\item \textbf{Thèse implicite} : elle n'est pas écrite, mais elle se déduit des arguments et des exemples ; le lecteur doit la reconstruire. Exemple : un texte qui accumule les accidents causés par le téléphone au volant laisse deviner que l'auteur condamne cet usage.
\item \textbf{Thèse suggérée} : elle est transmise par des moyens indirects (ironie, caricature, question rhétorique, récit symbolique). C'est le cas fréquent des contes et des romans comme \emph{Le Vieux nègre et la médaille} de Ferdinand Oyono, où la critique de la colonisation passe par l'humour et le récit.
\end{itemize}
\end{aretenir}

\begin{methode}[Repérer la thèse d'un texte en trois questions]
Pour trouver la thèse de n'importe quel texte argumentatif, posez-vous dans l'ordre trois questions :
\begin{enumerate}
\item \textbf{De quoi parle l'auteur ?} (le thème : la scolarisation des filles, le portable en classe...) ;
\item \textbf{Que veut-il prouver ?} (sa position : pour ou contre ? quelle idée défend-il ?) ;
\item \textbf{Comment le prouve-t-il ?} (les arguments et exemples qu'il aligne : cela confirme la thèse repérée).
\end{enumerate}
La réponse à la deuxième question, formulée en une phrase complète, est la thèse du texte.
\end{methode}

\courssub{La structure argumentative}

Un texte argumentatif bien construit ressemble à un escalier : on monte marche par marche, de la thèse vers la conclusion, en s'appuyant sur des arguments illustrés d'exemples.

\begin{aretenir}[Le plan type du texte argumentatif]
\begin{enumerate}
\item \textbf{Introduction} : elle présente le sujet et énonce (ou prépare) la thèse.
\item \textbf{Développement} : il aligne les arguments, chacun étayé par un ou plusieurs exemples. Les arguments sont reliés par des connecteurs logiques.
\item \textbf{Conclusion} : elle rappelle la thèse, la confirme et peut ouvrir sur une idée nouvelle.
\end{enumerate}
\end{aretenir}

\begin{popfigure}
\begin{center}
\begin{tikzpicture}[font=\small]
\tikzset{
    marche/.style={draw=popblue, fill=popblueL, rounded corners=2pt, minimum height=1cm, text width=3.2cm, align=center}
}
\node[marche, align=center] (t) at (0,0){\textbf{Introduction}\\ Thèse};
\node[marche, align=center] (a1) at (3.6,1.1){\textbf{Argument 1}\\ + exemple};
\node[marche, align=center] (a2) at (7.2,2.2){\textbf{Argument 2}\\ + exemple};
\node[marche, align=center] (c) at (10.8,3.3){\textbf{Conclusion}\\ retour sur la thèse};
\draw[-{Stealth}, thick, poporange] (t) -- (a1);
\draw[-{Stealth}, thick, poporange] (a1) -- (a2);
\draw[-{Stealth}, thick, poporange] (a2) -- (c);
\end{tikzpicture}
\end{center}
\legende{Schéma en escalier de la structure argumentative : chaque argument, illustré d'un exemple, fait progresser le lecteur de la thèse vers la conclusion.}
\end{popfigure}

\courssub{Les types d'arguments : logiques et affectifs}

Tous les arguments ne jouent pas sur le même registre. Certains s'adressent à notre raison, d'autres à notre cœur.

\begin{definition}[Arguments logiques et arguments affectifs]
\begin{itemize}
\item Les \cle{arguments logiques} (ou rationnels) s'appuient sur des \textbf{faits vérifiables}, des \textbf{statistiques}, des définitions, des raisonnements de cause à effet. Ils servent à \emph{convaincre}.
\item Les \cle{arguments affectifs} (ou émotionnels) font appel aux \textbf{valeurs} (justice, honneur, solidarité, religion), aux \textbf{sentiments} (peur, fierté, pitié), aux images frappantes. Ils servent à \emph{persuader}.
\end{itemize}
\end{definition}

\begin{exemplebox}[Les deux types d'arguments sur le même sujet]
Thèse : « Il faut lutter contre la déforestation au Cameroun. »
\begin{itemize}
\item Argument logique : « La forêt camerounaise perd chaque année des milliers d'hectares, ce qui menace l'équilibre du climat et les ressources en bois de demain. »
\item Argument affectif : « Voulons-nous laisser à nos enfants un pays nu, sans arbres, sans singes ni oiseaux, un pays qui a vendu son âme verte ? »
\end{itemize}
\end{exemplebox}

\courssub{Les procédés rhétoriques}

Pour rendre son argumentation plus efficace, l'auteur utilise des \cle{procédés rhétoriques}, c'est-à-dire des outils de langue qui frappent l'esprit ou le cœur du lecteur.

\begin{definition}[Principaux procédés rhétoriques]
\begin{itemize}
\item La \cle{question rhétorique} : une question qui n'attend pas de réponse, car la réponse est évidente. « Faut-il accepter que nos enfants grandissent sans instruction ? »
\item L'\cle{ironie} : dire le contraire de ce qu'on pense pour mieux ridiculiser. « Bravo, vraiment, à ceux qui jettent leurs ordures dans la rigole ! »
\item L'\cle{antithèse} : l'opposition de deux idées ou de deux mots contraires dans une même phrase. « L'école ouvre les portes de la liberté ; l'ignorance ferme celles de l'espoir. »
\item L'\cle{accumulation} : l'énumération de mots de même nature pour donner de l'ampleur. « Il ment, il triche, il vole, il corrompt. »
\item La \cle{gradation} : une énumération dont la force va en augmentant (gradation ascendante) ou en diminuant (gradation descendante). « Un retard, puis une absence, puis un échec, puis l'abandon des études. »
\end{itemize}
\end{definition}

Certains procédés visent surtout la raison, d'autres surtout l'émotion. Le tableau suivant les classe.

\begin{center}
\begin{tabularx}{\linewidth}{|X|X|}
\hline
\rowcolor{popblueL} \textbf{Procédés pour convaincre (raison)} & \textbf{Procédés pour persuader (émotion)} \\
\hline
Faits, chiffres, statistiques & Images frappantes, métaphores \\
\hline
Définitions précises & Appel aux valeurs (patrie, justice, famille) \\
\hline
Raisonnement logique (cause $\rightarrow$ conséquence) & Question rhétorique, ironie \\
\hline
Antithèse qui éclaire une opposition d'idées & Accumulation, gradation, répétition \\
\hline
Citations d'autorité (experts, rapports) & Témoignages, anecdotes émouvantes \\
\hline
\end{tabularx}
\end{center}

\courssub{Les connecteurs qui balisent l'argumentation}

Les \cle{connecteurs logiques} sont les panneaux indicateurs du texte argumentatif : ils montrent la relation entre les idées et guident le lecteur. Les repérer, c'est déjà comprendre l'architecture du texte — et c'est indispensable pour le résumer sans se tromper.

\begin{aretenir}[Les grandes familles de connecteurs]
\begin{itemize}
\item \textbf{Addition} : et, de plus, également, en outre, non seulement... mais aussi ;
\item \textbf{Cause} : car, parce que, en effet, puisque, comme ;
\item \textbf{Conséquence} : donc, par conséquent, c'est pourquoi, ainsi, de sorte que ;
\item \textbf{Opposition} : mais, cependant, pourtant, néanmoins, alors que, en revanche ;
\item \textbf{Concession} : certes, bien sûr, il est vrai que... mais ;
\item \textbf{Illustration} : par exemple, ainsi, c'est le cas de, notamment ;
\item \textbf{Conclusion} : enfin, en somme, finalement, en définitive.
\end{itemize}
\end{aretenir}

\begin{attention}[Ne pas confondre l'argument et l'exemple]
Erreur fréquente à l'examen : prendre un exemple pour une idée essentielle. L'\textbf{argument} est une idée générale qui soutient la thèse ; l'\textbf{exemple} n'est qu'une illustration secondaire, un cas particulier. Dans un résumé, on garde les arguments et on supprime presque toujours les exemples. Test pratique : si l'on peut retirer la phrase sans que le raisonnement s'effondre, c'est un exemple ; si le raisonnement perd un maillon, c'est un argument.
\end{attention}

\courssub{Application : analyse d'un éditorial}

Lisons cet extrait, inspiré d'un éditorial du quotidien \emph{Cameroon Tribune} consacré à la scolarisation des jeunes filles :

\begin{exemplebox}[Texte à analyser]
« Scolariser la jeune fille, c'est construire l'avenir du Cameroun. En effet, une mère instruite fait manger, soigner et étudier toute sa famille : les statistiques du MINSANTE montrent que les enfants des mères scolarisées sont mieux vaccinés et mieux nourris. De plus, l'éducation des filles retarde les mariages précoces, véritable fléau dans certaines régions. Certes, des coutumes résistent encore ; cependant, aucune tradition ne peut justifier que l'on prive la moitié du pays de la lumière du savoir. Faut-il rappeler que derrière chaque grande réussite camerounaise se cache souvent une mère lettrée ? Par conséquent, investir dans l'éducation des filles n'est pas une dépense : c'est le plus sûr des investissements. »
\end{exemplebox}

\begin{exoresolu}[Analyse argumentative de l'extrait]
Énoncé : 1) Dégagez la thèse du texte. 2) Relevez deux arguments en précisant leur type. 3) Identifiez trois connecteurs et leur fonction. 4) Relevez deux procédés rhétoriques.
\tcblower
Solution détaillée :
\begin{enumerate}
\item \textbf{Thèse} (explicite, énoncée dès la première phrase) : il faut scolariser les jeunes filles, car c'est construire l'avenir du Cameroun.
\item \textbf{Arguments} :
\begin{itemize}
\item « une mère instruite fait manger, soigner et étudier toute sa famille », appuyé par les statistiques du MINSANTE : argument \emph{logique} (fait vérifiable) ;
\item « l'éducation des filles retarde les mariages précoces » : argument \emph{logique} (relation de cause à effet) ;
\item « priver la moitié du pays de la lumière du savoir » : argument \emph{affectif} (image, valeur de justice).
\end{itemize}
\item \textbf{Connecteurs} : « En effet » (cause/justification), « De plus » (addition d'un argument), « Certes... cependant » (concession puis opposition), « Par conséquent » (conséquence, annonce la conclusion).
\item \textbf{Procédés rhétoriques} : la question rhétorique « Faut-il rappeler que... ? » ; l'antithèse finale « dépense / investissement » ; la métaphore « la lumière du savoir ».
\end{enumerate}
\end{exoresolu}

\begin{exercice}[À vous de jouer]
Voici un extrait sur l'usage du téléphone portable en classe au Cameroun : « Le téléphone portable est devenu le premier ennemi de la concentration en classe. Car comment un élève peut-il suivre une leçon de mathématiques tout en consultant ses messages ? Une enquête menée dans trois lycées de Douala montre que les élèves qui utilisent leur téléphone en classe obtiennent en moyenne trois points de moins. Certes, le téléphone peut servir à chercher une information ; mais en réalité, il distrait bien plus qu'il n'instruit. Il faut donc l'interdire purement et simplement pendant les cours. »
1) Formulez la thèse du texte. 2) Relevez deux arguments et dites s'ils sont logiques ou affectifs. 3) Relevez les connecteurs et indiquez leur fonction. 4) Quel procédé rhétorique est utilisé dans la deuxième phrase ?
\end{exercice}

\begin{corrige}[Exercice « À vous de jouer »]
1) Thèse (explicite, dernière phrase) : il faut interdire le téléphone portable pendant les cours.
2) Arguments : le téléphone empêche la concentration (logique) ; les élèves qui l'utilisent perdent en moyenne trois points (logique, appuyé sur une enquête = exemple statistique). L'expression « premier ennemi de la concentration » relève aussi de l'affectif (image forte).
3) Connecteurs : « Car » (cause), « Certes... mais » (concession puis opposition), « donc » (conséquence/conclusion).
4) La deuxième phrase est une question rhétorique : la réponse est sous-entendue (il ne peut pas suivre correctement).
\end{corrige}

\begin{savaistu}[La structure argumentative au Baccalauréat technique]
Au Baccalauréat technique, la question portant sur la structure argumentative d'un texte (dégager la thèse, relever les arguments, identifier les connecteurs) rapporte souvent de 2 à 4 points à l'épreuve de Français. C'est aussi la compétence de base du résumé : un candidat qui sait repérer la thèse et les arguments ne recopie plus le texte, il le reformule. Maîtriser cette leçon, c'est donc préparer à la fois la question de lecture et l'épreuve de contraction de texte.
\end{savaistu}

\begin{aretenir}[L'essentiel de la section]
Un texte argumentatif défend une \cle{thèse} (explicite, implicite ou suggérée) pour \emph{convaincre} (raison) ou \emph{persuader} (émotion). Il se construit en trois parties : introduction (thèse), développement (arguments + exemples), conclusion. Les arguments sont \emph{logiques} (faits, statistiques) ou \emph{affectifs} (valeurs, émotions), et s'enchaînent grâce aux \cle{connecteurs logiques} (car, en effet, cependant, par conséquent...). Les procédés rhétoriques (question rhétorique, ironie, antithèse, accumulation, gradation) renforcent l'efficacité du propos. Repérer cette structure est la première étape indispensable du résumé.
\end{aretenir}

\coursec{Lecture méthodique n°4 et n°5 : Le Vieux nègre et la médaille}

Dans la section précédente, nous avons étudié la structure du texte argumentatif : une thèse, des arguments, des exemples, des procédés de persuasion. Mais un texte littéraire peut aussi argumenter \emph{sans avoir l'air d'argumenter}. C'est précisément le cas du roman de Ferdinand Oyono, \emph{Le Vieux nègre et la médaille} : derrière une histoire racontée avec humour se cache une véritable démonstration contre la colonisation. Pour préparer la contraction de texte (résumé), il faut donc d'abord apprendre à \cle{lire méthodiquement} ce roman, c'est-à-dire à dégager ce que le texte dit réellement sous ce qu'il raconte. C'est l'objet de cette section.

\courssub{Ferdinand Oyono et le contexte du roman}

\begin{definition}[Lecture méthodique]
La \cle{lecture méthodique} est une lecture organisée en étapes : on situe le texte (auteur, époque, genre), on repère la situation d'énonciation (qui parle ? à qui ? où ? quand ?), on dégage le thème et la problématique, puis on analyse les procédés (ironie, satire, registres) pour formuler les idées essentielles. Elle s'oppose à la lecture rapide et distraite.
\end{definition}

Ferdinand Oyono (1929-2010) est un écrivain et homme d'État camerounais, né à Ngoulemakong, dans la région du Sud. \emph{Le Vieux nègre et la médaille}, publié en 1956, est son deuxième roman après \emph{Une vie de boy}. À cette date, le Cameroun est encore sous administration française : le pays n'accédera à l'indépendance que le 1\textsuperscript{er} janvier 1960. Le roman se déroule donc en pleine période coloniale, dans un village fictif, Doum, où cohabitent l'administration coloniale (le commandant, les gardiens de la paix), la mission catholique (le père Vandermayer) et les villageois soumis.

Le héros, Meka, est un vieil homme qui a tout donné à la France : ses terres pour la mission, ses deux fils morts à la guerre, sa foi dans les promesses des Blancs. Un jour, l'administration décide de lui remettre une médaille pour récompenser sa « fidélité ». Cette récompense, qui devrait être un honneur, va devenir le point de départ de sa chute et de son éveil.

\begin{savaistu}[Un roman écrit avant l'indépendance]
Oyono a écrit \emph{Le Vieux nègre et la médaille} en 1956, quatre ans avant l'indépendance du Cameroun (1\textsuperscript{er} janvier 1960). Écrire une satire de la colonisation alors que celle-ci existait encore était un acte courageux : le roman a contribué, avec d'autres œuvres africaines de l'époque, à ouvrir les yeux des lecteurs sur la réalité du système colonial.
\end{savaistu}

\courssub{Lecture méthodique n°4 : la remise de la médaille}

\textbf{Situation d'énonciation.} L'extrait étudié se situe au cœur du roman : le jour de la fête du 14 juillet, à Doum. Le narrateur, extérieur à l'histoire (narration à la troisième personne), décrit la cérémonie officielle au cours de laquelle le commandant remet à Meka la médaille de l'amitié franco-africaine. Autour d'eux : les notables blancs, les missionnaires, les villageois curieux, les gardiens de la paix.

\textbf{Thème.} La cérémonie de récompense coloniale, c'est-à-dire la manière dont l'administration coloniale honore les Africains « fidèles ».

\textbf{Problématique.} En quoi cette récompense est-elle une illusion qui cache la réalité de l'exploitation coloniale ?

\textbf{Analyse.} Tout dans la scène est excessif : les discours pompeux, les uniformes, la chaleur écrasante, la sueur, les éloges exagérés. Meka, placé au centre de la fête, est présenté comme un « ami de la France ». Pourtant, le lecteur comprend vite que cette amitié est à sens unique : les Blancs célèbrent surtout leur propre image de « bons colonisateurs ». La médaille ne coûte rien à l'administration, alors que Meka, lui, a tout perdu.

\begin{exemplebox}[Repérage de l'ironie et de la satire]
L'\cle{ironie} consiste à dire le contraire de ce que l'on pense, ou à laisser voir l'écart entre les apparences et la réalité. La \cle{satire} est la critique moqueuse d'une institution ou d'un groupe. Dans cet extrait :
\begin{itemize}
\item l'administration coloniale est ridiculisée : le commandant prononce un discours grandiloquent pour récompenser un homme ruiné par la colonisation ;
\item la mission catholique est visée : le père Vandermayer bénit une cérémonie qui consacre l'inégalité entre les hommes, contrairement au message chrétien ;
\item le titre même du roman est ironique : « le vieux nègre » est l'expression méprisante des colons, et « la médaille » est un bout de métal sans valeur face aux sacrifices de Meka.
\end{itemize}
\end{exemplebox}

\textbf{Idées essentielles de l'extrait.} En dépouillant la scène de ses détails descriptifs, on obtient trois idées :
\begin{enumerate}
\item la récompense coloniale est une hypocrisie : elle masque l'exploitation ;
\item la médaille est une illusion : elle ne rend ni les terres, ni les fils, ni la dignité ;
\item la cérémonie sert les colonisateurs, non le colonisé : elle donne bonne conscience aux Blancs.
\end{enumerate}

\begin{exemplebox}[Formulation de l'idée essentielle en une phrase]
« En récompensant Meka d'une médaille sans valeur, l'administration coloniale masque hypocritement l'exploitation et les sacrifices subis par les Africains. » (20 mots)
\end{exemplebox}

\courssub{Lecture méthodique n°5 : la révolte et la chute de Meka}

\textbf{Situation d'énonciation.} Le second extrait se situe après la cérémonie. Ivre et égaré, Meka rentre chez lui dans la nuit ; surpris par une tornade, il est arrêté par les gardiens de la paix, roué de coups et jeté en prison, comme n'importe quel « indigène ». La médaille, accrochée à sa poitrine, ne le protège pas.

\textbf{Thème.} La désillusion : la découverte brutale, par Meka, de la véritable nature du système colonial.

\textbf{Problématique.} Comment la récompense se transforme-t-elle en humiliation, et que comprend alors le héros ?

\textbf{Analyse.} Le contraste est frappant : le matin, Meka était porté aux nues ; le soir, il est battu et enfermé. Ce renversement brutal révèle la vérité : pour l'administration, Meka n'a jamais été un « ami », mais un instrument de propagande. Une fois la cérémonie terminée, il redevient un colonisé parmi les autres. Cette chute provoque chez lui un \cle{éveil de la conscience} : il comprend qu'on s'est moqué de lui toute sa vie.

\begin{popfigure}
\begin{center}
\begin{tikzpicture}[scale=1, every node/.style={font=\small}]
% ligne du temps
\draw[thick, popdark, -{Stealth}] (0,0) -- (13.2,0);
% étapes
\foreach \x/\titre/\detail/\col in {
  0.9/{Invitation}/{à la fête du 14 juillet}/popblue,
  4.1/{Remise de la médaille}/{illusion de l'honneur}/popgreen,
  7.3/{Nuit d'ivresse}/{et tornade}/poporange,
  10.5/{Arrestation, prison}/{humiliation}/poppink,
  12.6/{Désillusion}/{éveil de la conscience}/poppurple}
{
  \fill[\col] (\x,0) circle (0.12);
  \draw[\col, thick] (\x,0) -- (\x,0.9);
  \node[above, align=center, text=\col] at (\x,0.95) {\textbf{\titre}\\ \detail};
}
\end{tikzpicture}
\end{center}
\legende{Frise chronologique du parcours de Meka : de l'invitation officielle à la désillusion finale. Le renversement entre la médaille (honneur apparent) et l'arrestation (réalité coloniale) est le cœur de la démonstration du roman.}
\end{popfigure}

\textbf{Idées essentielles de l'extrait.} Trois idées se dégagent :
\begin{enumerate}
\item la récompense coloniale n'accorde aucun droit réel au colonisé ;
\item la violence est la véritable relation entre l'administration et les Africains ;
\item la souffrance de Meka provoque sa prise de conscience : il cesse de croire aux promesses des Blancs.
\end{enumerate}

\begin{exoresolu}[Dégager l'idée essentielle d'un paragraphe]
Énoncé. Voici un paragraphe inspiré de l'extrait de la chute : « Les gardiens de la paix se jetèrent sur le vieil homme. La médaille rebondit sur sa poitrine à chaque coup. Personne ne reconnut en lui l'ami de la France célébré le matin même. » Formule l'idée essentielle en une phrase de moins de 25 mots.
\tcblower
Solution. On supprime les détails narratifs (les coups, la poitrine, le matin) et on garde l'opposition centrale : honneur le matin, brutalité le soir. Idée essentielle : « La médaille ne protège pas Meka : dès la cérémonie terminée, l'administration redevient violente et méprisante envers lui. » (19 mots)
\end{exoresolu}

\courssub{Les outils de liaison : conjonctions de coordination et de subordination}

\begin{definition}[Conjonctions de coordination et de subordination]
Les \cle{conjonctions de coordination} (mais, ou, et, donc, or, ni, car) relient deux éléments de même nature syntaxique (mots, groupes, propositions) sans créer de dépendance. Les \cle{conjonctions de subordination} (que, quand, comme, si, parce que, bien que, afin que, etc.) introduisent une proposition subordonnée qui dépend d'une proposition principale.
\end{definition}

Dans les extraits d'Oyono, ces outils structurent l'argumentation implicite :
\begin{itemize}
\item \textbf{Coordination (opposition) :} « Le commandant prononce un discours grandiloquent \textbf{mais} Meka est ruiné. » $\rightarrow$ \emph{mais} marque le contraste entre l'apparence et la réalité.
\item \textbf{Subordination (cause) :} « La médaille ne le protège pas \textbf{parce que} le système colonial est violent. » $\rightarrow$ \emph{parce que} explicite la cause profonde.
\item \textbf{Subordination (concession) :} « \textbf{Bien que} décoré, Meka est battu. » $\rightarrow$ \emph{bien que} souligne l'ironie du sort.
\item \textbf{Subordination (finalité) :} « L'administration organise la cérémonie \textbf{afin que} les villageois croient à sa bienveillance. » $\rightarrow$ \emph{afin que} révèle l'intention manipulatrice.
\end{itemize}

\begin{methode}[Repérer et utiliser les conjonctions pour la contraction]
\begin{enumerate}
\item Souligner toutes les conjonctions dans le paragraphe à résumer.
\item Identifier leur fonction : coordination (ajout, opposition, choix, conséquence) ou subordination (cause, conséquence, but, concession, temps, condition).
\item Garder les liens logiques essentiels (cause, conséquence, opposition) qui portent l'argumentation.
\item Supprimer les subordonnées illustratives (exemples, détails temporels) en conservant le lien principal.
\item Reformuler en utilisant des connecteurs variés pour fluidifier le résumé.
\end{enumerate}
\end{methode}

\courssub{Du contenu manifeste au contenu latent}

\begin{definition}[Contenu manifeste et contenu latent]
Le \cle{contenu manifeste} d'un texte est ce qu'il dit explicitement, à la surface : les événements, les paroles, les descriptions. Le \cle{contenu latent} est ce que le texte signifie en profondeur sans le dire directement : la critique, la dénonciation, le message implicite que le lecteur doit reconstruire à partir des procédés (ironie, contrastes, symboles).
\end{definition}

C'est ici que la lecture méthodique rejoint la section précédente sur l'argumentation : le roman d'Oyono est un texte argumentatif \emph{déguisé} en récit. Sa thèse (latente) est : « la colonisation est un système hypocrite et injuste » ; ses arguments sont les épisodes du parcours de Meka.

\begin{popfigure}
\begin{center}
\begin{tikzpicture}[font=\small]
% niveau manifeste
\node[draw=popblue, fill=popblueL, rounded corners, thick, minimum width=12.5cm, minimum height=2.1cm, align=center] (man) at (0,2.6)
{\textbf{Niveau 1 : contenu manifeste (ce que le texte raconte)}\\[2pt]
Une cérémonie officielle : le commandant décore Meka,\\
les discours célèbrent l'amitié franco-africaine, le village fête son héros.};
% niveau latent
\node[draw=poppink, fill=poppinkL, rounded corners, thick, minimum width=12.5cm, minimum height=2.1cm, align=center] (lat) at (0,-0.6)
{\textbf{Niveau 2 : contenu latent (ce que le texte dénonce)}\\[2pt]
L'hypocrisie coloniale : la médaille est un leurre qui masque\\
l'exploitation, le mépris et la violence réelle du système colonial.};
% flèche
\draw[-{Stealth}, very thick, poppurple] (man.south) -- (lat.north)
  node[midway, right, align=left, text=poppurple]{\;procédés : ironie,\\ \;contrastes, satire};
\end{tikzpicture}
\end{center}
\legende{Schéma à deux niveaux : sous le récit de la cérémonie (contenu manifeste), le roman dénonce la colonisation (contenu latent). L'ironie et les contrastes sont les passages qui permettent au lecteur de descendre d'un niveau à l'autre.}
\end{popfigure}

\courssub{Méthode : dégager l'idée essentielle pour préparer le résumé}

\begin{methode}[Dégager l'idée essentielle d'un paragraphe littéraire]
\begin{enumerate}
\item Lire le paragraphe entier et souligner le sujet principal (de qui ou de quoi parle-t-on ?).
\item Repérer les verbes et les mots qui portent le jugement ou l'opposition (mépris, illusion, hypocrisie, mais, pourtant).
\item Éliminer les détails : descriptions, dialogues, gestes, adjectifs pittoresques, exemples.
\item Se demander : « que veut montrer l'auteur à travers cette scène ? » (passage au contenu latent).
\item Formuler l'idée en une phrase simple : sujet + verbe + complément, sans recopier les mots du texte.
\item Vérifier : la phrase doit rester vraie même si l'on n'a pas lu l'extrait.
\end{enumerate}
\end{methode}

\begin{methode}[Préparer et présenter un exposé oral]
\begin{enumerate}
\item \textbf{Analyser le sujet :} dégager la problématique et le plan (introduction, 2-3 parties, conclusion).
\item \textbf{Collecter et hiérarchiser l'information :} sélectionner les idées essentielles (pas les détails), les ordonner logiquement.
\item \textbf{Rédiger les notes :} mots-clés, phrases courtes, connecteurs logiques (d'abord, ensuite, en revanche, par conséquent).
\item \textbf{Travailler l'oral :} articulation, débit, regard vers le public, appui sur les mots-clés, gestion du temps.
\item \textbf{Prévoir les questions :} anticiper 2-3 questions possibles et préparer des réponses courtes.
\end{enumerate}
\end{methode}

\begin{attention}[Ne pas confondre résumer un récit et résumer une argumentation]
Le résumé d'un texte \emph{narratif} garde la chaîne des événements (qui fait quoi, dans quel ordre) en les condensant. Le résumé d'un texte \emph{argumentatif} garde la thèse et les arguments, en supprimant les exemples et les illustrations. Or un roman comme celui d'Oyono est un récit à visée argumentative : pour le résumer correctement à l'examen, il faut dégager les \emph{idées} (hypocrisie coloniale, illusion de la récompense, éveil de la conscience) et non raconter toute l'histoire de Meka. Recopier les descriptions (la chaleur, les uniformes, la sueur) est la faute la plus fréquente : cela alourdit le résumé sans en exprimer le sens.
\end{attention}

\courssub{Activités d'intégration : reformulation et contraction}

\begin{exemplebox}[Exemple de contraction d'un paragraphe narratif en idée essentielle]
\textbf{Texte original (45 mots) :} « Le soleil brûlait la place du village. Le commandant, en uniforme blanc, s'avança vers Meka. Il lui épingla la médaille sur la poitrine en prononçant des mots chaleureux sur l'amitié éternelle entre la France et l'Afrique. La foule applaudit. »
\textbf{Idée essentielle (18 mots) :} « Lors d'une cérémonie officielle, l'administration coloniale décore Meka, célébrant une amitié franco-africaine fictive. »
\textbf{Technique :} suppression des détails sensoriels (soleil, uniforme blanc, mots chaleureux, foule), conservation du sujet, de l'action principale et de la portée ironique (fictive).
\end{exemplebox}

\begin{exercice}[À toi de jouer]
\begin{enumerate}
\item Relis la frise du parcours de Meka. Pour chaque étape, formule une idée essentielle en une phrase de moins de 20 mots.
\item Classe les phrases suivantes selon qu'elles relèvent du contenu manifeste (M) ou du contenu latent (L) : a) « Le commandant épingle la médaille sur la poitrine de Meka. » ; b) « La colonisation donne de fausses récompenses pour mieux exploiter. » ; c) « Les gardiens frappent Meka et l'enferment. » ; d) « La violence est la vraie langue du système colonial. »
\item Résume l'extrait de la remise de la médaille en trois phrases maximum, sans aucune description.
\item Dans les phrases ci-dessous, entoure les conjonctions de coordination et souligne les conjonctions de subordination : 
\begin{itemize}
\item Meka est décoré \textbf{mais} il reste un colonisé.
\item \textbf{Parce que} la médaille n'a pas de valeur, elle ne le protège pas.
\item Il a donné ses terres \textbf{et} ses fils \textbf{mais} il reçoit un bout de métal.
\end{itemize}
\item Prépare un exposé de 3 minutes sur le thème : « L'ironie comme arme anti-coloniale chez Ferdinand Oyono ». Note les grandes étapes de ton plan et les connecteurs logiques que tu utiliseras.
\end{enumerate}
\end{exercice}

\begin{corrige}[Exercice]
1. Invitation : l'administration choisit Meka comme vitrine de sa politique. Médaille : la récompense est un honneur illusoire qui ne coûte rien aux colons. Arrestation : la médaille ne protège pas Meka de la violence coloniale. Désillusion : Meka comprend l'hypocrisie du système et sa conscience s'éveille.
2. a) M ; b) L ; c) M ; d) L.
3. Exemple de résumé : « L'administration coloniale décore Meka pour sa fidélité. Cette récompense, sans valeur réelle, masque l'exploitation dont il a été victime toute sa vie. La cérémonie sert surtout à donner bonne conscience aux colonisateurs. »
4. Coordination : \textbf{mais}, \textbf{et}, \textbf{mais}. Subordination : \textbf{Parce que}.
5. Plan suggéré : Introduction (accroche, présentation de l'œuvre, problématique) ; I. L'ironie situationnelle : la médaille comme symbole creux (connecteurs : tout d'abord, en effet) ; II. L'ironie dramatique : le lecteur sait, pas le héros (connecteurs : ensuite, de surcroît) ; III. La satire des institutions : administration et mission (connecteurs : enfin, par conséquent) ; Conclusion (bilan, ouverture).
\end{corrige}

\begin{aretenir}[L'essentiel de la section]
\begin{itemize}
\item \emph{Le Vieux nègre et la médaille} (1956) de Ferdinand Oyono dénonce la colonisation à travers le destin de Meka.
\item La lecture méthodique suit des étapes : situation d'énonciation, thème, problématique, procédés, idées essentielles.
\item L'ironie et la satire visent l'administration coloniale et la mission catholique.
\item Trois idées essentielles : hypocrisie coloniale, illusion de la récompense, éveil de la conscience.
\item Le contenu manifeste raconte ; le contenu latent dénonce. Un bon résumé formule les idées, il ne recopie pas les descriptions.
\item Les conjonctions de coordination (mais, ou, et, donc, or, ni, car) et de subordination (que, quand, comme, si, parce que, bien que, afin que) structurent la logique du texte ; les repérer aide à contracter sans perdre le sens.
\item L'exposé oral demande une préparation rigoureuse : problématique, plan, notes synthétiques, connecteurs, entraînement à l'oral.
\end{itemize}
\end{aretenir}

\coursec{Les techniques de réduction du texte}

Dans la section précédente, la lecture méthodique du \emph{Vieux nègre et la médaille} de Ferdinand Oyono (lectures n°4 et n°5) nous a appris à repérer les idées essentielles d'un texte argumentatif : la thèse de l'auteur, ses arguments, ses exemples, l'enchaînement logique. Mais repérer les idées ne suffit pas. À l'examen du Probatoire et du Baccalauréat technique, l'épreuve de contraction de texte demande de \cle{reformuler} ces idées en un nombre de mots imposé, en respectant la structure argumentative. C'est précisément l'objet de cette section : transformer un texte long en un texte court, sans trahir la pensée de l'auteur. Autrement dit, nous passons de la lecture analytique à l'écriture de synthèse.

\courssub{Qu'est-ce que la contraction de texte ?}

Imaginons une situation de la vie courante. Ton grand frère te raconte pendant dix minutes un match de football. Toi, tu dois résumer ce match en une phrase à ton père qui arrive : « Le Cameroun a gagné deux buts à zéro grâce à une meilleure défense. » Tu viens, sans le savoir, de faire une \cle{contraction} : tu as gardé l'essentiel (le résultat, la cause) et abandonné les détails (les occasions manquées, les remplacements, la météo). La contraction de texte est exactement ce geste, appliqué à un texte écrit et soumis à des règles strictes.

\begin{definition}[Contraction de texte et réduction]
La \cle{contraction de texte} est un exercice qui consiste à réduire un texte à une fraction imposée de sa longueur (généralement le \cle{quart}), tout en conservant l'intégralité de ses idées essentielles et de sa structure argumentative. La \cle{réduction} désigne l'ensemble des techniques (suppression, substitution, généralisation, fusion) qui permettent d'atteindre ce nombre de mots sans perdre le sens. Le texte obtenu s'appelle le \cle{résumé} ou \cle{texte contracté}.
\end{definition}

\begin{propriete}[La règle du quart]
Au Probatoire et au Baccalauréat technique, le résumé doit respecter le \cle{quart} du nombre de mots du texte original, avec une marge de tolérance de plus ou moins 10 \%. Ainsi :
\begin{itemize}
\item un texte de 500 mots donne un résumé d'environ \textbf{125 mots} (entre 112 et 138 mots) ;
\item un texte de 400 mots donne un résumé d'environ \textbf{100 mots} ;
\item un texte de 240 mots donne un résumé d'environ \textbf{60 mots}.
\end{itemize}
Le candidat doit indiquer le nombre de mots de son résumé à la fin de sa copie. Un dépassement important de la limite est sanctionné.
\end{propriete}

\begin{attention}[Compter les mots, une habitude à acquérir]
Beaucoup de candidats perdent des points non pas parce que leur résumé est mauvais, mais parce qu'il est trop long ou trop court. Entraîne-toi à compter les mots rapidement : compte les mots d'une ligne, multiplie par le nombre de lignes, puis ajuste. Les mots composés à trait d'union (« peut-être », « aujourd'hui ») comptent pour un seul mot. Les nombres écrits en chiffres comptent pour un mot.
\end{attention}

\courssub{Les quatre techniques de réduction}

Réduire un texte n'est pas couper au hasard. C'est appliquer méthodiquement quatre techniques complémentaires, que l'on peut représenter comme un entonnoir : le texte entre large en haut et ressort étroit en bas, après être passé par chaque filtre.

\begin{popfigure}
\begin{center}
\begin{tikzpicture}[
etape/.style={draw=popblue, fill=popblueL, rounded corners=2pt, align=center, font=\small, text=popdark},
etapeF/.style={draw=popgreen, fill=popgreenL, rounded corners=2pt, align=center, font=\small, text=popdark},
fleche/.style={-{Stealth[length=3mm]}, thick, poporange}
]
\node[etape, minimum width=11cm, minimum height=0.9cm] (t0) at (0,0) {\textbf{Texte original} (thèse, arguments, exemples, répétitions, connecteurs)};
\node[etape, minimum width=8.5cm, minimum height=0.9cm] (t1) at (0,-1.4) {\textbf{Suppression} : exemples, détails, répétitions, citations, énumérations};
\node[etape, minimum width=6.5cm, minimum height=0.9cm] (t2) at (0,-2.8) {\textbf{Substitution} : énumérations $\rightarrow$ termes génériques ; périphrases $\rightarrow$ mots simples};
\node[etape, minimum width=5.2cm, minimum height=0.9cm] (t3) at (0,-4.2) {\textbf{Généralisation \& fusion} : particulier $\rightarrow$ général ; phrases $\rightarrow$ phrase unique};
\node[etapeF, minimum width=3.5cm, minimum height=0.9cm] (t4) at (0,-5.6) {\textbf{Résumé au quart} (autonome, cohérent, fidèle)};
\draw[fleche] (t0) -- (t1);
\draw[fleche] (t1) -- (t2);
\draw[fleche] (t2) -- (t3);
\draw[fleche] (t3) -- (t4);
\end{tikzpicture}
\end{center}
\legende{L'entonnoir de la contraction : le texte original traverse successivement les filtres de la suppression, de la substitution, de la généralisation et de la fusion pour devenir un résumé au quart de sa longueur, en conservant la thèse, les arguments et l'articulation logique.}
\end{popfigure}

\textbf{Technique n\textsuperscript{o} 1 : la suppression.}\par
C'est la technique la plus simple et la plus rentable. Elle consiste à éliminer purement et simplement tout ce qui n'est pas indispensable à la compréhension de la pensée de l'auteur :
\begin{itemize}
\item les \cle{exemples} : ils illustrent l'argument, ils ne le constituent pas (ex. : « Par exemple, à Douala... » $\rightarrow$ suppression) ;
\item les \cle{répétitions} et les reformulations : un auteur dit souvent deux fois la même idée sous des formes différentes ;
\item les \cle{détails} descriptifs ou narratifs : dates précises, noms propres secondaires, descriptions pittoresques ;
\item les \cle{énumérations} longues : on les supprime ou on les remplace (voir la substitution) ;
\item les \cle{citations} et dialogues : les paroles rapportées disparaissent du résumé ou sont rapportées indirectement.
\end{itemize}

\textbf{Technique n\textsuperscript{o} 2 : la substitution.}\par
Elle consiste à remplacer un groupe de mots par un seul mot de sens équivalent. Le cas le plus fréquent est le remplacement d'une énumération par un \cle{terme générique} (hyperonyme) : « mangues, oranges, bananes et ananas » devient « fruits » ; « haches, machettes et pioches » devient « outils ». On peut aussi remplacer une expression longue par un mot unique : « d'une manière rapide » devient « rapidement » ; « le fait que » devient « que ».

\textbf{Technique n\textsuperscript{o} 3 : la généralisation.}\par
Proche de la substitution, elle consiste à passer du \cle{particulier} au \cle{général}. Si le texte dit : « À Douala, à Yaoundé et à Bafoussam, les embouteillages paralysent les habitants », on généralise : « Dans les grandes villes camerounaises, les embouteillages paralysent les habitants. » L'idée est conservée, la précision géographique inutile est sacrifiée. Cette technique est cruciale pour contracter les exemples locaux en arguments généraux.

\textbf{Technique n\textsuperscript{o} 4 : la fusion.}\par
Elle consiste à regrouper deux ou trois phrases en une seule, souvent grâce à une conjonction de subordination, une relative ou un participe présent. Par exemple : « Le paludisme tue beaucoup d'enfants. Il sévit surtout en saison des pluies. » devient : « Le paludisme, qui sévit surtout en saison des pluies, tue beaucoup d'enfants. » Deux phrases de 13 mots deviennent une phrase de 11 mots, sans perte d'information. La fusion préserve les \cle{articulations logiques} (cause, conséquence, opposition, concession) grâce aux connecteurs.

\begin{exemplebox}[Tableau comparatif : phrase originale et phrase réduite]
\begin{center}
\begin{tabularx}{\linewidth}{|X|X|c|c|}
\hline
\rowcolor{popblueL} \textbf{Phrase originale} & \textbf{Phrase réduite} & \textbf{Avant} & \textbf{Après} \\
\hline
Le vieux Meka, ému, les larmes aux yeux, reçut des mains du commandant une médaille brillante. & Meka reçut ému une médaille du commandant. & 17 mots & 8 mots \\
\hline
Les mangues, les oranges, les bananes et les ananas abondent sur les marchés de la région du Centre. & Les fruits abondent sur les marchés du Centre. & 17 mots & 8 mots \\
\hline
La déforestation progresse rapidement. Elle détruit l'habitat des animaux sauvages. & La déforestation rapide détruit l'habitat des animaux sauvages. & 12 mots & 9 mots \\
\hline
\end{tabularx}
\end{center}
\legende{Application combinée des quatre techniques : suppression des détails (larmes aux yeux), substitution (fruits), généralisation (région du Centre), fusion (deux phrases en une).}
\end{exemplebox}

\courssub{Ce qu'il faut conserver et ce qu'il faut éviter}

\begin{aretenir}[Les quatre éléments à conserver obligatoirement]
\begin{enumerate}
\item La \cle{thèse} de l'auteur : c'est l'idée directrice, le cœur du texte argumentatif.
\item Les \cle{arguments} : ils justifient la thèse ; les supprimer revient à vider le texte de sa substance.
\item Les \cle{articulations logiques} : les connecteurs (mais, donc, cependant, en effet, par conséquent) qui montrent l'enchaînement des idées.
\item Le \cle{ton} de l'auteur : si le texte est indigné, ironique ou admiratif, le résumé doit le laisser sentir par le choix du vocabulaire.
\end{enumerate}
\end{aretenir}

\begin{attention}[Ce qu'il faut absolument éviter]
\begin{itemize}
\item Les \cle{citations} et les dialogues : ils doivent disparaître ou être rapportés indirectement.
\item Les \cle{exemples} : ils sont les premiers sacrifiés.
\item Les \cle{jugements personnels} : le résumé n'est pas un commentaire ; n'écris jamais « je pense que l'auteur a raison ».
\item La \cle{première personne} : on ne dit pas « je résume ce texte », on rédige directement le résumé à la troisième personne.
\item Le \cle{recopiage} : recopier des phrases entières du texte au lieu de reformuler est l'erreur la plus fréquente et la plus sanctionnée. Le correcteur veut voir TES phrases, pas celles de l'auteur. Reformuler, c'est prouver que tu as compris.
\end{itemize}
\end{attention}

\begin{propriete}[Les trois qualités d'un bon résumé]
Un résumé réussi est :
\begin{enumerate}
\item \cle{autonome} : il se lit et se comprend seul, sans avoir le texte original sous les yeux ;
\item \cle{cohérent} : les phrases s'enchaînent logiquement grâce aux connecteurs (subordination, coordination) ;
\item \cle{fidèle} : il restitue exactement la pensée de l'auteur, sans rien ajouter ni déformer.
\end{enumerate}
\end{propriete}

\courssub{Méthode de travail en cinq étapes}

\begin{methode}[Réussir une contraction de texte en 5 étapes]
\begin{enumerate}
\item \textbf{Lire} le texte deux fois : une première fois pour le sens global, une deuxième fois en identifiant la thèse, le plan argumentatif et les connecteurs logiques.
\item \textbf{Souligner} les idées essentielles : thèse, arguments, connecteurs logiques. Ignorer exemples, détails, citations. Utiliser un code couleur : thèse (rouge), arguments (bleu), connecteurs (vert).
\item \textbf{Hiérarchiser} : classer les idées de la plus importante à la moins importante, afin de savoir quoi sacrifier si le compte de mots l'exige.
\item \textbf{Réduire} en appliquant les quatre techniques : suppression, substitution, généralisation, fusion. Compter les mots régulièrement (tous les 2-3 phrases).
\item \textbf{Reformuler} : rédiger le résumé final avec ses propres mots, vérifier la cohérence (relecture à voix haute), compter les mots exactement et indiquer leur nombre entre parenthèses à la fin.
\end{enumerate}
\end{methode}

\courssub{Entraînement guidé : du texte au résumé}

\begin{exoresolu}[Réduire un paragraphe de 80 mots à 20 mots]
\textbf{Énoncé.} Réduis le paragraphe suivant (80 mots) à environ 20 mots, en conservant l'idée essentielle et la structure argumentative.

« Le paludisme demeure au Cameroun un problème majeur de santé publique. Chaque année, cette maladie, transmise par la piqûre d'un moustique appelé anophèle, frappe des milliers de personnes, en particulier les enfants de moins de cinq ans et les femmes enceintes. Par exemple, dans la seule région du Centre, on a enregistré l'an dernier des centaines de cas graves. Pourtant, des moyens simples existent : les moustiquaires imprégnées, la destruction des eaux stagnantes, le traitement rapide des fièvres. »
\tcblower
\textbf{Solution détaillée.}

\textbf{Étape 1 : repérage.} Thèse : le paludisme est un grave problème de santé au Cameroun. Argument 1 : il touche surtout les personnes fragiles. Argument 2 (concession/opposition : \emph{Pourtant}) : des moyens de prévention simples existent. Connecteurs : \emph{en particulier}, \emph{Par exemple}, \emph{Pourtant}.

\textbf{Étape 2 : suppression.} On élimine l'exemple (« dans la seule région du Centre... »), les détails techniques (« appelé anophèle »), les précisions temporelles (« l'an dernier »), les quantificatifs vagues (« des centaines de cas »).

\textbf{Étape 3 : substitution et généralisation.} « les enfants de moins de cinq ans et les femmes enceintes » $\rightarrow$ « les personnes vulnérables » (généralisation) ; « les moustiquaires imprégnées, la destruction des eaux stagnantes, le traitement rapide des fièvres » $\rightarrow$ « des mesures de prévention simples » (substitution + généralisation).

\textbf{Étape 4 : fusion et rédaction.}

\emph{Grave problème de santé publique au Cameroun, le paludisme frappe surtout les personnes vulnérables, alors que des mesures de prévention simples existent.} (20 mots)

Le sens est intégralement conservé : thèse (problème grave), argument 1 (victimes vulnérables), argument 2 avec opposition (prévention possible), articulation logique (\emph{alors que}).
\end{exoresolu}

\begin{exercice}[À toi de jouer]
Réduis le passage suivant (60 mots) à environ 15 mots, en conservant la thèse, l'argumentation et le connecteur logique : « L'éducation des jeunes filles est une condition indispensable du développement de notre pays. En effet, une femme instruite prend soin de la santé de ses enfants, les envoie à l'école, gère mieux les ressources de la famille et participe activement à la vie économique de son village ou de son quartier. »
\end{exercice}

\begin{corrige}[Exercice : l'éducation des jeunes filles]
\textbf{Analyse.} Thèse : l'éducation des filles est indispensable au développement. Connecteur : \emph{En effet} (explication). Argument unique développé en quatre points (santé, scolarisation, gestion, économie).

\textbf{Réduction.} Suppression de « En effet » (simple annonce), fusion des quatre bénéfices en un terme générique par généralisation.

\emph{L'éducation des jeunes filles, condition du développement national, profite à toute la famille et à la communauté.} (15 mots)
\end{corrige}

\begin{savaistu}[Le résumé, une compétence professionnelle]
La contraction de texte n'est pas qu'un exercice scolaire. Dans la vie professionnelle — secrétariat, administration, journalisme, compte rendu de réunion, rapport de stage ou de chantier —, savoir résumer fidèlement un document long en quelques lignes est une compétence très recherchée. Le technicien qui rédige une note de synthèse ou un rapport d'activité pratique exactement les mêmes techniques : suppression des détails, généralisation des données, fusion des informations, reformulation en langage propre.
\end{savaistu}

En résumé, la contraction repose sur un principe simple : garder l'essentiel (thèse, arguments, articulation), sacrifier l'accessoire (exemples, détails, répétitions), et tout dire avec ses propres mots. Les quatre techniques — suppression, substitution, généralisation, fusion — sont tes outils ; la règle du quart est ta contrainte ; la fidélité à la pensée de l'auteur est ta boussole. La section suivante nous permettra d'étudier les outils grammaticaux qui rendent cette reformulation possible : les conjonctions de coordination et de subordination, indispensables à la fusion des phrases et à la cohérence du résumé.

\coursec{Les outils de liaison : conjonctions de coordination et de subordination}

Dans la section précédente, nous avons appris à \cle{réduire} un texte : supprimer les répétitions, les exemples, les détails inutiles. Mais supprimer ne suffit pas. Pour \emph{fusionner} des phrases sans perdre le sens, il faut des outils capables d'exprimer en un seul mot des relations entières : la cause, la conséquence, l'opposition, le but. Ces outils, ce sont les \cle{conjonctions}. Bien les connaître, c'est gagner des mots au résumé et des points à l'épreuve de contraction de texte.

\courssub{Rappel : phrase simple, phrase complexe, proposition}

\begin{definition}[Phrase simple et phrase complexe]
Une \cle{phrase simple} ne contient qu'un seul verbe conjugué : elle est formée d'une seule \cle{proposition}. Une \cle{phrase complexe} contient au moins deux verbes conjugués, donc au moins deux propositions.
\end{definition}

Une \cle{proposition} est un groupe de mots organisé autour d'un verbe conjugué. C'est l'unité de base de la phrase complexe.

\begin{exemplebox}[Observer pour comprendre]
\begin{enumerate}
\item \emph{Meka reçoit une médaille.} $\rightarrow$ 1 verbe conjugué (\emph{reçoit}) : phrase \textbf{simple}.
\item \emph{Meka reçoit une médaille car il a donné ses fils.} $\rightarrow$ 2 verbes conjugués (\emph{reçoit}, \emph{a donné}) : phrase \textbf{complexe} à deux propositions.
\end{enumerate}
\end{exemplebox}

Lorsqu'une phrase complexe contient plusieurs propositions, celles-ci peuvent être liées de deux grandes manières : la \cle{coordination} et la \cle{subordination}. C'est toute la matière de cette section.

\courssub{La coordination : juxtaposition et conjonctions de coordination}

\textbf{La juxtaposition et la coordination.}\par
Deux propositions sont \cle{juxtaposées} lorsqu'elles sont simplement placées l'une à côté de l'autre, séparées par une virgule ou un point-virgule, sans mot de liaison : \emph{Meka vieillit, ses forces l'abandonnent.}

Deux propositions sont \cle{coordonnées} lorsqu'elles sont reliées par une conjonction de coordination. Elles restent grammaticalement \emph{indépendantes} : chacune pourrait exister seule.

\begin{definition}[Conjonction de coordination]
Une \cle{conjonction de coordination} est un mot invariable qui relie deux mots, deux groupes de mots ou deux propositions de \textbf{même fonction grammaticale}. Il en existe sept : \textbf{mais, ou, et, donc, or, ni, car} (moyen mnémotechnique : \emph{« Mais où est donc Ornicar ? »}).
\end{definition}

\begin{aretenir}[Sens des sept conjonctions de coordination]
\begin{center}
\begin{tabularx}{\linewidth}{|l|l|X|}
\hline
\rowcolor{popblueL} \textbf{Conjonction} & \textbf{Relation exprimée} & \textbf{Exemple} \\
\hline
\cle{et} & addition, énumération & Meka a donné ses fils \textbf{et} sa terre. \\
\hline
\cle{ou} & alternative, choix & Obéiras-tu \textbf{ou} résisteras-tu ? \\
\hline
\cle{mais} & opposition, restriction & Il est vieux, \textbf{mais} il reste digne. \\
\hline
\cle{donc} & conséquence, conclusion & Il a tout perdu, \textbf{donc} il mérite récompense. \\
\hline
\cle{or} & passage à un argument nouveau & Les Blancs le flattent ; \textbf{or} ils le méprisent. \\
\hline
\cle{ni} & addition négative & Il n'a \textbf{ni} terre \textbf{ni} fils. \\
\hline
\cle{car} & cause, justification & Il est honoré \textbf{car} il a beaucoup sacrifié. \\
\hline
\end{tabularx}
\end{center}
\end{aretenir}

\courssub{La subordination : conjonctions et propositions subordonnées}

\begin{definition}[Conjonction de subordination]
Une \cle{conjonction de subordination} est un mot invariable qui introduit une proposition \textbf{dépendante} d'une proposition principale : la proposition subordonnée ne peut pas exister seule, elle \emph{complète} la principale. Les plus fréquentes sont : \textbf{que, parce que, puisque, comme, bien que, quoique, si, lorsque, quand, afin que, pour que}.
\end{definition}

La proposition qui dépend est appelée \cle{proposition subordonnée conjonctive} ; celle dont elle dépend est la \cle{proposition principale}. On distingue :

\begin{itemize}
\item les subordonnées \textbf{complétives}, introduites par \emph{que}, qui complètent un verbe : \emph{Meka comprend \textbf{que} la médaille ne nourrit pas.} (fonction : COD du verbe \emph{comprend}) ;
\item les subordonnées \textbf{circonstancielles}, qui précisent une circonstance (cause, temps, but, condition, concession) : \emph{On l'honore \textbf{parce qu'}il a donné ses fils.} ;
\item les subordonnées \textbf{relatives} (rappel), introduites par un pronom relatif (\emph{qui, que, où, dont}) et non par une conjonction : \emph{la médaille \textbf{qu'}on lui remet}.
\end{itemize}

\begin{methode}[Distinguer proposition coordonnée et proposition subordonnée : le test de suppression]
\begin{enumerate}
\item Repérer le mot de liaison entre les deux propositions.
\item \textbf{Supprimer} la seconde proposition : la phrase restante a-t-elle encore un sens complet ?
\item Si \textbf{oui}, les propositions sont indépendantes : c'est de la \textbf{coordination}. \emph{Meka reçoit une médaille car il a donné ses fils} $\rightarrow$ \emph{Meka reçoit une médaille} : phrase complète.
\item Si \textbf{non} (la phrase reste en suspens), c'est de la \textbf{subordination}. \emph{Meka comprend que la médaille ne nourrit pas} $\rightarrow$ \emph{Meka comprend\ldots} : incomplet, il manque le complément.
\item Vérifier aussi la place : une subordonnée circonstancielle peut souvent être \textbf{déplacée} (\emph{Parce qu'il a donné ses fils, on l'honore}), ce qu'une coordonnée après \emph{car} ne peut pas.
\end{enumerate}
\end{methode}

\begin{attention}[Ne pas confondre \emph{car} et \emph{parce que}]
Les deux expriment la cause, mais leur nature grammaticale est différente : \cle{car} est une conjonction de \textbf{coordination} (il relie deux propositions indépendantes), tandis que \cle{parce que} est une conjonction de \textbf{subordination} (il introduit une proposition dépendante). Conséquences pratiques : on peut commencer une phrase par \emph{Parce que\ldots}, jamais par \emph{Car\ldots} en liaison interne ; et à l'écrit, \emph{car} est plus soutenu. Au résumé, les deux sont précieux car ils condensent une cause en un seul mot.
\end{attention}

\begin{exemplebox}[Analyse grammaticale complète]
Phrase : \emph{\textbf{Meka reçoit une médaille car il a donné ses fils et sa terre.}}
\begin{itemize}
\item \textbf{P1} : \emph{Meka reçoit une médaille} — proposition indépendante (sujet : \emph{Meka} ; verbe : \emph{reçoit} ; COD : \emph{une médaille}).
\item \textbf{P2} : \emph{il a donné ses fils et sa terre} — proposition \textbf{coordonnée} à P1 par \emph{car} (elle exprime la \textbf{cause}).
\item À l'intérieur de P2, \emph{et} coordonne les deux COD : \emph{ses fils} / \emph{sa terre}.
\end{itemize}
Test de suppression : \emph{Meka reçoit une médaille} se comprend seul $\rightarrow$ coordination confirmée.
\end{exemplebox}

\courssub{Vue d'ensemble : l'arbre des liaisons}

\begin{popfigure}
\begin{tikzpicture}[
  grow=down,
  level distance=1.5cm,
  level 1/.style={sibling distance=5cm},
  level 2/.style={sibling distance=2.5cm},
  every node/.style={font=\small},
  racine/.style={rectangle, rounded corners, draw=popdark, fill=popblueL, font=\bfseries, align=center, text width=3.5cm},
  branche/.style={rectangle, rounded corners, draw=popblue, fill=popblue!15, font=\bfseries, align=center, text width=3cm},
  feuille/.style={rectangle, rounded corners, draw=popmuted, fill=popcream, align=center, font=\footnotesize, text width=2.5cm},
  edge from parent/.style={draw=popmuted, thick}
]
\node[racine, align=center]{Phrase complexe\\ (plusieurs propositions)}
  child { node[branche, align=center]{Coordination / juxtaposition\\ (propositions indépendantes)}
    child { node[feuille]{addition : \textbf{et, ni}} }
    child { node[feuille]{alternative : \textbf{ou}} }
    child { node[feuille]{opposition : \textbf{mais}} }
    child { node[feuille]{conséquence : \textbf{donc}} }
    child { node[feuille]{argument : \textbf{or}} }
    child { node[feuille]{cause : \textbf{car}} }
  }
  child { node[branche, align=center]{Subordination\\ (proposition dépendante)}
    child { node[feuille]{cause : \textbf{parce que, puisque, comme}} }
    child { node[feuille]{temps : \textbf{quand, lorsque, dès que}} }
    child { node[feuille]{but : \textbf{pour que, afin que}} }
    child { node[feuille]{condition : \textbf{si}} }
    child { node[feuille]{concession : \textbf{bien que, quoique}} }
    child { node[feuille]{complétive : \textbf{que}} }
  };
\end{tikzpicture}
\legende{Arbre de classement des liaisons dans la phrase complexe : à gauche, les sept conjonctions de coordination classées par sens ; à droite, les principales conjonctions de subordination classées par relation logique.}
\end{popfigure}

\courssub{Les relations logiques : tableau de synthèse}

Pour le résumé, il faut raisonner en termes de \cle{relations logiques} : à chaque relation correspondent plusieurs outils, qu'il faut savoir choisir selon le niveau de langue et l'économie de mots.

\begin{aretenir}[Relations logiques et connecteurs associés]
\begin{center}
\begin{tabularx}{\linewidth}{|l|X|X|}
\hline
\rowcolor{popblueL} \textbf{Relation} & \textbf{Conjonctions} & \textbf{Autres connecteurs utiles} \\
\hline
Cause & car ; parce que, puisque, comme & c'est pourquoi (effet), grâce à, à cause de \\
\hline
Conséquence & donc & alors, ainsi, par conséquent, c'est pourquoi \\
\hline
Opposition & mais & cependant, pourtant, en revanche, alors que, tandis que \\
\hline
Concession & bien que, quoique, même si & malgré (+ nom), certes\ldots mais \\
\hline
But & pour que, afin que (+ subjonctif) & pour, afin de (+ infinitif) \\
\hline
Temps & quand, lorsque, dès que, pendant que & puis, ensuite, enfin, d'abord \\
\hline
Condition & si, à condition que & à moins que, dans le cas où \\
\hline
\end{tabularx}
\end{center}
\end{aretenir}

\begin{attention}[Pièges fréquents]
\begin{itemize}
\item Après \emph{bien que, quoique, afin que, pour que, à condition que}, le verbe est au \textbf{subjonctif} : \emph{bien qu'il \textbf{soit} vieux} (et non \emph{bien qu'il est vieux}).
\item \emph{Si} exprimant la condition est suivi de l'imparfait, jamais du conditionnel : \emph{s'il \textbf{avait}} (et non \emph{s'il aurait}).
\item Ne pas multiplier les connecteurs dans le résumé : un connecteur juste vaut mieux que trois connecteurs approximatifs.
\end{itemize}
\end{attention}

\courssub{Le rôle des conjonctions dans le résumé : dire plus avec moins}

Voilà le lien direct avec la contraction de texte. Une conjonction bien choisie permet de \textbf{fusionner deux phrases en une seule}, donc de gagner des mots tout en \emph{conservant la relation logique} du texte de départ. C'est exactement ce qu'attend le correcteur : un résumé fidèle \emph{et} cohérent.

\begin{exoresolu}[Fusion de phrases pour gagner des mots]
\textbf{Énoncé.} Texte de départ (36 mots) : \emph{Meka est très vieux. Il a donné ses deux fils à la guerre. Il a aussi donné sa terre à la mission. Le commandant lui remet une médaille. Cette médaille ne le nourrit pas. Meka reste pauvre.} Fusionnez ces six phrases en deux phrases complexes à l'aide de conjonctions, sans perdre aucune idée essentielle.
\tcblower
\textbf{Solution détaillée.}
\begin{enumerate}
\item On repère les relations logiques : \emph{vieillissement + dons} (addition) ; \emph{dons $\rightarrow$ médaille} (cause) ; \emph{médaille $\rightarrow$ pauvreté} (opposition).
\item On choisit les connecteurs : \emph{et} (addition), \emph{car} ou \emph{parce que} (cause), \emph{mais} (opposition).
\item Rédaction (24 mots) : \emph{Meka, très vieux, reçoit une médaille \textbf{car} il a donné ses fils \textbf{et} sa terre, \textbf{mais} cette récompense ne le nourrit pas : il reste pauvre.}
\end{enumerate}
Bilan : 36 mots $\rightarrow$ 24 mots, soit un tiers gagné, sans perte d'idée et avec des relations logiques explicites. C'est la technique de base de la contraction.
\end{exoresolu}

\begin{exercice}[À vous de jouer]
Fusionnez chaque groupe de phrases en une seule phrase complexe, en choisissant la conjonction adaptée à la relation logique.
\begin{enumerate}
\item Le vieux Meka marche jusqu'à la ville. Le chemin est long et pénible. (opposition ou concession)
\item Les villageois admirent la médaille. Ils ne comprennent pas sa valeur réelle. (opposition)
\item Meka a perdu ses fils. Meka a perdu sa terre. (addition négative avec \emph{ni})
\item On organise une grande cérémonie. On veut honorer le vieillard devant tout le village. (but)
\end{enumerate}
\end{exercice}

\begin{corrige}[Exercice : fusion de phrases]
\begin{enumerate}
\item \emph{Bien que} le chemin \textbf{soit} long et pénible, le vieux Meka marche jusqu'à la ville. (concession, subjonctif après \emph{bien que}) — ou : Le chemin est long et pénible, \emph{mais} Meka marche jusqu'à la ville.
\item Les villageois admirent la médaille, \emph{mais} ils ne comprennent pas sa valeur réelle. (opposition, coordination)
\item Meka n'a \emph{ni} fils \emph{ni} terre. (addition négative, coordination de deux COD)
\item On organise une grande cérémonie \emph{afin d'}honorer le vieillard devant tout le village. (but ; \emph{afin de} + infinitif car le sujet est le même ; sinon : \emph{afin que le village honore\ldots})
\end{enumerate}
\end{corrige}

\begin{savaistu}[Les connecteurs et la note de langue au Bac]
Au Baccalauréat technique, les barèmes de contraction de texte et de dissertation réservent une part importante de la note à la \textbf{cohérence textuelle} : enchaînement logique des idées, correction de la langue, précision des connecteurs. Un résumé où chaque phrase est reliée par le bon connecteur (\emph{car} pour la cause, \emph{mais} pour l'opposition, \emph{donc} pour la conséquence) peut gagner jusqu'à deux points par rapport à une suite de phrases juxtaposées sans lien. À l'inverse, un connecteur mal choisi (par exemple \emph{donc} à la place de \emph{mais}) déforme la pensée de l'auteur : c'est une faute de fidélité, pénalisée deux fois, en langue et en contenu. Maîtriser sept conjonctions de coordination et une dizaine de conjonctions de subordination est donc l'un des investissements les plus rentables de l'année de Terminale.
\end{savaistu}

\begin{aretenir}[L'essentiel de la section]
\begin{itemize}
\item Phrase simple = 1 verbe conjugué ; phrase complexe = plusieurs propositions.
\item Coordination : propositions \textbf{indépendantes} reliées par \emph{mais, ou, et, donc, or, ni, car}.
\item Subordination : proposition \textbf{dépendante} introduite par \emph{que, parce que, bien que, si, lorsque, afin que\ldots}
\item Test de suppression : si la phrase restante garde un sens complet, c'est de la coordination.
\item \emph{Car} = coordination ; \emph{parce que} = subordination.
\item Dans le résumé, chaque conjonction exprime une relation logique (cause, conséquence, opposition, but, temps, condition) et permet de fusionner des phrases pour gagner des mots.
\end{itemize}
\end{aretenir}

\coursec{Contenus manifestes et contenus latents}

Dans la section précédente, nous avons étudié les outils de liaison — conjonctions de coordination et de subordination — qui relient les idées entre elles de manière \emph{visible} dans la phrase. Mais tout ce qu'un texte communique n'est pas écrit noir sur blanc. Les conjonctions relient des idées \emph{exprimées} ; or un discours véhicule aussi des idées \emph{non exprimées}, que le lecteur doit reconstruire. C'est précisément l'objet de cette section : apprendre à distinguer ce qui est dit de ce qui est suggéré, compétence indispensable pour comprendre un texte en profondeur et pour le résumer sans le trahir.

\courssub{Définitions : le dit et le non-dit}

Commençons par une situation de la vie quotidienne. Un élève rentre tard au lycée. Le surveillant général lui dit : « \emph{Te voilà enfin ! Le soleil s'est levé depuis longtemps, toi tu arrives maintenant. Très bien, très bien...} » Pris au premier degré, le surveillant félicite l'élève. Mais tout le monde comprend qu'il le critique. Le message réel n'est pas dans les mots : il est \emph{derrière} les mots. Cette expérience banale montre qu'une communication comporte toujours deux niveaux.

\begin{definition}[Contenu manifeste et contenu latent]
Le \cle{contenu manifeste} d'un message est ce qui est dit \textbf{explicitement}, de façon directe et observable : les mots, les phrases, les informations clairement énoncées. C'est le sens de surface, accessible à une lecture littérale.\par
Le \cle{contenu latent} est ce qui est \textbf{suggéré, sous-entendu, implicite} : les intentions cachées, les jugements non formulés, les critiques indirectes, les valeurs que le texte défend ou dénonce sans les nommer. C'est le sens profond, que le lecteur reconstruit à partir d'indices.
\end{definition}

On peut comparer un texte à un iceberg : la partie visible au-dessus de l'eau est petite ; la partie immergée, invisible, est beaucoup plus grande. De même, dans un discours habile, le contenu latent est souvent plus riche et plus important que le contenu manifeste.

\begin{popfigure}
\begin{center}
\begin{tikzpicture}[scale=1]
% Eau
\fill[popblueL] (-5.2,0) rectangle (5.2,-3.6);
\draw[popblue, thick] (-5.2,0) -- (5.2,0);
\node[popdark, font=\small] at (4.1,0.28) {surface (lecture littérale)};
% Iceberg émergé
\fill[white, draw=popline, thick] (-1.1,0) -- (-0.7,1.5) -- (0,2.1) -- (0.8,1.4) -- (1.2,0) -- cycle;
% Iceberg immergé
\fill[popblue!40, draw=popblue, thick] (-1.2,0) -- (-2.2,-1.6) -- (-1.4,-3.2) -- (0.2,-3.5) -- (1.8,-2.4) -- (2.1,-0.9) -- (1.2,0) -- cycle;
% Étiquettes
\node[align=center, font=\small\bfseries, popdark] at (0.05,1.15) {Contenu\\manifeste};
\node[align=center, font=\small, popdark] at (0.05,0.55) {(ce qui est dit)};
\node[align=center, font=\small\bfseries, white] at (0,-1.9) {Contenu latent};
\node[align=center, font=\small, popdark] at (0,-2.7) {(ce qui est suggéré : ironie,\\implicite, présupposés, non-dits)};
% Flèches
\draw[-{Stealth}, poporange, thick] (2.6,1.2) -- (1.0,1.1);
\node[poporange, font=\small, align=left] at (3.9,1.3) {lecture\\superficielle};
\draw[-{Stealth}, poporange, thick] (3.4,-1.9) -- (1.6,-1.9);
\node[poporange, font=\small, align=left] at (4.35,-1.9) {lecture\\approfondie};
\end{tikzpicture}
\end{center}
\legende{Le modèle de l'iceberg : le contenu manifeste est la partie visible du message ; le contenu latent, plus vaste, exige un travail de décodage.}
\end{popfigure}

\begin{aretenir}[L'essentiel]
Comprendre un texte, ce n'est pas seulement comprendre les mots : c'est saisir ce que l'auteur \textbf{veut dire} au-delà de ce qu'il \textbf{dit}. Le contenu manifeste se lit ; le contenu latent se \cle{décode}.
\end{aretenir}

\courssub{Les indices du contenu latent}

Le contenu latent n'est pas deviné au hasard : il est signalé par des \cle{indices} précis que le lecteur doit apprendre à repérer.

\textbf{1. L'ironie.} Dire le contraire de ce que l'on pense, pour mieux critiquer. Quand le surveillant dit « Très bien, très bien... » à l'élève en retard, l'écart entre les mots (éloge) et la situation (faute) signale l'ironie.

\textbf{2. L'implicite.} Une information que le lecteur doit déduire lui-même. « Meka a encore salué le commandant avec un grand sourire » implique, selon le contexte, que ce sourire est forcé ou hypocrite.

\textbf{3. Les présupposés.} Ce qu'une phrase admet comme vrai sans le dire. « As-tu \emph{enfin} fini tes devoirs ? » présuppose que tu es lent et que tu tardes toujours. Le petit mot « enfin » porte tout le jugement.

\textbf{4. Les connotations.} Les valeurs affectives attachées aux mots. Dire d'un homme qu'il est « un fonctionnaire zélé » ou « un domestique docile » ne produit pas le même effet : les mots choisis révèlent le regard de l'énonciateur.

\textbf{5. Les non-dits.} Ce que le texte tait délibérément : un silence sur les souffrances, sur les humiliations, sur les résistances. Ce qui n'est pas dit peut être aussi significatif que ce qui est dit.

\courssub{Le rôle du contexte}

Un même énoncé peut changer de sens selon le \cle{contexte}. Pour décoder correctement le contenu latent, il faut interroger quatre éléments :

\begin{itemize}
\item \textbf{L'époque} : un discours prononcé pendant la colonisation ne se lit pas comme un discours d'aujourd'hui ;
\item \textbf{L'auteur} : qui parle ? un colon, un colonisé, un narrateur ironique ? sa position sociale éclaire ses intentions ;
\item \textbf{Le destinataire} : à qui s'adresse le message ? on ne dit pas la même chose à un supérieur, à un ami ou au public ;
\item \textbf{La situation de communication} : où, quand, dans quelles circonstances le message est-il produit ?
\end{itemize}

\begin{attention}[Piège fréquent]
Décoder l'implicite ne signifie pas inventer. Une \textbf{interprétation fondée} s'appuie sur des indices présents dans le texte (mots, tons, répétitions, contexte) ; une \textbf{opinion personnelle} projette sur le texte des idées qui n'y sont pas. Au Probatoire et au Baccalauréat technique, toute lecture du contenu latent doit être \textbf{justifiée par le texte}. Dire « l'auteur dénonce la colonisation » sans citer aucun indice est une surinterprétation ; montrer que le discours du commandant emploie des mots paternalistes (« nos enfants », « la bonne œuvre de la France ») est une analyse fondée.
\end{attention}

\courssub{Application au \emph{Vieux nègre et la médaille}}

Dans le roman de Ferdinand Oyono, cette distinction est capitale. Le discours officiel des colonisateurs — le commandant, le père Vandermayer, les invités de la cérémonie — affiche un contenu manifeste généreux : la France « récompense » Meka, la médaille couronne « une vie de dévouement », la colonisation est présentée comme une œuvre civilisatrice. Mais le contenu latent, que le lecteur reconstruit grâce à l'ironie du narrateur, aux détails humiliants et aux silences du texte, dénonce la domination coloniale : la médaille est une récompense dérisoire pour une vie spoliée ; la « reconnaissance » masque l'exploitation ; la fête officielle cache le mépris.

\begin{center}
\begin{tabularx}{\linewidth}{|X|X|}
\hline
\rowcolor{popblueL} \textbf{Ce qui est dit (contenu manifeste)} & \textbf{Ce qui est dénoncé (contenu latent)} \\
\hline
« Meka est un bon nègre, fidèle et dévoué. » & Le paternalisme colonial : l'Africain est infantilisé, jugé comme un serviteur, non comme un homme. \\
\hline
La médaille « récompense une vie de dévouement à la France ». & Une médaille ne compense ni les terres perdues, ni les fils morts à la guerre, ni une vie de soumission. \\
\hline
La cérémonie est grandiose : discours, drapeaux, invités. & La mise en scène sert la propagande coloniale, non le bonheur de Meka, qui reste un objet d'exposition. \\
\hline
Le commandant embrasse Meka devant les photographes. & Geste calculé pour l'image ; l'affection affichée dissimule le mépris réel. \\
\hline
Meka, ivre, est arrêté et maltraité le soir même. & Le contraste révèle la vérité : la médaille ne change rien à la condition du colonisé. \\
\hline
\end{tabularx}
\end{center}

Ce tableau montre la méthode : à gauche, on cite ou on reformule fidèlement ce que le texte dit ; à droite, on explicite le sens profond en s'appuyant sur les indices (choix des mots, ironie, contrastes de la narration).

\courssub{Décoder les messages des médias camerounais}

Cette compétence sert aussi hors de l'école. Les messages publicitaires et les discours politiques diffusés à la radio et à la télévision camerounaises (CRTV, chaînes privées, radios communautaires) sont construits pour convaincre : leur contenu latent est souvent plus efficace que leur contenu manifeste.

\begin{exemplebox}[Analyse d'un slogan publicitaire]
Slogan télévisé : « \emph{Avec la bière X, sois un vrai lion, un vrai Indomptable !} »\par
\textbf{Contenu manifeste} : une invitation à acheter une boisson associée à l'équipe nationale de football.\par
\textbf{Contenu latent} : boire cette bière prouverait la force, le courage et la virilité (« vrai lion ») ; ne pas la boire, ce serait ne pas être un « vrai » supporter ni un « vrai » homme. Le slogan exploite l'amour des Lions Indomptables et la fierté nationale pour vendre un produit, en suggérant une identité que le produit seul ne peut évidemment pas donner.\par
\textbf{Décodage} : connotation des mots « lion », « vrai » ; présupposé (la consommation = preuve de valeur) ; non-dit (les dangers de l'alcool ne sont jamais mentionnés).
\end{exemplebox}

Le même travail s'applique aux discours politiques : quand un orateur déclare « nous avons \emph{enfin} apporté la paix et le développement dans cette région », le mot « enfin » présuppose que rien n'existait avant lui ; l'énumération flatteuse occulte (non-dit) les problèmes non résolus. Le citoyen averti, comme le lecteur averti, interroge le contexte : qui parle, à qui, dans quelle circonstance (campagne électorale, inauguration), et avec quelle intention ?

\begin{savaistu}[Le corps aussi parle]
Les contenus latents ne passent pas seulement par les mots : la \cle{communication non verbale} — gestes, regards, silences, posture, ton de la voix — en véhicule aussi. Un orateur qui détourne les yeux en promettant, un silence prolongé après une question gênante, un sourire forcé lors d'une poignée de main officielle : autant de signaux que le public décode spontanément. Dans \emph{Le Vieux nègre et la médaille}, le sourire de Meka et les gestes théâtraux du commandant disent autant que leurs paroles. Observer, c'est déjà lire l'implicite.
\end{savaistu}

\courssub{Méthode et importance pour le résumé}

\begin{methode}[Décoder l'implicite en trois étapes]
\begin{enumerate}
\item \textbf{Repérer les indices} : souligner les mots à connotation, les tournures ironiques, les adverbes comme « enfin », « encore », « même », les répétitions, les silences du texte (ce qui n'est pas dit alors qu'on l'attendrait).
\item \textbf{Interroger le contexte} : qui énonce ? à qui ? quand ? dans quelle situation ? quelle intention probable (informer, convaincre, dénoncer, séduire) ?
\item \textbf{Vérifier la cohérence} : confronter l'interprétation à l'ensemble du texte. Si le sens latent proposé contredit d'autres passages ou repose sur aucun indice, il faut l'abandonner. Une bonne lecture de l'implicite est \textbf{contrôlée par le texte}.
\end{enumerate}
\end{methode}

Pour la \cle{contraction de texte}, cette compétence est décisive mais délicate. Le résumé doit restituer le \textbf{sens profond} du texte : si l'auteur dénonce la colonisation par l'ironie, le résumé doit le dire clairement, car c'est l'idée essentielle. Mais il faut respecter deux limites : ne pas \textbf{inventer} (ajouter des idées absentes du texte) et ne pas \textbf{surinterpréter} (faire dire au texte plus qu'il ne dit). La règle d'or : dans le résumé, on explicite l'implicite \emph{de l'auteur}, jamais le sien.

\begin{exoresolu}[Repérer manifeste et latent dans un extrait]
\textbf{Énoncé.} Voici un extrait adapté du \emph{Vieux nègre et la médaille} : « Le commandant posa la médaille sur la poitrine du vieux Meka et déclara : "Voilà un homme qui a donné toute sa vie à la France. La France sait reconnaître ses enfants fidèles." » Indiquez le contenu manifeste et le contenu latent de ce passage, en justifiant par deux indices.
\tcblower
\textbf{Solution détaillée.} \emph{Contenu manifeste} : le commandant décore Meka et le félicite publiquement pour sa fidélité à la France ; le discours officiel est un éloge. \emph{Contenu latent} : le texte dénonce le paternalisme colonial et le caractère dérisoire de la récompense. \emph{Indice 1} : le mot « enfants » appliqué à un vieillard révèle le regard infantilisant du colonisateur (connotation paternaliste). \emph{Indice 2} : « donné toute sa vie » mis en regard d'une simple médaille crée un contraste disproportionné qui suggère l'iniquité de l'échange (implicite par contraste). L'interprétation reste fondée : elle s'appuie sur les mots du texte et sur le contexte colonial de l'œuvre.
\end{exoresolu}

\begin{exercice}[À vous de décoder]
1. Un candidat déclare à la radio : « Contrairement à d'autres, moi je n'ai jamais menti au peuple. » Relevez le présupposé et le non-dit de cette phrase.\par
2. Dans un extrait du roman, le narrateur écrit que Meka, médaille au cou, « ressemblait à un vieux perroquet savant ». Que suggère cette comparaison ? Quel procédé est utilisé ?\par
3. Pourquoi, dans un résumé, faut-il parfois expliciter le contenu latent ? Donnez un exemple tiré de l'œuvre.
\end{exercice}

\begin{corrige}[Exercice : décoder l'implicite]
1. \emph{Présupposé} : « contrairement à d'autres » présuppose que les autres candidats ont menti au peuple, sans que l'orateur le prouve. \emph{Non-dit} : le candidat ne cite aucun fait, aucun programme ; il tait ses propres erreurs éventuelles et ne définit pas « le peuple ».\par
2. La comparaison avec « un vieux perroquet savant » suggère que Meka, paré de la médaille, est devenu un objet de curiosité, un animal dressé qu'on exhibe : la récompense le ridiculise au lieu de l'honorer. Le procédé est l'\textbf{ironie} (appuyée par la comparaison animalière dévalorisante).\par
3. Il faut expliciter le contenu latent quand il porte l'idée essentielle du texte : Oyono ne dit jamais directement « la colonisation est injuste », il le fait comprendre par l'ironie et les contrastes. Un résumé qui ne retiendrait que le manifeste (« Meka reçoit une médaille et fait la fête ») trahirait l'œuvre. On écrira donc, par exemple : « À travers la récompense dérisoire de Meka, l'auteur dénonce l'hypocrisie et l'injustice de la domination coloniale », ce qui restitue le sens profond sans rien inventer.
\end{corrige}

\begin{aretenir}[Pour le Probatoire et le Baccalauréat technique]
Dans les questions de compréhension, les formulations « Que suggère l'auteur ? », « Quel est le sens implicite de ce passage ? », « Relevez l'ironie » visent précisément le contenu latent. Répondez toujours en deux temps : (1) formuler le sens profond, (2) citer l'indice du texte qui le justifie. Cette double exigence — interpréter \emph{et} prouver — distingue la bonne copie de la copie moyenne.
\end{aretenir}

\coursec{Reformulation, exposé et activités d'intégration}

Dans la section précédente, nous avons appris à distinguer les \cle{contenus manifestes} (ce que le texte dit explicitement) des \cle{contenus latents} (ce qu'il suggère sans le dire). Cette compétence de lecture fine trouve maintenant son aboutissement pratique : une fois les idées essentielles repérées et hiérarchisées, il faut les \emph{dire autrement}, avec ses propres mots, sans trahir la pensée de l'auteur. C'est le cœur de la \cle{reformulation}, dernière grande étape de la contraction de texte. Nous verrons ensuite comment cette même compétence — hiérarchiser puis reformuler — se transpose à l'oral dans l'\cle{exposé}, avant de nous entraîner collectivement lors d'une activité d'intégration.

\courssub{La reformulation : dire autrement les idées essentielles}

\textbf{Intuition.} Imaginez qu'un camarade absent vous demande de lui raconter le cours d'histoire de la matinée. Vous ne récitez pas mot pour mot les phrases du professeur : vous retenez les idées principales et vous les exprimez avec votre propre vocabulaire, vos propres tournures. C'est exactement ce qu'exige la reformulation dans le résumé : restituer la \emph{substance} du texte, jamais sa \emph{forme}.

\begin{definition}[Reformulation]
La \cle{reformulation} est l'opération qui consiste à exprimer avec ses propres mots les idées essentielles d'un texte, en respectant scrupuleusement la pensée de l'auteur, sans rien ajouter (ni avis personnel, ni exemple nouveau) et sans rien omettre d'essentiel. Elle intervient après le repérage et la réduction des idées, et avant la vérification finale du résumé.
\end{definition}

La reformulation répond à une exigence double. D'une part, le correcteur du Probatoire et du Baccalauréat technique sanctionne le \emph{plagiat} : recopier des phrases entières du texte, même bien choisies, prouve que le candidat n'a pas assimilé le contenu. D'autre part, reformuler est la meilleure preuve de compréhension : on ne peut dire autrement que ce que l'on a réellement compris.

\begin{attention}[L'avis personnel est interdit]
L'erreur la plus grave et la plus fréquente en contraction de texte est d'introduire son \textbf{opinion personnelle} dans le résumé : « je pense que l'auteur a raison », « cette idée est fausse car... », « heureusement que... ». Le résumé n'est ni un commentaire ni une dissertation : c'est la restitution \emph{neutre et fidèle} de la pensée d'autrui. Toute marque de jugement (« malheureusement », « à juste titre », « prétendument ») qui n'est pas dans le texte constitue une faute méthodologique éliminatoire. De même, il est interdit d'ajouter des exemples ou des informations extérieures au texte.
\end{attention}

\courssub{Les techniques de reformulation}

Pour reformuler efficacement, trois techniques principales sont à votre disposition. Elles se combinent le plus souvent.

\textbf{1. La substitution par synonymes.} On remplace les mots du texte par des équivalents : « l'école » devient « l'institution scolaire » ; « riche » devient « fortuné » ; « lutter » devient « combattre ». Attention : le synonyme doit convenir au contexte (« riche en idées » ne devient pas « fortuné en idées » !).

\textbf{2. Le changement de tournure.} On transforme la construction de la phrase :
\begin{itemize}
\item passage de la voix active à la voix passive (ou inversement) : « Le colonisateur méprisa les traditions » $\rightarrow$ « Les traditions furent méprisées par le colonisateur » ;
\item transformation d'une subordonnée en proposition indépendante, ou inversement : « Parce qu'il était courageux, Meka résista » $\rightarrow$ « Meka, courageux, résista » ;
\item fusion de deux phrases courtes en une seule, ou découpage d'une phrase longue.
\end{itemize}

\textbf{3. La nominalisation.} C'est la technique la plus puissante pour condenser.

\begin{definition}[Nominalisation]
La \cle{nominalisation} consiste à transformer un verbe ou un adjectif en nom (substantif), ce qui permet de condenser une proposition entière en un simple groupe nominal. Exemple : « Les paysans \emph{cultivent} le cacao et cela les \emph{enrichit} » $\rightarrow$ « La \emph{culture} du cacao assure l'\emph{enrichissement} des paysans » (deux propositions réduites à une).
\end{definition}

\begin{exemplebox}[Les trois techniques appliquées à une même phrase]
Phrase originale : « Le gouvernement a décidé de construire des écoles parce que l'éducation développe le pays. »
\begin{itemize}
\item Synonymes : « Les autorités ont choisi d'édifier des établissements scolaires car l'instruction fait progresser la nation. »
\item Changement de tournure : « Des écoles seront construites par le gouvernement, l'éducation étant un facteur de développement. »
\item Nominalisation : « La construction d'écoles par le gouvernement vise le développement du pays par l'éducation. »
\end{itemize}
Dans les trois cas, le sens est strictement identique ; seule la forme change.
\end{exemplebox}

\begin{attention}[Pièges de la reformulation]
Trois pièges guettent le candidat : (1) le \emph{faux synonyme} qui déforme le sens (« diminuer » n'est pas « supprimer ») ; (2) la reformulation qui \emph{généralise} abusivement (« certains jeunes » devenant « les jeunes ») ; (3) la reformulation qui \emph{affaiblit} ou \emph{renforce} la nuance de l'auteur (« il semble que » devenant « il est certain que »). Après chaque phrase reformulée, posez-vous la question : « L'auteur signerait-il cette phrase ? »
\end{attention}

\courssub{La vérification finale du résumé}

La reformulation achevée, le travail n'est pas terminé : il reste à relire et à contrôler le texte produit. Cette étape, trop souvent négligée par manque de temps, fait pourtant la différence entre une copie moyenne et une excellente copie.

\begin{methode}[Vérifier son résumé en quatre questions]
Avant de rendre votre résumé, interrogez votre texte :
\begin{enumerate}
\item \textbf{Est-il fidèle ?} Toutes les idées essentielles du texte sont-elles présentes, sans déformation, sans ajout, sans avis personnel ?
\item \textbf{Est-il cohérent ?} Les phrases s'enchaînent-elles logiquement grâce aux connecteurs (d'abord, ensuite, en effet, cependant, ainsi, enfin) ? Le texte se lit-il comme un tout, et non comme une liste de fragments ?
\item \textbf{Est-il autonome ?} Un lecteur qui n'a pas lu le texte original peut-il comprendre votre résumé seul ? (Pas de « comme le dit l'auteur dans le deuxième paragraphe », pas de renvoi au texte, pas de pronoms sans référent.)
\item \textbf{Le nombre de mots est-il respecté ?} Comptez vos mots et indiquez le total à la fin. La consigne (par exemple 100 mots, soit le quart d'un texte de 400 mots) admet une marge de tolérance d'environ 10 \% : entre 90 et 110 mots pour une consigne de 100.
\end{enumerate}
Ajoutez enfin un contrôle de la correction grammaticale : accords, conjugaisons, orthographe, ponctuation.
\end{methode}

L'organigramme suivant synthétise le parcours complet de la contraction de texte, de la première lecture à la vérification finale.

\begin{popfigure}
\begin{center}
\begin{tikzpicture}[
  node distance=0.55cm,
  etape/.style={rectangle, rounded corners=3pt, draw=popblue, fill=popblueL, text width=4.6cm, align=center, minimum height=1cm, font=\small},
  detail/.style={rectangle, draw=popmuted, fill=popcream, text width=5.6cm, align=left, font=\footnotesize, right=0.5cm of #1},
  fleche/.style={-{Stealth[length=3mm]}, thick, popdark}
]
\node[etape, align=center] (e1){\textbf{1. Lecture}\\ Lire deux fois le texte, saisir le sens global et la thèse};
\node[etape, below=of e1, align=center] (e2){\textbf{2. Repérage}\\ Souligner les idées essentielles, dégager le plan};
\node[etape, below=of e2, align=center] (e3){\textbf{3. Réduction}\\ Supprimer exemples, répétitions, détails ; regrouper};
\node[etape, below=of e3, align=center] (e4){\textbf{4. Reformulation}\\ Dire autrement : synonymes, tournures, nominalisation};
\node[etape, below=of e4, fill=popgreenL, draw=popgreen, align=center] (e5){\textbf{5. Vérification}\\ Fidèle ? Cohérent ? Autonome ? Nombre de mots ?};
\draw[fleche] (e1) -- (e2);
\draw[fleche] (e2) -- (e3);
\draw[fleche] (e3) -- (e4);
\draw[fleche] (e4) -- (e5);
\draw[fleche, dashed, poporange] (e5.west) -- ++(-1.2,0) |- (e4.west)
  node[midway, left, font=\footnotesize, text=poporange, align=center]{si échec,\\ reprendre};
\end{tikzpicture}
\end{center}
\legende{Organigramme des cinq étapes de la contraction de texte. La vérification (étape 5) peut renvoyer à la reformulation (étape 4) si l'un des quatre critères n'est pas satisfait.}
\end{popfigure}

\courssub{Exemple résolu : un résumé complet, corrigé et commenté}

\begin{exoresolu}[Réduire un texte de 400 mots à 100 mots]
\textbf{Énoncé.} Voici un texte argumentatif de 400 mots (reconstitué pour l'exercice). Rédigez-en le résumé en 100 mots environ.

\smallskip
\emph{« Dans nombre de nos villages, l'école est encore perçue comme une étrangère. Les parents, il faut le reconnaître, envoient leurs enfants en classe sans toujours comprendre pourquoi. Pourtant, l'éducation demeure le seul ascenseur social véritablement accessible aux familles modestes. Prenons l'exemple de ce jeune de Bangangté, fils de cultivateurs, devenu ingénieur grâce à sa bourse d'excellence : son histoire, admirable, illustre ce que l'école peut offrir. Certes, l'école coûte cher : fournitures, uniformes, contributions diverses pèsent lourd sur les budgets familiaux. Certains parents préfèrent d'ailleurs garder les enfants aux champs, où leurs bras sont immédiatement utiles. Mais ce calcul à courte vue appauvrit durablement les familles. Un enfant instruit, c'est demain un salaire, une exploitation agricole modernisée, des soins de santé mieux compris. L'école, en outre, forme des citoyens capables de lire un contrat, de défendre leurs droits, de participer à la vie publique. Investir dans l'éducation, c'est donc investir dans l'avenir de toute la communauté. »} (Texte reconstitué à des fins pédagogiques.)

\smallskip
\tcblower
\textbf{Étape 1 — Repérage des idées essentielles} (la thèse est : l'école est un investissement indispensable) :
\begin{itemize}
\item Idée 1 : l'école est parfois mal comprise des parents dans les villages.
\item Idée 2 : elle est pourtant le seul moyen de promotion sociale des familles modestes.
\item Idée 3 : elle coûte cher, ce qui décourage certains parents (concession).
\item Idée 4 : mais garder les enfants aux champs appauvrit les familles à long terme.
\item Idée 5 : l'instruction rapporte : salaire, agriculture modernisée, santé.
\item Idée 6 : elle forme aussi des citoyens conscients de leurs droits.
\item Conclusion : investir dans l'éducation, c'est investir dans l'avenir collectif.
\end{itemize}
\textbf{Étape 2 — Suppressions} : l'exemple du jeune de Bangangté (illustration), les détails (fournitures, uniformes), les répétitions et les formules d'insistance (« il faut le reconnaître », « admirable »).

\smallskip
\textbf{Étape 3 — Résumé reformulé (102 mots)} :

\smallskip
« Dans les villages, l'école reste mal comprise de certains parents. Pourtant, elle constitue le seul moyen d'ascension sociale offert aux familles pauvres. Si son coût pousse parfois les parents à garder leurs enfants aux champs, ce choix les appauvrit durablement : l'instruction garantit en effet un revenu futur, une agriculture modernisée et une meilleure santé. Elle forme en outre des citoyens capables de défendre leurs droits. Éduquer les enfants, c'est donc investir dans l'avenir de toute la communauté. » (102 mots)

\smallskip
\textbf{Commentaire du corrigé.} Le résumé respecte les quatre critères : il est \emph{fidèle} (les six idées essentielles y figurent, la concession « si son coût... » est conservée) ; il est \emph{cohérent} (connecteurs : « pourtant », « si », « en effet », « en outre », « donc ») ; il est \emph{autonome} (aucune référence au texte) ; il respecte le \emph{nombre de mots} (102 pour 100 exigés). La nominalisation transforme « devenir ingénieur grâce à sa bourse » en « moyen d'ascension sociale ». Aucun avis personnel n'a été ajouté.
\end{exoresolu}

\courssub{L'exposé : de la préparation à la présentation}

Le résumé est un exercice écrit ; l'\cle{exposé} en est le frère oral. Dans les deux cas, il s'agit de \textbf{hiérarchiser} une information (distinguer l'essentiel de l'accessoire) puis de la \textbf{reformuler} pour un destinataire. Celui qui sait résumer sait préparer un exposé ; celui qui sait présenter un exposé a déjà l'aisance de celui qui reformule.

\textbf{1. La préparation de l'exposé.}
\begin{itemize}
\item \textbf{Définir le sujet et l'objectif} : que doit savoir l'auditoire à la fin ?
\item \textbf{Rechercher et sélectionner les informations} : comme dans le résumé, on élimine le superflu.
\item \textbf{Construire un plan} en trois parties : introduction (accroche, sujet, annonce du plan), développement (deux ou trois parties ordonnées), conclusion (bilan, ouverture).
\item \textbf{Rédiger des fiches} : jamais de texte lu intégralement ! On note sur de petites fiches les idées forces, les chiffres, les transitions — autant de reformulations condensées.
\item \textbf{Préparer les supports} : affiche, tableau, schéma, diaporama sobre. Un bon support illustre, il ne remplace pas la parole.
\item \textbf{Répéter} à voix haute, en chronométrant.
\end{itemize}

\textbf{2. La présentation orale.}
\begin{itemize}
\item \textbf{La voix} : débit ni trop rapide ni trop lent, volume audible au fond de la salle, articulation soignée, intonations variées pour éviter la monotonie.
\item \textbf{Le regard} : balayer l'auditoire, ne pas fixer ses fiches ni le plafond ; le regard capte l'attention et mesure la compréhension de l'auditoire.
\item \textbf{La gestuelle et la posture} : se tenir droit, gestes naturels qui soulignent les idées, éviter les tics (mains dans les poches, balancements).
\item \textbf{La gestion du temps} : respecter la durée impartie, annoncer les transitions (« j'en viens maintenant à... »), conclure sans s'éterniser.
\end{itemize}

\begin{center}
\begin{tabularx}{\linewidth}{|l|X|c|}
\hline
\rowcolor{popblueL} \textbf{Critère} & \textbf{Indicateurs observés} & \textbf{Points} \\
\hline
Contenu & Exactitude des informations ; hiérarchisation (essentiel/accessoire) ; plan clair et annoncé ; reformulation personnelle (pas de lecture) & /6 \\
\hline
Expression orale & Voix audible et posée ; articulation ; vocabulaire correct ; fluidité ; transitions explicites & /6 \\
\hline
Posture & Regard vers l'auditoire ; gestuelle naturelle ; tenue ; maîtrise du stress & /4 \\
\hline
Temps et supports & Durée respectée ; supports clairs et utilisés à bon escient ; réponses aux questions & /4 \\
\hline
\rowcolor{popgreenL} \textbf{Total} & & \textbf{/20} \\
\hline
\end{tabularx}
\end{center}

\begin{savaistu}[L'exposé, une compétence professionnelle]
Savoir exposer clairement une idée est l'une des compétences les plus recherchées dans la vie professionnelle technique. Le technicien qui présente un diagnostic de panne à son chef d'atelier, l'infirmier qui transmet une consigne au changement d'équipe, l'informaticien qui défend son projet devant un client, l'agent commercial qui présente un produit : tous pratiquent, sans le nommer, la même compétence que le résumé et l'exposé scolaires — sélectionner l'essentiel, l'ordonner, le formuler clairement dans un temps limité. Les enquêtes auprès des employeurs placent régulièrement la « communication orale » parmi les trois qualités les plus attendues des jeunes diplômés techniques.
\end{savaistu}

\courssub{Activité d'intégration : reformuler en groupe}

\begin{experience}[Activité d'intégration en classe (2 séances)]
\textbf{Objectif} : reformuler collectivement les idées essentielles d'un texte à contracter, puis présenter le résultat à l'oral.

\smallskip
\textbf{Matériel} : un texte argumentatif de 400 mots (extrait de presse ou texte fourni par le professeur), une copie par groupe, des fiches cartonnées, le tableau.

\smallskip
\textbf{Protocole.}
\begin{enumerate}
\item La classe est divisée en groupes de quatre à six élèves. Chaque groupe lit le texte deux fois, en silence.
\item Chaque groupe souligne les idées essentielles, puis les rédige sur ses fiches, une idée par fiche, \emph{reformulée} (interdiction de recopier une phrase du texte).
\item Les fiches sont ordonnées : le groupe construit le plan du résumé et choisit les connecteurs logiques.
\item Le groupe rédige le résumé complet en 100 mots, le compte et le vérifie avec les quatre questions de la méthode.
\item Chaque groupe désigne un rapporteur qui \emph{présente} le résumé à la classe en deux minutes, selon les critères de la grille d'évaluation.
\end{enumerate}
\textbf{Observations attendues} : les résumés des différents groupes, bien que formulés différemment, doivent contenir les mêmes idées essentielles — preuve que le sens d'un texte est objectif alors que sa formulation est libre.
\end{experience}

\courssub{Compte-rendu collectif, correction et remédiation}

L'activité se clôt par un \textbf{compte-rendu collectif} au tableau. Le professeur confronte les productions des groupes et fait émerger, avec la classe, la liste des idées essentielles attendues. Chaque écart est discuté : une idée essentielle oubliée ? un détail conservé à tort ? une phrase recopiée ? un avis personnel glissé dans le résumé ?

La \textbf{remédiation} porte sur les erreurs fréquentes recensées pendant les séances de l'unité :

\begin{center}
\begin{tabularx}{\linewidth}{|l|X|X|}
\hline
\rowcolor{popblueL} \textbf{Erreur fréquente} & \textbf{Pourquoi c'est une faute} & \textbf{Remédiation} \\
\hline
Recopier des phrases du texte & Plagiat : ne prouve pas la compréhension & Imposer la règle « texte fermé pendant la rédaction » \\
\hline
Donner son avis & Le résumé restitue la pensée d'autrui & Repérer et biffer toute marque de jugement \\
\hline
Conserver exemples et détails & Le résumé ne garde que l'essentiel & S'interroger : « cette phrase est-elle indispensable à la thèse ? » \\
\hline
Enchaîner sans connecteurs & Le texte devient une liste incohérente & Réutiliser les conjonctions et connecteurs étudiés dans l'unité \\
\hline
Dépasser le nombre de mots & Non-respect de la consigne & Compter, puis couper les redondances et nominaliser \\
\hline
Lire ses fiches à l'oral & L'exposé exige le contact avec l'auditoire & Fiches réduites à des mots-clés ; répétition chronométrée \\
\hline
\end{tabularx}
\end{center}

\begin{aretenir}[L'essentiel de la section]
\begin{itemize}
\item \textbf{Reformuler}, c'est dire autrement les idées essentielles avec son propre vocabulaire, sans trahir ni compléter la pensée de l'auteur ; l'avis personnel est \textbf{interdit}.
\item Trois techniques : les \textbf{synonymes}, le \textbf{changement de tournure} (actif/passif, subordonnée/indépendante) et la \textbf{nominalisation} (verbe ou adjectif transformé en nom).
\item La vérification finale tient en quatre questions : le résumé est-il \textbf{fidèle}, \textbf{cohérent}, \textbf{autonome}, au bon \textbf{nombre de mots} ?
\item L'\textbf{exposé} transpose à l'oral les mêmes compétences : hiérarchiser, reformuler ; il se prépare (plan, fiches, supports) et se présente (voix, regard, gestuelle, temps).
\item Résumé écrit et exposé oral sont deux faces d'une même compétence, précieuse à l'examen comme dans la vie professionnelle technique.
\end{itemize}
\end{aretenir}

\begin{exercice}[Pour s'entraîner]
1. Reformulez la phrase suivante par nominalisation : « Les élèves travaillent régulièrement, ce qui leur permet de réussir. »

2. Un candidat écrit dans son résumé : « L'auteur dénonce le mépris du commandant, et il a bien raison car le colonialisme fut une injustice. » Relevez les deux fautes méthodologiques.

3. Voici un texte de 120 mots. Rédigez-en le résumé en 30 mots environ, puis vérifiez-le avec les quatre questions de la méthode : « Le sport scolaire est souvent sacrifié au profit des matières dites sérieuses. C'est une erreur. D'abord, l'exercice physique améliore la concentration et donc les résultats scolaires. Ensuite, il enseigne le respect des règles et l'esprit d'équipe, qualités utiles dans la vie professionnelle. Enfin, il préserve la santé des adolescents, de plus en plus sédentaires. Supprimer le sport à l'école, c'est donc appauvrir la formation de l'élève tout entier. »
\end{exercice}

\begin{corrige}[Exercice]
1. « Le travail régulier des élèves permet leur réussite » (ou : « grâce à leur travail régulier, les élèves réussissent »). Les verbes \emph{travailler} et \emph{réussir} deviennent les noms \emph{travail} et \emph{réussite}.

2. Première faute : le jugement de valeur « il a bien raison » exprime l'avis personnel du candidat, interdit dans un résumé. Deuxième faute : « car le colonialisme fut une injustice » ajoute une information extérieure au texte ; le résumé ne doit contenir que la pensée de l'auteur.

3. Résumé possible (31 mots) : « Souvent négligé, le sport scolaire est pourtant précieux : il développe concentration, esprit d'équipe et santé. Le supprimer reviendrait à appauvrir la formation complète de l'élève. » Vérification : \emph{fidèle} (les trois arguments et la conclusion sont présents, la concession initiale aussi) ; \emph{cohérent} (« pourtant », deux-points énumératifs) ; \emph{autonome} (compréhensible sans le texte) ; \emph{nombre de mots} respecté (31 pour 30 exigés, dans la marge de 10 \%).
\end{corrige}

\newpage

\coursec{Activité d'intégration}

\coursec{Leçon 21 : Formulation des idées essentielles d'un texte à résumer}

\courssub{Objectifs de l'activité}

À l'issue de cette activité d'intégration, l'élève doit être capable de :

\begin{itemize}
\item repérer dans un texte argumentatif de presse la \cle{thèse}, les \cle{arguments} et les \cle{exemples} ;
\item distinguer le \cle{contenu manifeste} (ce qui est dit explicitement) du \cle{contenu latent} (ce qui est suggéré, sous-entendu) ;
\item dégager les \cle{idées essentielles} et appliquer les techniques de réduction : \cle{suppression}, \cle{substitution} et \cle{fusion} ;
\item rédiger une \cle{contraction de texte} fidèle, cohérente et personnelle, en respectant le nombre de mots imposé ;
\item employer correctement les \cle{conjonctions de coordination} et de \cle{subordination} pour articuler les idées ;
\item préparer et présenter un \cle{exposé oral} bref, clair et structuré ;
\item évaluer la production d'autrui à l'aide d'une grille et participer à la remédiation collective.
\end{itemize}

\begin{aretenir}[Compétence visée]
L'activité mobilise en une seule tâche l'ensemble des savoir-faire de la leçon : \emph{lire} méthodiquement, \emph{analyser} la structure argumentative, \emph{réduire}, \emph{reformuler} et \emph{présenter} à l'oral. C'est exactement la démarche exigée à l'épreuve écrite du Probatoire et du Baccalauréat technique pour le résumé de texte.
\end{aretenir}

\courssub{Contexte programmatique : lien avec les lectures méthodiques}

\begin{savaistu}[Le Vieux nègre et la médaille -- Lectures n°4 et n°5]
Les lectures méthodiques n°4 et n°5 de l'œuvre \emph{Le Vieux nègre et la médaille} (Sembène Ousmane) ont permis d'étudier la structure narrative et argumentative d'un texte littéraire (thèse implicite sur la colonisation, arguments à travers le parcours de M. Ndiaye, exemples concrets). Cette activité d'intégration transpose ces compétences sur un texte journalistique contemporain : l'analyse de la thèse, des arguments et des exemples, la distinction manifeste/latent, et la reformulation synthétique. Le passage du littéraire au journalistique ancre la compétence dans la réalité professionnelle du technicien (rapports, notes de synthèse, veille documentaire).
\end{savaistu}

\courssub{Situation}

Le club de presse du Lycée technique de Nkolbisson prépare son bulletin trimestriel intitulé \emph{L'Étudiant technique}. Le rédacteur en chef, élève de Terminale F4, a sélectionné un éditorial paru dans un quotidien camerounais et consacré à l'entrepreneuriat des jeunes diplômés techniques. Le bulletin ne disposant que d'une demi-page, l'article original de \textbf{500 mots} doit être ramené à \textbf{125 mots} (soit le quart du texte initial, avec une marge tolérée de 10 \% : entre 113 et 137 mots). Le résumé retenu sera présenté oralement lors de la séance du club, en \textbf{3 minutes} exactement.

Voici l'éditorial à contracter (extraits représentatifs, reconstitués à des fins pédagogiques ; le texte ci-dessous compte environ 160 mots pour l'exercice en classe, il représente la structure de l'original de 500 mots) :

\begin{exemplebox}[Éditorial à résumer : « Oser entreprendre après le Baccalauréat technique »]
\begin{itshape}Chaque année, des milliers de jeunes quittent nos lycées techniques, le précieux diplôme en poche, et frappent aux portes des administrations et des grandes entreprises. Beaucoup attendent des mois, parfois des années, un emploi salarié qui ne vient pas. Pourtant, une autre voie existe : l'entrepreneuriat.

Le jeune diplômé technique possède un trésor que beaucoup ignorent : un savoir-faire concret. Le titulaire d'un BAC F4 sait construire ; celui d'un BAC F3 répare et installe ; celui d'un BAC G sait gérer et comptabiliser. Ces compétences, le marché camerounais les réclame chaque jour. À Yaoundé comme à Bafoussam, les ateliers de menuiserie, de froid et climatisation ou de maintenance informatique créés par d'anciens élèves des lycées techniques prospèrent.

Certes, les obstacles sont réels : le manque de financement, la peur de l'échec, le mépris social dont souffre encore le travail manuel. Mais ces obstacles ne sont pas insurmontables. Les microfinances, les incubateurs et les programmes gouvernementaux comme le PIAASI offrent des appuis. Surtout, le regroupement en GIC permet aux jeunes de mutualiser leurs moyens.

Attendre indéfiniment un emploi public, c'est laisser dormir son diplôme. Entreprendre, c'est le faire fructifier. L'État ne pourra jamais embaucher tous les diplômés ; le secteur privé formel non plus. La solution passe donc par la création d'activités. Nos jeunes techniciens doivent cesser d'être des demandeurs d'emploi pour devenir des créateurs d'emplois. C'est à ce prix que le Cameroun transformera ses diplômés en bâtisseurs de son émergence.\end{itshape}
\end{exemplebox}

\courssub{Consignes}

Travail en groupes de quatre élèves. Chaque groupe traite les tâches suivantes, dans l'ordre, et consigne ses réponses sur une feuille de travail.

\begin{enumerate}
\item \textbf{Lecture méthodique.} Lisez l'éditorial deux fois. À la deuxième lecture, annotez le texte : soulignez la thèse en rouge, encadrez les arguments, numérotez les exemples.
\item \textbf{Structure argumentative.} Répondez par écrit : (a) quelle est la thèse défendue par l'éditorialiste ? (b) citez trois arguments qui la soutiennent ; (c) citez deux exemples concrets qui illustrent ces arguments.
\item \textbf{Manifeste et latent.} (a) Relevez deux informations explicitement affirmées (contenu manifeste). (b) Dégagez deux idées sous-entendues (contenu latent) : que suggère l'auteur sur l'emploi public ? sur le regard social porté sur les métiers techniques ?
\item \textbf{Idées essentielles.} Listez les cinq idées essentielles du texte, c'est-à-dire celles dont la suppression rendrait le résumé infidèle.
\item \textbf{Réduction.} Appliquez les trois techniques : indiquez (a) deux passages à supprimer (répétitions, détails, exemples secondaires) ; (b) deux segments à substituer par un mot ou une expression générique ; (c) deux phrases à fusionner.
\item \textbf{Contraction.} Rédigez le résumé en \textbf{125 mots} ($\pm$ 10 \%), en respectant les contraintes suivantes : aucune phrase recopiée du texte ; emploi d'au moins \textbf{trois conjonctions de subordination} (par exemple \emph{parce que, puisque, bien que, si, lorsque}) et \textbf{deux conjonctions de coordination} (\emph{mais, ou, et, donc, or, ni, car}) ; aucun jugement personnel.
\item \textbf{Exposé.} Désignez un rapporteur. Préparez sa prise de parole : introduction (présentation de l'article et de son auteur), annonce du plan, lecture posée du résumé, conclusion (intérêt du sujet pour les élèves). Durée : 3 minutes.
\item \textbf{Évaluation croisée.} Pendant chaque exposé, les autres groupes remplissent la grille d'évaluation ci-dessous et attribuent une note sur 20.
\item \textbf{Compte-rendu et remédiation.} À l'issue des exposés, la classe dresse la liste des erreurs les plus fréquentes et propose collectivement des corrections.
\end{enumerate}

\courssub{Grille d'évaluation}

\begin{center}
\begin{tabularx}{\linewidth}{|l|X|c|}
\hline
\rowcolor{popblueL} \textbf{Critère} & \textbf{Indicateurs} & \textbf{Barème} \\
\hline
Fidélité & Thèse et idées essentielles respectées, aucune idée ajoutée, aucun jugement personnel & /6 \\
\hline
Cohérence & Enchaînement logique des idées, connecteurs correctement employés & /4 \\
\hline
Langue & Orthographe, grammaire, conjugaison, ponctuation & /5 \\
\hline
Contraintes & Nombre de mots respecté (113--137), aucune phrase recopiée & /3 \\
\hline
Présentation orale & Voix, posture, respect des 3 minutes, plan annoncé & /2 \\
\hline
\end{tabularx}
\end{center}

\courssub{Ressources à mobiliser}

\begin{aretenir}[Structure du texte argumentatif]
Un texte argumentatif se compose d'une \cle{thèse} (idée défendue), d'\cle{arguments} (raisons qui la justifient) et d'\cle{exemples} (faits concrets qui illustrent les arguments). Le résumé conserve la thèse et les arguments ; il élimine la plupart des exemples, des répétitions et des figures de style.
\end{aretenir}

\begin{aretenir}[Les trois techniques de réduction]
\begin{itemize}
\item \textbf{Suppression} : éliminer répétitions, détails, exemples, énumérations, citations, figures de style.
\item \textbf{Substitution} : remplacer un groupe de mots par un mot générique (\emph{les ateliers de menuiserie, de froid et de maintenance} $\rightarrow$ \emph{les ateliers techniques}) ; remplacer une proposition par un nom (\emph{parce qu'il a créé son atelier} $\rightarrow$ \emph{grâce à la création de son atelier}).
\item \textbf{Fusion} : réunir en une seule phrase deux phrases exprimant des idées proches, à l'aide d'un connecteur.
\end{itemize}
\end{aretenir}

\begin{aretenir}[Conjonctions utiles]
\begin{itemize}
\item \textbf{Coordination} : \emph{mais, ou, et, donc, or, ni, car}. Elles relient des éléments de même nature.
\item \textbf{Subordination} : \emph{que, parce que, puisque, bien que, quoique, si, lorsque, comme, afin que, pour que}. Elles introduisent une proposition dépendante et expriment la cause, la concession, la condition, le temps ou le but.
\end{itemize}
\end{aretenir}

\begin{attention}[Pièges fréquents]
Ne jamais recopier de phrases du texte (c'est de la compilation, sanctionnée) ; ne jamais donner son avis (\emph{je pense que}, \emph{à mon avis} sont interdits dans un résumé) ; compter les mots avec soin ; ne pas confondre un exemple (supprimable) et un argument (à conserver) ; employer \emph{car} (coordination) et non \emph{parce que} (subordination) pour relier deux propositions indépendantes.
\end{attention}

\courssub{Démarche et corrigé commenté}

\begin{exoresolu}[Contraction de l'éditorial : corrigé type]
\textbf{Énoncé.} Résumez l'éditorial « Oser entreprendre après le Baccalauréat technique » (500 mots) en 125 mots ($\pm$ 10 \%), sans phrase recopiée ni jugement personnel, en employant au moins trois conjonctions de subordination et deux de coordination.

\tcblower

\textbf{Étape 1 : repérage de la structure argumentative.}

\begin{itemize}
\item \textbf{Thèse} : les jeunes diplômés techniques doivent créer leur propre activité plutôt que d'attendre un emploi salarié.
\item \textbf{Argument 1} : ils possèdent un savoir-faire concret et recherché par le marché.
\item \textbf{Argument 2} : l'emploi public et privé formel ne peut pas absorber tous les diplômés.
\item \textbf{Argument 3} : les obstacles (financement, peur de l'échec, mépris du travail manuel) peuvent être surmontés grâce aux microfinances, aux incubateurs et aux GIC.
\item \textbf{Exemples} (à supprimer ou substituer) : les ateliers de Yaoundé et Bafoussam, les spécialités F4, F3, G, le PIAASI.
\end{itemize}

\textbf{Étape 2 : contenus manifeste et latent.}

Manifeste : l'attente de l'emploi salarié est longue ; des appuis financiers existent. Latent : l'auteur suggère que l'État est défaillant dans l'insertion des jeunes et que la société dévalorise injustement les métiers manuels ; il invite implicitement à changer de mentalité. Le résumé ne retient que le manifeste.

\textbf{Étape 3 : idées essentielles retenues.}

(1) Beaucoup de diplômés techniques attendent vainement un emploi salarié. (2) Leur savoir-faire est pourtant demandé. (3) L'État et le secteur formel ne peuvent pas tous les embaucher. (4) Des appuis existent pour entreprendre. (5) Ils doivent devenir créateurs d'emplois.

\textbf{Étape 4 : application des techniques.}

\begin{itemize}
\item \emph{Suppression} : l'énumération des villes, des spécialités et du nom du programme ; l'antithèse finale (« demandeurs / créateurs ») est reformulée brièvement.
\item \emph{Substitution} : « les microfinances, les incubateurs et les programmes gouvernementaux » $\rightarrow$ « des dispositifs d'appui — microfinances, incubateurs — » ; « frappent aux portes des administrations » $\rightarrow$ « cherchent un emploi salarié ».
\item \emph{Fusion} : les phrases sur les obstacles et sur les solutions sont réunies par \emph{mais} (ou subordination concessive \emph{bien que}).
\end{itemize}

\textbf{Étape 5 : résumé rédigé (122 mots).}

\emph{Chaque année, de nombreux jeunes diplômés des lycées techniques camerounais attendent vainement un emploi salarié, parce que l'État et le secteur privé formel ne peuvent absorber tous les sortants. Pourtant, leur savoir-faire technique concret est très recherché sur le marché, et d'anciens élèves réussissent déjà en créant leurs propres ateliers de production ou de service. Bien que le manque de financement, la peur de l'échec et le mépris du travail manuel constituent des obstacles réels, des dispositifs d'appui — microfinances, incubateurs — et le regroupement en GIC permettent de les surmonter. L'éditorialiste invite donc ces jeunes à entreprendre plutôt qu'à attendre : lorsqu'ils créent leur activité, ils font fructifier leur diplôme et deviennent créateurs d'emplois, participant ainsi activement à l'émergence du pays.}

Vérification des contraintes : 122 mots (dans la fourchette 113--137) ; trois subordonnantes (\emph{parce que}, \emph{bien que}, \emph{lorsqu'ils}) ; trois coordinations (\emph{et}, \emph{ou}, \emph{donc}) ; aucune phrase recopiée ; aucun avis personnel.

\textbf{Étape 6 : préparation de l'exposé (3 minutes).}

\begin{enumerate}
\item Introduction (30 s) : présenter l'article, son support, son thème.
\item Annonce du plan (15 s) : « Nous présenterons la thèse, puis les arguments, puis notre résumé. »
\item Lecture du résumé (1 min 30), voix posée, regard vers l'auditoire.
\item Conclusion (30 s) : intérêt du sujet pour les élèves de Terminale technique.
\end{enumerate}

\textbf{Étape 7 : remédiation.} Erreurs relevées en séance : phrases recopiées (reformuler avec ses propres mots), oubli du comptage (compter par lignes de dix), mélange des registres (ne pas écrire \emph{je trouve que}), confusion entre \emph{car} (coordination) et \emph{parce que} (subordination), omission d'idées essentielles au profit de détails.

\end{exoresolu}

\courssub{Bilan de l'activité}

Cette activité a permis de mettre en évidence que la contraction de texte n'est pas un simple raccourcissement mécanique, mais une chaîne d'opérations ordonnées : \textbf{comprendre} (lecture méthodique, distinction manifeste/latent), \textbf{sélectionner} (thèse, arguments, idées essentielles), \textbf{réduire} (suppression, substitution, fusion) et \textbf{reformuler} (phrases personnelles articulées par des connecteurs variés). Les élèves ont constaté qu'un bon résumé se reconnaît à quatre qualités : la \cle{fidélité} à la pensée de l'auteur, la \cle{cohérence} de l'enchaînement, la correction de la \cle{langue} et le respect des \cle{contraintes} imposées. L'exposé oral a en outre montré que savoir résumer, c'est aussi savoir transmettre : le rapporteur qui présente en 3 minutes le travail de son groupe exerce une compétence de communication directement utile dans la vie professionnelle d'un technicien (compte rendu de chantier, rapport de stage, présentation de projet). La grille d'évaluation croisée et la remédiation collective ont enfin développé l'esprit critique : évaluer la production d'autrui, c'est apprendre à relire et à améliorer la sienne, démarche décisive pour l'épreuve de contraction de texte au Probatoire et au Baccalauréat technique.

\newpage

\coursec{Méthodes et exercices résolus}

\coursec{Formulation des idées essentielles d'un texte à résumer}

\courssub{Le texte argumentatif : structure et procédés}

\begin{definition}[Texte argumentatif]
Un \cle{texte argumentatif} a pour but de convaincre ou de persuader un destinataire en défendant une \cle{thèse} à l'aide d'\cle{arguments} organisés selon une logique rigoureuse. Il s'appuie sur des procédés rhétoriques (exemples, analogies, citations, figures de style) et des connecteurs logiques pour assurer la cohérence du raisonnement.
\end{definition}

\begin{propriete}[Structure canonique d'un texte argumentatif]
Tout texte argumentatif bien construit comporte quatre moments identifiables :
\begin{enumerate}
\item \textbf{L'accroche (ou exorde)} : amorce le sujet, capte l'attention, pose le contexte.
\item \textbf{La thèse} : position claire que l'auteur défend (souvent en début ou fin de texte).
\item \textbf{Le développement} : enchaînement d'arguments (arguments d'autorité, logiques, pathétiques) illustrés par des exemples.
\item \textbf{La conclusion} : réaffirme la thèse, en tire les conséquences, ouvre parfois sur une perspective.
\end{enumerate}
\end{propriete}

\begin{methode}[Repérer la structure argumentative d'un texte]
Pour analyser la structure argumentative d'un texte avant de le résumer :
\begin{enumerate}
\item \textbf{Lire le texte deux fois} : une première lecture globale pour saisir le sujet, une seconde lecture attentive en soulignant les articulations logiques (connecteurs, paragraphes).
\item \textbf{Identifier la thèse} : chercher la position défendue par l'auteur. Se demander : « Que veut me faire admettre l'auteur ? »
\item \textbf{Repérer les arguments} : ce sont les raisons qui soutiennent la thèse. Ils sont souvent introduits par \cle{d'abord}, \cle{ensuite}, \cle{en effet}, \cle{car}, \cle{de plus}, \cle{par ailleurs}.
\item \textbf{Distinguer les exemples/illustrations} : ils concrétisent les arguments mais sont \cle{sacrificiels} dans un résumé (supprimés ou généralisés).
\item \textbf{Tracer le schéma} : Thèse $\rightarrow$ Argument 1 (+ illustration) $\rightarrow$ Argument 2 (+ illustration) $\rightarrow$ Conclusion. Vérifier la cohérence logique (déduction, induction, analogie).
\end{enumerate}
\end{methode}

\begin{aretenir}[Formule clé : Thèse, Argument, Exemple]
\begin{itemize}
\item \textbf{Thèse} = Ce que l'auteur affirme (son point de vue).
\item \textbf{Argument} = Pourquoi il l'affirme (la raison, la preuve logique).
\item \textbf{Exemple} = Preuve concrète, fait précis, anecdote (à éliminer ou condenser au résumé).
\end{itemize}
\end{aretenir}

\begin{exoresolu}[Analyse de la structure d'un paragraphe argumentatif]
Énoncé. Voici un extrait adapté : « L'école doit rester le pilier de la promotion sociale. En effet, elle offre à tous les enfants, quelle que soit leur origine, les mêmes chances de réussite. Ainsi, un enfant de village pauvre peut devenir médecin ou ingénieur grâce à ses études. De plus, l'école transmet les valeurs de travail et de discipline indispensables à la vie professionnelle. Il est donc indispensable de renforcer le système éducatif. » Dégagez la thèse, les arguments et l'exemple.
\tcblower
Solution détaillée.
\begin{enumerate}
\item \textbf{La thèse} : « L'école doit rester le pilier de la promotion sociale. » (Position défendue).
\item \textbf{Argument 1} : Introduit par \emph{en effet} : « elle offre à tous les enfants... les mêmes chances de réussite. » (Égalité des chances).
\item \textbf{L'exemple} : Introduit par \emph{ainsi} : « un enfant de village pauvre peut devenir médecin... » (Illustration concrète de l'argument 1 $\rightarrow$ \textbf{à supprimer au résumé}).
\item \textbf{Argument 2} : Introduit par \emph{de plus} : « l'école transmet les valeurs de travail et de discipline... » (Socialisation).
\item \textbf{La conclusion} : Marquée par \emph{donc} : « il est indispensable de renforcer le système éducatif. » (Reprise et extension de la thèse).
\end{enumerate}
\textbf{Schéma} : Thèse (pilier social) $\rightarrow$ Arg 1 (égalité chances) + Exemple $\rightarrow$ Arg 2 (valeurs) $\rightarrow$ Conclusion (renforcer l'école).
\end{exoresolu}

\courssub{Lecture méthodique n°4 et n°5 : Le Vieux nègre et la médaille}

\begin{methode}[Démarche de la lecture méthodique (Roman)]
L'étude d'une œuvre intégrale au programme suit trois temps :
\begin{enumerate}
\item \textbf{Découverte (Lecture cursive)} : Lecture personnelle à la maison ; repérage : auteur, contexte historique (Cameroun, années 50, fin colonisation), genre (roman satirique), personnages principaux (Meka, Kelara, Commandant, Père Vandermayer), intrigue.
\item \textbf{Analyse guidée (Lectures méthodiques)} : En classe, étude d'extraits choisis (lectures n°4 et n°5) pour approfondir : \textbf{thèmes} (hypocrisie coloniale, vanité des honneurs, conflit générations/tradition), \textbf{procédés} (ironie, satire, contraste, oralité), \textbf{portrait} (Meka : naïveté $\rightarrow$ lucidité), \textbf{style} (phrases courtes, dialogues, proverbes).
\item \textbf{Synthèse} : Rédaction de fiches de lecture (résumé, analyse personnages, thèmes, citation clé, avis critique).
\end{enumerate}
\end{methode}

\begin{exemplebox}[Repères pour les Lectures 4 et 5 (Extrait type : Remise de la médaille / Nuit de la perte)]
\textbf{Lecture 4 (Remise de la médaille)} : Focus sur l'ironie dramatique. Le discours du Commandant (contenu manifeste vs latent), la naïveté de Meka, le regard critique du narrateur. Travail sur le \cle{contenu latent} (voir Méthode 5).
\textbf{Lecture 5 (La perte de la médaille / Désillusion)} : Focus sur le basculement. La nuit, la recherche vaine, la prise de conscience de Meka (« La médaille est perdue... la France ne sait pas reconnaître »). Structure narrative : climax $\rightarrow$ chute $\rightarrow$ résolution amère. Préparation à la contraction (sélection idées essentielles).
\end{exemplebox}

\begin{exoresolu}[Analyse Manifeste / Latent sur un extrait de la Lecture 4]
Énoncé. Le Commandant dit : « Tu es un brave homme. La France sait reconnaître ses fidèles serviteurs. » (Extrait adapté). Distinguez contenu manifeste et latent.
\tcblower
Solution détaillée.
\begin{enumerate}
\item \textbf{Manifeste} : Éloge officiel ; promesse de reconnaissance institutionnelle pour services rendus.
\item \textbf{Indices du latent} : « Brave homme » (paternalisme, infantilisations) ; « Fidèles serviteurs » (statut de subalterne, non d'égal) ; Contexte : médaille de « mérite indigène » (honneur dérisoire).
\item \textbf{Latent} : L'auteur (Oyono) suggère le \cle{mépris paternaliste} et l'\cle{hypocrisie coloniale} : la domination se masque sous l'honneur. L'ironie naît du décalage entre la solennité du discours et la réalité de l'objet (médaille) et du statut (serviteur).
\end{enumerate}
\end{exoresolu}

\courssub{Les techniques de réduction du texte (Contraction)}

\begin{definition}[Contraction de texte]
La \cle{contraction} (ou résumé contracté) est une réécriture qui réduit un texte à un quart de sa longueur (généralement 25 \% $\pm$ 10 \%) en ne conservant que les \cle{idées essentielles} (thèse + arguments principaux), reformulées avec ses propres mots, dans un texte cohérent et autonome.
\end{definition}

\begin{methode}[Appliquer les techniques de réduction pas à pas]
\begin{enumerate}
\item \textbf{Compter et calculer} : Compter les mots du texte original ($N$). Cible $= N / 4$. Marge tolérée $\pm 10 \%$. Exemple : 400 mots $\rightarrow$ cible 100 mots (90 à 110 mots).
\item \textbf{Sélectionner (Élagage)} : \textbf{Supprimer} exemples, anecdotes, énumérations, répétitions, figures de style, adjectifs qualificatifs non essentiels, connecteurs de simple liaison (certains « et », « aussi »). \textbf{Garder} : Thèse, arguments majeurs, articulations logiques, conclusion.
\item \textbf{Généraliser} : Remplacer une énumération par un hyperonyme (« mangues, goyaves, papayes » $\rightarrow$ « fruits tropicaux »).
\item \textbf{Condenser (Nominalisation / Subordination)} : Transformer une proposition en groupe nominal (« parce qu'il était courageux » $\rightarrow$ « par son courage » / « grâce à son courage »), une relative en participe (« qui habitait au village » $\rightarrow$ « habitant au village »), une conjonction en préposition (« comme il pleuvait » $\rightarrow$ « à cause de la pluie »).
\item \textbf{Reformuler} : Dire avec son vocabulaire et sa syntaxe (changement de voix, de structure, synonymes). \textbf{Interdiction} : recopier une suite de 4 mots et plus du texte original.
\item \textbf{Relier et Vérifier} : Assurer la cohérence par des connecteurs logiques (voir Méthode 4). Compter les mots du brouillon final. Vérifier fidélité (pas d'ajout, pas d'opinion).
\end{enumerate}
\end{methode}

\begin{aretenir}[Check-list du bon résumé contracté]
\begin{itemize}
\item[$\square$] Longueur respectée (quart $\pm$ 10 \%).
\item[$\square$] Idées essentielles toutes présentes (thèse + arguments principaux).
\item[$\square$] Aucune idée ajoutée, aucune opinion personnelle.
\item[$\square$] Reformulation effective (vocabulaire + syntaxe changés).
\item[$\square$] Cohérence textuelle (connecteurs, temps verbaux homogènes, progression logique).
\item[$\square$] Orthographe, grammaire, ponctuation soignées.
\end{itemize}
\end{aretenir}

\begin{exoresolu}[Contraction d'un paragraphe (Entraînement Bac)]
Énoncé. Résumez en \textbf{25 mots environ} le passage suivant (\textbf{environ 75 mots}) : « Dans les villages, les jeunes désertent les champs. Ils partent vers les villes, attirés par les lumières, les boutiques, les cinémas, les bars, espérant trouver un emploi facile et une vie confortable. Mais la ville, loin de tenir ses promesses, les plonge souvent dans le chômage, la misère et la délinquance. »
\tcblower
Solution détaillée.
\begin{enumerate}
\item \textbf{Comptage} : Texte source $\approx$ 55 mots. Quart $\approx$ 14 mots. Consigne impose 25 mots $\rightarrow$ viser 23--27 mots.
\item \textbf{Idées essentielles} : (1) Exode rural des jeunes vers la ville. (2) Motif : attirance/espérance vie meilleure. (3) Réalité : échec (chômage, misère).
\item \textbf{Réduction} : Énumérations ville $\rightarrow$ « attraits urbains » ; liste maux $\rightarrow$ « désillusions » / « précarité ».
\item \textbf{Reformulation} : « Attirés par les attraits urbains où ils espèrent une vie meilleure, les jeunes quittent les campagnes pour y connaître souvent de graves désillusions. » (24 mots).
\item \textbf{Vérification} : 3 idées conservées $\checkmark$ ; Reformulation $\checkmark$ ; Cohérence $\checkmark$ ; Compte $\checkmark$.
\end{enumerate}
\end{exoresolu}

\courssub{Les outils de liaison : conjonctions de coordination et de subordination}

\begin{definition}[Connecteurs logiques]
Les \cle{conjonctions} sont des mots-outils invariables qui relient des unités syntaxiques (mots, groupes, propositions) en explicitant le \cle{rapport logique} (addition, opposition, cause, conséquence, but, condition, temps, concession). Leur maîtrise est indispensable pour articuler les idées du résumé.
\end{definition}

\begin{propriete}[Conjonctions de coordination]
Relient deux éléments de \textbf{même nature/fonction} (propositions indépendantes). Les 7 principales : \cle{mais}, \cle{ou}, \cle{et}, \cle{donc}, \cle{or}, \cle{ni}, \cle{car}.
\emph{Moyen mnémotechnique} : « \textbf{Mais où est donc Ornicar ?} »
\begin{itemize}
\item \textbf{Mais} : opposition/restriction.
\item \textbf{Ou} : alternative/choix.
\item \textbf{Et} : addition/enchaînement.
\item \textbf{Donc} : conséquence/déduction.
\item \textbf{Or} : opposition forte / justification (souvent en milieu de phrase).
\item \textbf{Ni} : addition négative (souvent corrélé : ni... ni...).
\item \textbf{Car} : cause/explication (jamais en début de phrase).
\end{itemize}
\end{propriete}

\begin{propriete}[Conjonctions de subordination]
Introduisent une \textbf{proposition subordonnée} (dépendante). Le choix impose souvent le \textbf{mode} (Indicatif / Subjonctif).
\begin{center}
\begin{tabularx}{\linewidth}{|l|X|l|}\hline
\textbf{Rapport} & \textbf{Conjonctions principales} & \textbf{Mode} \\ \hline
Cause & \cle{parce que}, \cle{puisque}, \cle{comme}, \cle{vu que} & Indicatif \\ \hline
Conséquence & \cle{si bien que}, \cle{de sorte que} & Indicatif \\ \hline
But & \cle{pour que}, \cle{afin que}, \cle{de peur que} & \textbf{Subjonctif} \\ \hline
Condition & \cle{si}, \cle{au cas où}, \cle{pourvu que} & Indicatif (si) / Subjonctif (pourvu que) \\ \hline
Temps & \cle{quand}, \cle{lorsque}, \cle{dès que}, \cle{avant que}, \cle{après que} & Ind. (quand, dès que) / Subj. (avant que) / Ind. (après que) \\ \hline
Concession & \cle{bien que}, \cle{quoique}, \cle{même si} & \textbf{Subjonctif} (bien que, quoique) / Ind. (même si) \\ \hline
Relatif / Complément & \cle{que} (COD), \cle{qui} (Sujet) & Indicatif \\ \hline
\end{tabularx}
\end{center}
\end{propriete}

\begin{attention}[Pièges Bac : Mode après conjonctions de subordination]
\begin{itemize}
\item \textbf{Subjonctif OBLIGATOIRE} après : \emph{pour que, afin que, de peur que, avant que, bien que, quoique, pourvu que}.
\item \textbf{Indicatif} après : \emph{parce que, puisque, comme, quand, lorsque, dès que, après que, si} (condition), \emph{même si}.
\item \textbf{Erreur fatale} : Écrire « \emph{bien qu'il est} » au lieu de « \emph{bien qu'il soit} » ; « \emph{pour qu'il a} » au lieu de « \emph{pour qu'il ait} ».
\end{itemize}
\end{attention}

\begin{exoresolu}[Choix et orthographe des conjonctions]
Énoncé. Reliez par la conjonction appropriée (attention au mode) : (a) Meka est vieux. Il marche. (Opposition) (b) Il a donné ses terres. Il veut la médaille. (But) (c) La médaille est perdue. Meka comprend. (Temps / Conséquence)
\tcblower
Solution détaillée.
\begin{enumerate}
\item (a) \textbf{Coordination} : « Meka est vieux, \emph{mais} il marche. » \textbf{Subordination (Concession)} : « \emph{Bien qu'il soit} vieux, Meka marche. » (Subjonctif \emph{soit}).
\item (b) \textbf{But} : « Il a donné ses terres \emph{pour que} la mission \emph{lui donne} la médaille. » (Subjonctif \emph{donne}). Ou infinitif : « \emph{Pour} obtenir la médaille, il a donné ses terres. »
\item (c) \textbf{Temps} : « \emph{Quand} la médaille est perdue, Meka comprend. » \textbf{Conséquence} : « La médaille est perdue, \emph{donc} Meka comprend. » / « \emph{Si bien que} Meka comprend. »
\end{enumerate}
\end{exoresolu}

\courssub{Contenus manifestes et contenus latents}

\begin{definition}[Niveaux de lecture]
\begin{itemize}
\item \textbf{Contenu \cle{manifeste} (Explicite)} : Ce qui est \textbf{dit textuellement}, le sens littéral, l'information de surface.
\item \textbf{Contenu \cle{latent} (Implicite)} : Ce qui est \textbf{suggéré, sous-entendu}, non dit directement. Il se déduit par l'analyse du contexte, du ton, des connotations, des figures de style (ironie, litote), des silences.
\end{itemize}
\end{definition}

\begin{methode}[Repérer et formuler le contenu latent]
\begin{enumerate}
\item \textbf{Énoncer le manifeste} : Résumer le sens littéral de l'énoncé.
\item \textbf{Traquer les indices} : Mots à double sens, connotations péjoratives/laudatives, ironie (dire le contraire), litote (dire moins pour dire plus), contraste forme/fond, non-dits, point de vue narratif.
\item \textbf{Formuler l'implicite} : Rédiger une phrase claire : « L'auteur suggère que... » / « Derrière cet éloge, se profile une critique de... »
\item \textbf{Justifier par le texte} : Citer le(s) mot(s) ou passage(s) précis qui fondent l'interprétation.
\item \textbf{Valider la cohérence} : Vérifier que le latent ne contredit pas le contexte global de l'œuvre.
\end{enumerate}
\end{methode}

\begin{exemplebox}[Application à \emph{Le Vieux nègre et la médaille} (Lecture 4)]
\textbf{Texte} : Le Commandant : « Tu es un brave homme. La France sait reconnaître ses fidèles serviteurs. »
\begin{itemize}
\item \textbf{Manifeste} : Félicitations officielles ; reconnaissance de la fidélité par la métropole.
\item \textbf{Indices latent} : « Brave homme » $\rightarrow$ condescendance ; « Fidèles serviteurs » $\rightarrow$ relation maître/serviteur, non citoyen/patrie ; Ironie narrative (Oyono) : la médaille sera perdue, la « reconnaissance » est vide.
\item \textbf{Latent} : \textbf{Dénonciation de l'hypocrisie coloniale} : l'administration instrumentalise l'honneur pour assurer la soumission. Le « service » rendu (don de terres) est spoliation ; la « récompense » est dérisoire.
\end{itemize}
\end{exemplebox}

\courssub{Reformulation, exposé et activités d'intégration}

\begin{methode}[Reformuler sans trahir : Procédure en 5 étapes]
\begin{enumerate}
\item \textbf{Comprendre} : Isoler l'idée nucléaire de la phrase (Sujet + Verbe + Complément essentiel). Ignorer les accessoires.
\item \textbf{Substituer le lexique} : Synonymes, antonymes (négation), hyperonymes (généralisation), termes techniques $\leftrightarrow$ termes courants.
\item \textbf{Restructurer la syntaxe} :
\begin{itemize}
\item Voix active $\leftrightarrow$ Passive / Pronominal.
\item Phrase verbale $\leftrightarrow$ Phrase nominale (nominalisation : « il décide » $\rightarrow$ « sa décision »).
\item Affirmation $\leftrightarrow$ Négation (litote : « il n'a pas refusé »).
\item Subordination $\leftrightarrow$ Coordination / Juxtaposition / Participes / Infinitifs.
\end{itemize}
\item \textbf{Conserver registre et nuance} : Ne pas passer du soutenu au familier ; ne pas amplifier (« très content » $\neq$ « ravi » si original « content ») ; ne pas affaiblir.
\item \textbf{Contrôler} : Relecture croisée Original vs Reformulation. \textbf{Règle d'or} : Aucune suite de 4 mots identiques.
\end{enumerate}
\end{methode}

\begin{exoresolu}[Reformulation de phrases du roman]
Énoncé. Reformulez : (a) « Le vieux Meka fut profondément ému par la médaille que lui remit le commandant. » (b) « Les villageois admirent celui que les Blancs ont honoré. »
\tcblower
Solution détaillée.
\begin{enumerate}
\item (a) \textbf{Idée} : Émotion de Meka causée par remise médaille par Commandant.
\textbf{Reformulations} : « La remise de la médaille par le commandant a profondément ému le vieux Meka. » (Passif $\rightarrow$ Actif + Nominalisation). « En recevant la médaille des mains du commandant, le vieux Meka a ressenti une vive émotion. » (Subordination temporelle + Synonymie \emph{ému/ressenti émotion}).
\item (b) \textbf{Idée} : Admiration villageois $\leftarrow$ Honneur Blancs.
\textbf{Reformulations} : « L'honneur reçu des Blancs vaut à Meka l'admiration des villageois. » (Nominalisation cause). « C'est parce que les Blancs l'ont honoré que les villageois l'admirent. » (Cleft sentence / Insistance).
\end{enumerate}
\end{exoresolu}

\begin{methode}[Réussir un exposé oral (5--10 min)]
\begin{enumerate}
\item \textbf{Préparation écrite} : Sujet précis $\rightarrow$ Plan détaillé (Intro / 2--3 Parties / Conclusion) $\rightarrow$ \textbf{Fiches bristol} : Mots-clés uniquement (jamais phrases entières), citations courtes, chiffres, articulation connecteurs.
\item \textbf{Introduction (1 min)} : Accroche $\rightarrow$ Sujet $\rightarrow$ Problématique $\rightarrow$ Annonce du plan.
\item \textbf{Développement (3--6 min)} : 2 ou 3 parties équilibrées. Chaque partie : Idée directrice $\rightarrow$ Argument(s) $\rightarrow$ Exemple/Preuve (texte) $\rightarrow$ Mini-transition. \textbf{Regarder l'auditoire}, voix posée, articulation.
\item \textbf{Conclusion (1 min)} : Bilan synthétique (réponse à la problématique) $\rightarrow$ Ouverture (question, autre œuvre, actualité).
\item \textbf{Questions} : Écoute, reformulation question, réponse structurée, honnêteté si ignorance.
\end{enumerate}
\end{methode}

\begin{exoresolu}[Plan d'exposé : « La satire coloniale dans Le Vieux nègre et la médaille »]
Énoncé. Construisez le plan détaillé (fiches) pour un exposé de 5 min.
\tcblower
Solution détaillée.
\textbf{Fiche 1 -- Introduction} : Oyono, 1956, Cameroun. Roman satire. Problématique : Comment la médaille devient symbole de la vanité coloniale ? Plan : 1. Procédés satiriques 2. Cibles 3. Portée.
\textbf{Fiche 2 -- I. Procédés} : Ironie (discours Commandant vs réalité) ; Contraste (apparence honneur / réalité spoliation) ; Naïveté Meka (révélateur absurdité système) ; Langage (proverbes, oralité, style direct).
\textbf{Fiche 3 -- II. Cibles} : Administration (Commandant : paternalisme) ; Église (Père Vandermayer : complicité) ; Villageois (crédulité, course honneurs) ; Meka lui-même (aliénation).
\textbf{Fiche 4 -- III. Portée} : Déconstruction mythe « mission civilisatrice » ; Éveil conscience (fin : lucidité Meka) ; Prémonition indépendances.
\textbf{Fiche 5 -- Conclusion} : Médaille = symbole vide. Satire = arme lucidité. Ouverture : actualité néocoloniale / autres œuvres (Betty, Mongo Beti).
\end{exoresolu}

\courssub{Activité d'intégration : De la lecture à la contraction}

\begin{exercice}[Activité intégrée type Bac -- Texte : Extrait Lecture 5 (Perte médaille)]
\textbf{Texte support (environ 120 mots)} : « Cette nuit-là, Meka ne dormit pas. Il chercha la médaille partout. Dans sa case, sur le chemin, au bord de la rivière. Il interrogea les passants, les enfants, les vieillards. Personne ne l'avait vue. Le soleil se leva. Meka comprit. La médaille était perdue à jamais. Avec elle, s'envolait l'honneur. Il avait tout donné : ses terres, sa paix, sa fierté. Pour un morceau de ferraille que le vent avait emporté. Il pleura, non pas sur l'objet, mais sur sa propre aveuglement. La France ne savait pas reconnaître. Elle ne savait que prendre. »
\begin{enumerate}
\item \textbf{Analyse} : Dégagez la thèse, les 2 arguments principaux et l'exemple/illustration de ce paragraphe narratif (traité comme texte argumentatif implicite).
\item \textbf{Contraction} : Résumez ce texte en \textbf{30 mots exactement} ($\pm$ 2 mots).
\item \textbf{Reformulation} : Reformulez cette phrase : « Il pleura, non pas sur l'objet, mais sur sa propre aveuglement. »
\item \textbf{Grammaire} : Reliez par une subordonnée de concession : « Meka a tout donné. Il n'a rien reçu. » (Utiliser \emph{bien que} + subjonctif).
\end{enumerate}
\end{exercice}

\begin{corrige}[Exercice intégration]
\begin{enumerate}
\item \textbf{Analyse structurelle} :
\begin{itemize}
\item \textbf{Thèse implicite} : La quête des honneurs coloniaux mène à la désillusion et à la perte de soi.
\item \textbf{Argument 1} : La recherche vaine de la médaille symbolise la vanité de l'honneur colonial (détails : case, chemin, rivière, interrogatoires).
\item \textbf{Argument 2} : Le bilan est une spoliation totale (terres, paix, fierté) pour un objet sans valeur (ferraille).
\item \textbf{Conclusion/Chute} : Lucidité tardive : « La France ne savait que prendre. » (Inversion du manifeste « reconnaître »).
\end{itemize}
\item \textbf{Contraction (30 mots $\pm$ 2)} : « Après une nuit de recherches vaines, Meka réalise la perte définitive de sa médaille. Il mesure alors l'étendue de son sacrifice (terres, paix) pour une ferraille vaine, découvrant l'illusion de la reconnaissance coloniale. » (31 mots -- OK).
\item \textbf{Reformulation} : « Ses larmes ne pleuraient pas l'objet perdu, mais sa propre cécité. » / « Ce n'est pas la médaille qu'il pleurait, mais son aveuglement. »
\item \textbf{Grammaire} : « \emph{Bien qu'il ait} tout donné, Meka n'a rien reçu. » (Subjonctif \emph{ait} obligatoire après \emph{bien que}).
\end{enumerate}
\end{corrige}

\courssub{Compte-rendu et remédiation}

\begin{aretenir}[Bilan des compétences visées (Leçon 21)]
\begin{itemize}
\item \textbf{Reception} : Maîtriser la lecture méthodique d'un roman (Oyono) : repérer manifeste/latent, ironie, structure narrative.
\item \textbf{Langue} : Maîtriser les conjonctions (coordination/subordination) et la concordance des modes (Indicatif/Subjonctif) pour écrire cohérent.
\item \textbf{Production} : Appliquer le protocole rigoureux de la contraction (Comptage $\rightarrow$ Sélection $\rightarrow$ Condensation $\rightarrow$ Reformulation $\rightarrow$ Vérification). Réussir un exposé oral structuré.
\end{itemize}
\end{aretenir}

\begin{attention}[Les 5 erreurs fatales au Probatoire / Bac Tech -- Résumé/Contraction]
\begin{enumerate}
\item \textbf{Le "Copier-Coller"} : Reprendre des segments du texte sans reformulation (sanction : note 0 ou pénalité lourde). \textbf{Parade} : Fermer le texte, écrire de mémoire ses idées, puis reformuler.
\item \textbf{L'Intrusion personnelle} : Ajouter « Je pense que... », « C'est vrai car... », des connaissances extérieures. \textbf{Parade} : S'effacer derrière l'auteur ; utiliser « L'auteur montre que... », « Le texte explique que... ».
\item \textbf{Le Non-respect du quantitatif} : Hors fourchette (quart $\pm$ 10 \%). \textbf{Parade} : Compter mots source $\rightarrow$ Calculer cible $\rightarrow$ Compter mots brouillon $\rightarrow$ Ajuster (condenser/développer) $\rightarrow$ Compter mots final.
\item \textbf{La Confusion Argument/Exemple} : Garder l'exemple (souvent long) et couper l'argument. \textbf{Parade} : Identifier la thèse $\rightarrow$ Lister arguments $\rightarrow$ Barrer exemples/énumérations $\rightarrow$ Généraliser si nécessaire.
\item \textbf{La Rupture de cohérence / Fautes de liaison} : Phrases juxtaposées sans logique, erreurs de mode (subjonctif/indicatif), anacoluthes. \textbf{Parade} : Relire à voix haute ; vérifier chaque connecteur ; réviser les déclencheurs de subjonctif (\emph{pour que, bien que, avant que, pourvu que}...).
\end{enumerate}
\end{attention}

\begin{savaistu}[Contexte camerounais : Ferdinand Oyono et la satire]
Ferdinand \textbf{Oyono} (1929--2010), figure majeure des lettres camerounaises, diplomate et homme d'État. \emph{Le Vieux nègre et la médaille} (1956) est le premier roman d'une trilogie (suivi de \emph{Une vie de boy}, \emph{Le Chemin d'Europe}). Publié aux Éditions de Minuit (collection « Le Chemin »), il a obtenu le prix Sainte-Beuve. L'œuvre illustre la \cle{littérature de la négritude} et la \cle{littérature engagée} pré-indépendance : elle utilise l'ironie et le rire pour déconstruire l'idéologie coloniale « civilisatrice » et révéler l'aliénation des « évolués ». Au programme du Bac Tech pour son accessibilité, sa portée historique et sa richesse stylistique (oralité, proverbes, style indirect libre).
\end{savaistu}

\newpage

\coursec{Exercices}

\courssub{Je vérifie mes connaissances}

\begin{exercice}[QCM : le texte argumentatif et le résumé]
Pour chaque question, entoure la bonne réponse.
\begin{enumerate}
\item La thèse d'un texte argumentatif est :
\begin{enumerate}
\item une anecdote personnelle de l'auteur ;
\item l'idée principale que l'auteur défend ;
\item la conclusion grammaticale du texte.
\end{enumerate}
\item Un résumé-contraction de texte doit :
\begin{enumerate}
\item contenir des exemples personnels ;
\item respecter le nombre de mots imposé et rester fidèle aux idées de l'auteur ;
\item critiquer les idées de l'auteur.
\end{enumerate}
\item Dans la phrase « Il pleut, donc je reste », le mot « donc » est :
\begin{enumerate}
\item une conjonction de subordination ;
\item une préposition ;
\item une conjonction de coordination.
\end{enumerate}
\item Le contenu latent d'un texte est :
\begin{enumerate}
\item ce qui est dit explicitement ;
\item ce qui est suggéré, sous-entendu ;
\item le titre du texte.
\end{enumerate}
\end{enumerate}
\end{exercice}

\begin{exercice}[Vrai ou faux justifié]
Réponds par vrai ou faux, puis justifie ta réponse en une phrase.
\begin{enumerate}
\item Dans un résumé, on peut citer une phrase du texte entre guillemets.
\item « Parce que » est une conjonction de subordination.
\item Dans \emph{Le Vieux nègre et la médaille} de Ferdinand Oyono, la médaille reçue par Meka symbolise véritablement la reconnaissance et l'égalité entre les hommes.
\item Supprimer les exemples et les répétitions est une technique de réduction du texte.
\end{enumerate}
\end{exercice}

\begin{exercice}[Définitions]
Définis en une ou deux phrases chacun des termes suivants : \cle{thèse}, \cle{argument}, \cle{contraction de texte}, \cle{reformulation}, \cle{contenu manifeste}.
\end{exercice}

\begin{exercice}[Reconnaissance directe des outils de liaison]
Dans les phrases suivantes, souligne la conjonction et précise s'il s'agit d'une conjonction de coordination ou de subordination.
\begin{enumerate}
\item Meka a reçu une médaille, mais il n'est pas devenu l'égal du colonisateur.
\item Les élèves utilisent les réseaux sociaux parce qu'ils veulent communiquer.
\item Tu résumes le texte ou tu recopies le texte ?
\item Bien qu'il soit fatigué, le vieillard traverse le village à pied.
\end{enumerate}
\end{exercice}

\courssub{Je m'entraîne}

\begin{exercice}[Classer les conjonctions]
Classe les douze conjonctions suivantes dans un tableau à deux colonnes : « coordination » et « subordination ».
\begin{center}
mais — parce que — ou — lorsque — et — puisque — donc — quoique — or — si — ni — car
\end{center}
\end{exercice}

\begin{exercice}[Transformer des phrases simples en phrases complexes]
Transforme chaque groupe de phrases simples en une seule phrase complexe, en utilisant la conjonction indiquée entre parenthèses, afin de réduire le paragraphe.
\begin{enumerate}
\item Les élèves passent beaucoup de temps sur les réseaux sociaux. Ils négligent leurs devoirs. (si bien que)
\item Meka est vieux. Meka est faible. Meka marche jusqu'à la ville. (bien que)
\item Le commandant donne une médaille. Il veut montrer sa générosité. (pour que / car)
\item La pluie tombe. La cérémonie continue. (quoique)
\item Tu lis le texte deux fois. Tu comprends les idées essentielles. (quand)
\end{enumerate}
\end{exercice}

\begin{exercice}[Contenu manifeste et contenu latent]
Lis l'extrait suivant, inspiré de la situation du \emph{Vieux nègre et la médaille} : « Meka, la médaille sur la poitrine, rentre sous la pluie. Ses pieds nus s'enfoncent dans la boue. Les chaussures offertes par les Blancs le blessent ; il les a laissées au village. »
\begin{enumerate}
\item Dans un tableau à deux colonnes, distingue trois éléments du \cle{contenu manifeste} (ce qui est dit) et trois éléments du \cle{contenu latent} (ce qui est suggéré).
\item Formule en une seule phrase l'idée essentielle de cet extrait.
\end{enumerate}
\end{exercice}

\begin{exercice}[Appliquer les techniques de réduction]
Voici un paragraphe de 80 mots : « Les réseaux sociaux, c'est-à-dire Facebook, WhatsApp, TikTok et bien d'autres applications du même genre, occupent une place vraiment très importante, une place énorme, dans la vie quotidienne des élèves camerounais d'aujourd'hui. Par exemple, un élève de Terminale peut passer jusqu'à cinq heures par jour sur son téléphone portable. »
\begin{enumerate}
\item Réduis ce paragraphe à environ 30 mots en appliquant au moins deux techniques de réduction (suppression des exemples, suppression des répétitions, substitution, généralisation).
\item Justifie chaque suppression ou substitution opérée.
\end{enumerate}
\end{exercice}

\courssub{Je me prépare au Bac}

\begin{exercice}[Contraction de texte : les réseaux sociaux]
Texte : « Aujourd'hui, dans nos lycées et collèges, les réseaux sociaux sont devenus les compagnons inséparables des élèves. Dès la sortie des classes, et parfois même pendant les cours, les écrans s'allument et les pouces se mettent à défiler. Ce phénomène, loin d'être anodin, mérite qu'on s'y arrête, car il comporte à la fois des avantages réels et des dangers certains.

Il faut d'abord reconnaître que les réseaux sociaux rendent de fiers services aux élèves. Ils facilitent la communication entre camarades de classe : un devoir oublié, une leçon incomprise, et le groupe WhatsApp de la classe vient à la rescousse. Ils donnent aussi accès à une masse considérable d'informations : des cours en ligne, des vidéos éducatives, des encyclopédies gratuites sont désormais à portée de main. Un élève de Terminale peut ainsi préparer son Baccalauréat en consultant des annales corrigées sans dépenser un franc.

Cependant, ces outils présentent des dangers qu'on ne saurait négliger. Le premier est la perte de temps : combien d'élèves sacrifient leurs heures de sommeil et d'étude à des vidéos sans intérêt ? Le second danger est la distraction en classe : l'élève dont l'esprit reste rivé sur son téléphone ne peut suivre attentivement les explications du professeur. S'y ajoutent les mauvaises rencontres et les fausses informations qui circulent librement.

En définitive, les réseaux sociaux sont un outil : tout dépend de l'usage qu'on en fait. Aux élèves de savoir s'en servir avec modération et discernement, afin qu'ils restent un moyen de réussite et non une cause d'échec. » (environ 280 mots)
\begin{enumerate}
\item Dégage la thèse défendue par l'auteur et les trois arguments principaux.
\item Rédige un résumé de ce texte en 70 mots (à 10 \% près), sans aucune citation, sans avis personnel.
\item Indique le nombre de mots de ton résumé à la fin.
\end{enumerate}
\end{exercice}

\begin{exercice}[Contraction et grammaire : la déforestation au Cameroun]
Texte : « La forêt camerounaise, l'une des plus riches d'Afrique, disparaît chaque année un peu plus. Les causes de cette déforestation sont multiples. D'abord, l'exploitation industrielle du bois, menée par de grandes sociétés, abat chaque année des milliers d'arbres centenaires. Ensuite, l'agriculture itinérante sur brûlis pousse les paysans à défricher toujours plus loin pour nourrir leurs familles. Enfin, la recherche du bois de chauffe, principale source d'énergie des ménages, achève de fragiliser les zones forestières proches des villes.

Les conséquences sont déjà visibles. Le climat se dérègle : les pluies deviennent irrégulières et les saisons sèches s'allongent. La faune s'appauvrit : éléphants, gorilles et chimpanzés perdent leur habitat naturel. Les sols, privés de la protection des arbres, sont emportés par l'érosion.

Pourtant, des solutions existent. Le reboisement, l'éducation des populations et l'utilisation de foyers améliorés peuvent inverser la tendance. Encore faut-il que les autorités et les citoyens agissent ensemble, avant qu'il ne soit trop tard. » (environ 180 mots)
\begin{enumerate}
\item Releve dans le texte deux conjonctions ou locutions conjonctives de subordination et deux connecteurs d'énumération.
\item Rédige un résumé du texte en 45 mots (à 10 \% près) en appliquant les quatre techniques de réduction.
\item Justifie par écrit deux suppressions et une substitution que tu as opérées.
\end{enumerate}
\end{exercice}

\begin{exercice}[Remédiation et exposé : corriger un résumé fautif]
Un élève a rédigé le résumé suivant (consigne : 50 mots maximum) : « Je trouve que ce texte est très intéressant et je suis d'accord avec l'auteur. Comme le dit le texte, "les réseaux sociaux sont les compagnons inséparables des élèves". Les réseaux sociaux aident les élèves à communiquer et à trouver des informations, mais ils font perdre du temps et distraient en classe. Il faut donc les utiliser avec modération. Personnellement, j'utilise WhatsApp tous les jours avec mes amis du quartier. » (74 mots)
\begin{enumerate}
\item Releve dans ce résumé quatre erreurs par rapport aux règles de la contraction de texte et corrige-les.
\item Réécris un résumé correct de 50 mots maximum.
\item Prépare un exposé oral de cinq minutes intitulé « Les étapes d'une contraction de texte réussie » : tu indiqueras ton plan en trois parties et le support visuel que tu utiliseras.
\end{enumerate}
\end{exercice}

\newpage

\coursec{Corrigés des exercices}

\begin{corrige}[Exercice 1 — QCM : le texte argumentatif et le résumé]
\begin{enumerate}
\item \textbf{b)} La thèse est l'idée principale que l'auteur défend. Une anecdote ne peut être qu'un exemple, et la conclusion n'est pas la thèse : elle la rappelle.
\item \textbf{b)} Le résumé-contraction doit respecter le nombre de mots imposé et rester fidèle aux idées de l'auteur. Les exemples personnels et la critique sont interdits : le résumé est neutre.
\item \textbf{c)} « Donc » est une conjonction de coordination (comme \emph{mais, ou, et, or, ni, car}). Il relie deux propositions de même niveau et exprime ici la conséquence.
\item \textbf{b)} Le contenu latent est ce qui est suggéré, sous-entendu. Ce qui est dit explicitement relève du contenu manifeste.
\end{enumerate}
\textbf{Point de vigilance :} ne pas confondre « donc » avec une conjonction de subordination : il n'introduit pas une proposition dépendante.
\end{corrige}

\begin{corrige}[Exercice 2 — Vrai ou faux justifié]
\begin{enumerate}
\item \textbf{Faux.} Le résumé doit être rédigé avec les mots \textbf{de l'élève} (ses propres mots) : aucune citation entre guillemets n'est admise, car cela reviendrait à recopier le texte.
\item \textbf{Vrai.} « Parce que » introduit une proposition subordonnée circonstancielle de cause, qui dépend de la proposition principale.
\item \textbf{Faux.} La médaille est un symbole trompeur : elle donne à Meka l'illusion de la reconnaissance, mais il reste méprisé et humilié par les colonisateurs ; c'est le contenu latent du roman.
\item \textbf{Vrai.} La suppression des exemples et des répétitions est l'une des quatre techniques de réduction, avec la substitution et la généralisation.
\end{enumerate}
\end{corrige}

\begin{corrige}[Exercice 3 — Définitions]
\begin{itemize}
\item \textbf{Thèse :} idée principale, point de vue que l'auteur d'un texte argumentatif défend et cherche à faire admettre au lecteur.
\item \textbf{Argument :} raison avancée pour appuyer la thèse et convaincre le lecteur ; il est souvent illustré par un exemple.
\item \textbf{Contraction de texte :} exercice qui consiste à réduire un texte à une fraction de sa longueur (souvent le quart) en n'en conservant que les idées essentielles, avec ses propres mots.
\item \textbf{Reformulation :} reprise d'une idée du texte dans d'autres termes, sans en changer le sens, pour éviter la citation et le recopiage.
\item \textbf{Contenu manifeste :} ensemble des informations explicitement énoncées dans le texte, ce qui est dit directement, par opposition au contenu latent (suggéré).
\end{itemize}
\end{corrige}

\begin{corrige}[Exercice 4 — Reconnaissance directe des outils de liaison]
\begin{enumerate}
\item « Meka a reçu une médaille, \textbf{mais} il n'est pas devenu l'égal du colonisateur. » → \textbf{mais} : conjonction de \cle{coordination} (opposition).
\item « Les élèves utilisent les réseaux sociaux \textbf{parce qu'}ils veulent communiquer. » → \textbf{parce que} : conjonction de \cle{subordination} (cause).
\item « Tu résumes le texte \textbf{ou} tu recopies le texte ? » → \textbf{ou} : conjonction de \cle{coordination} (alternative).
\item « \textbf{Bien qu'}il soit fatigué, le vieillard traverse le village à pied. » → \textbf{bien que} : conjonction de \cle{subordination} (concession ; attention, elle exige le subjonctif : « soit »).
\end{enumerate}
\textbf{Point de vigilance :} une conjonction de coordination relie deux unités de même niveau ; une conjonction de subordination introduit une proposition dépendante.
\end{corrige}

\begin{corrige}[Exercice 5 — Classer les conjonctions]
\begin{center}
\begin{tabularx}{\linewidth}{|X|X|}
\hline
\rowcolor{popblueL} \textbf{Coordination} & \textbf{Subordination} \\
\hline
mais & parce que \\
\hline
ou & lorsque \\
\hline
et & puisque \\
\hline
donc & quoique \\
\hline
or & si \\
\hline
ni & \\
\cline{1-1}
car & \\
\hline
\end{tabularx}
\end{center}
\textbf{Astuce :} les sept conjonctions de coordination se retiennent grâce à la phrase-mémoire « \emph{Mais où est donc Ornicar ?} » : \textbf{m}ais, \textbf{ou}, \textbf{e}st (et), \textbf{d}onc, \textbf{or}, \textbf{ni}, \textbf{car}. Toutes les autres conjonctions de la liste sont des conjonctions de subordination.
\end{corrige}

\begin{corrige}[Exercice 6 — Transformer des phrases simples en phrases complexes]
\begin{enumerate}
\item Les élèves passent beaucoup de temps sur les réseaux sociaux, \textbf{si bien qu'}ils négligent leurs devoirs. (conséquence)
\item \textbf{Bien qu'}il soit vieux et faible, Meka marche jusqu'à la ville. (concession ; « bien que » + subjonctif)
\item Le commandant donne une médaille \textbf{pour que} chacun admire sa générosité, \textbf{car} il veut montrer sa générosité. — Formulation la plus naturelle : « Le commandant donne une médaille \textbf{car} il veut montrer sa générosité. » (cause)
\item \textbf{Quoique} la pluie tombe, la cérémonie continue. (concession ; on peut aussi écrire « La cérémonie continue quoique la pluie tombe ».)
\item \textbf{Quand} tu lis le texte deux fois, tu comprends les idées essentielles. (temps)
\end{enumerate}
\textbf{Points de vigilance :} « bien que » et « quoique » exigent le \emph{subjonctif} ; la transformation ne doit pas changer le sens du paragraphe, seulement le condenser.
\end{corrige}

\begin{corrige}[Exercice 7 — Contenu manifeste et contenu latent]
\textbf{1.} Tableau de distinction :
\begin{center}
\begin{tabularx}{\linewidth}{|X|X|}
\hline
\rowcolor{popblueL} \textbf{Contenu manifeste (dit)} & \textbf{Contenu latent (suggéré)} \\
\hline
Meka porte une médaille sur la poitrine. & La médaille est un honneur trompeur, sans valeur réelle. \\
\hline
Il rentre sous la pluie, les pieds nus dans la boue. & Sa misère n'a pas changé malgré la décoration. \\
\hline
Les chaussures offertes le blessent ; il les a laissées au village. & Les « cadeaux » des Blancs sont inadaptés : le colonisateur ne comprend pas le colonisé. \\
\hline
\end{tabularx}
\end{center}
\textbf{2.} Idée essentielle : « Malgré les honneurs apparents que lui accorde le colonisateur, Meka demeure pauvre et humilié, ce qui révèle l'hypocrisie de la prétendue reconnaissance coloniale. »
\textbf{Point de vigilance :} le contenu latent doit rester fondé sur des indices du texte (pluie, boue, chaussures abandonnées), non sur une invention personnelle.
\end{corrige}

\begin{corrige}[Exercice 8 — Appliquer les techniques de réduction]
\textbf{1.} Résumé possible (30 mots) : « Les réseaux sociaux occupent une place très importante dans la vie quotidienne des élèves camerounais, qui peuvent y consacrer plusieurs heures par jour. » (27 mots — autre version acceptée : « Les réseaux sociaux occupent une place énorme dans la vie quotidienne des élèves camerounais, certains y passant jusqu'à cinq heures par jour. »)

\textbf{2.} Justification des opérations :
\begin{itemize}
\item \textbf{Suppression de l'exemple} : « c'est-à-dire Facebook, WhatsApp, TikTok et bien d'autres applications du même genre » est une énumération d'exemples qui illustre « réseaux sociaux » ; elle est supprimée sans perte d'idée.
\item \textbf{Suppression de la répétition} : « une place vraiment très importante, une place énorme » dit deux fois la même chose ; on ne garde qu'une seule formulation.
\item \textbf{Substitution / condensation} : « un élève de Terminale peut passer jusqu'à cinq heures par jour sur son téléphone portable » devient « qui peuvent y consacrer plusieurs heures par jour » : le pronom relatif « qui » évite de répéter « élèves » et la généralisation « plusieurs heures » remplace le cas particulier chiffré.
\end{itemize}
\textbf{Point de vigilance :} conserver l'idée chiffrée « cinq heures » est admis si le détail est jugé essentiel, mais la généralisation est préférable quand le nombre de mots est limité.
\end{corrige}

\begin{corrige}[Exercice 9 — Contraction de texte : les réseaux sociaux]
\textbf{1.} Thèse et arguments :
\begin{itemize}
\item \textbf{Thèse :} les réseaux sociaux sont un outil à double face pour les élèves : utiles mais dangereux ; tout dépend de l'usage qu'on en fait (il faut les utiliser avec modération et discernement).
\item \textbf{Argument 1 :} ils facilitent la communication entre camarades (entraide scolaire).
\item \textbf{Argument 2 :} ils donnent accès à une masse d'informations et de ressources éducatives gratuites.
\item \textbf{Argument 3 :} ils présentent des dangers : perte de temps, distraction en classe, mauvaises rencontres et fausses informations.
\end{itemize}

\textbf{2.} Résumé attendu (70 mots $\pm$ 10 \%, soit 63 à 77 mots) :

« Les réseaux sociaux sont devenus les compagnons quotidiens des élèves. Ils facilitent la communication entre camarades et donnent accès gratuitement à de nombreuses ressources éducatives utiles pour les études. Cependant, ils font perdre du temps, distraient en classe et exposent aux mauvaises rencontres et aux fausses informations. Ce sont donc de simples outils : les élèves doivent les utiliser avec modération et discernement pour réussir. » (66 mots)

\textbf{3.} Nombre de mots : 66 (compris dans la fourchette 63--77).

\textbf{Barème indicatif (sur 10) :} thèse et arguments correctement dégagés : 3 pts ; respect du nombre de mots : 2 pts ; fidélité au texte (aucune idée ajoutée, aucun avis personnel, aucune citation) : 2 pts ; reformulation et liaison des idées : 2 pts ; correction de la langue : 1 pt.

\textbf{Points de vigilance :} compter les mots avec soin et l'indiquer à la fin ; ne pas écrire « je » ni « à mon avis » ; ne pas recopier de phrase entière du texte ; conserver la progression thèse → avantages → dangers → conclusion.
\end{corrige}

\begin{corrige}[Exercice 10 — Contraction et grammaire : la déforestation au Cameroun]
\textbf{1.} Relevés :
\begin{itemize}
\item Conjonctions / locutions de subordination : « \textbf{Encore faut-il que} les autorités et les citoyens agissent ensemble » ; « \textbf{avant qu'}il ne soit trop tard ». (On accepte aussi « principale source d'énergie des ménages » non, ce n'en est pas une ; seuls ces deux relevés sont corrects.)
\item Connecteurs d'énumération : « \textbf{D'abord} », « \textbf{Ensuite} », « \textbf{Enfin} » (deux suffisent).
\end{itemize}

\textbf{2.} Résumé attendu (45 mots $\pm$ 10 \%, soit 40 à 50 mots) :

« Au Cameroun, la forêt disparaît à cause de l'exploitation industrielle, de l'agriculture sur brûlis et du bois de chauffe. Le climat se dérègle, la faune s'appauvrit et les sols s'érodent. Seuls le reboisement, l'éducation et les foyers améliorés, appliqués par tous ensemble, peuvent inverser cette tendance. » (46 mots)

\textbf{3.} Justification des opérations :
\begin{itemize}
\item \textbf{Suppression 1 :} « l'une des plus riches d'Afrique » est une précision élogieuse non indispensable à l'argumentation ; elle est supprimée.
\item \textbf{Suppression 2 :} l'énumération « éléphants, gorilles et chimpanzés » est un exemple de la faune menacée : on garde le terme générique « la faune ».
\item \textbf{Substitution :} « la recherche du bois de chauffe, principale source d'énergie des ménages » est remplacée par « le bois de chauffe » : l'apposition explicative est supprimée car elle ne fait que préciser le terme.
\end{itemize}

\textbf{Barème indicatif (sur 10) :} relevés grammaticaux exacts : 2 pts ; respect du nombre de mots : 2 pts ; fidélité et exhaustivité des idées essentielles (causes, conséquences, solutions) : 3 pts ; justifications des techniques : 2 pts ; langue : 1 pt.

\textbf{Points de vigilance :} ne pas confondre connecteur d'énumération (« d'abord ») et conjonction de subordination ; « avant que » exige le subjonctif (« soit ») ; le résumé doit conserver la structure causes → conséquences → solutions.
\end{corrige}

\begin{corrige}[Exercice 11 — Remédiation et exposé : corriger un résumé fautif]
\textbf{1.} Les quatre erreurs et leur correction :
\begin{enumerate}
\item \textbf{Avis personnel} : « Je trouve que ce texte est très intéressant et je suis d'accord avec l'auteur » → à supprimer : le résumé est neutre, on n'émet aucun jugement.
\item \textbf{Citation entre guillemets} : « "les réseaux sociaux sont les compagnons inséparables des élèves" » → à reformuler avec ses propres mots : « les réseaux sociaux accompagnent les élèves au quotidien ».
\item \textbf{Exemple personnel} : « Personnellement, j'utilise WhatsApp tous les jours avec mes amis du quartier » → à supprimer : aucun ajout étranger au texte n'est admis.
\item \textbf{Dépassement du nombre de mots} : 74 mots pour 50 maximum → il faut appliquer les techniques de réduction (suppressions, substitutions) pour atteindre la consigne.
\end{enumerate}

\textbf{2.} Résumé corrigé (50 mots maximum) :

« Les réseaux sociaux accompagnent les élèves au quotidien. Ils les aident à communiquer et à trouver des informations, mais ils font perdre du temps et distraient en classe. Il faut donc les utiliser avec modération et discernement pour qu'ils restent un moyen de réussite. » (45 mots)

\textbf{3.} Exposé oral « Les étapes d'une contraction de texte réussie » (5 minutes) :
\begin{itemize}
\item \textbf{Plan en trois parties :}
\begin{enumerate}
\item \emph{Lire et comprendre} : deux lectures, repérage du thème et de la thèse, soulignage des idées essentielles.
\item \emph{Réduire} : supprimer exemples, répétitions et détails ; substituer et généraliser ; reformuler avec ses propres mots.
\item \emph{Rédiger et vérifier} : rédiger un texte lié grâce aux connecteurs, compter les mots, vérifier la fidélité et la neutralité.
\end{enumerate}
\item \textbf{Support visuel :} un tableau ou un schéma en trois flèches (Lire → Réduire → Rédiger) affiché au tableau ou sur un carton, avec un exemple de phrase avant/après réduction.
\end{itemize}

\textbf{Barème indicatif (sur 10) :} identification des quatre erreurs : 4 pts ; résumé corrigé conforme (mots, neutralité, fidélité) : 3 pts ; plan d'exposé cohérent et support annoncé : 2 pts ; langue : 1 pt.

\textbf{Points de vigilance :} les erreurs à relever concernent les \emph{règles} de la contraction (neutralité, interdiction de citer, consigne de longueur, fidélité) ; à l'oral, annoncer le plan, parler lentement et s'appuyer sur le support sans le lire.
\end{corrige}

\newpage

\coursec{Fiche bilan}

\begin{popfigure}
\begin{tikzpicture}[
    mindmap,
    concept color=popblue,
    level 1 concept/.append style={level distance=4.5cm, sibling angle=60},
    level 2 concept/.append style={level distance=3cm, sibling angle=45},
    every node/.style={font=\small},
    text=white,
    grow cyclic
]
\node[concept, scale=1.2, align=center]{Formulation des idées\\ essentielles \& Résumé}
    child[concept color=poporange] {
        node[concept] {Texte argumentatif}
        child { node[concept, scale=0.7] {Structure (Thèse, Arguments)} }
        child { node[concept, scale=0.7] {Procédés rhétoriques} }
    }
    child[concept color=popgreen] {
        node[concept, align=center]{\emph{Le Vieux Nègre}\\ et la Médaille}
        child { node[concept, scale=0.7] {Lecture méthodique n°4} }
        child { node[concept, scale=0.7] {Lecture méthodique n°5} }
    }
    child[concept color=poppurple] {
        node[concept] {Techniques de réduction}
        child { node[concept, scale=0.7] {Sélection / Suppression} }
        child { node[concept, scale=0.7] {Généralisation} }
        child { node[concept, scale=0.7] {Reformulation / Fusion} }
    }
    child[concept color=poppink] {
        node[concept] {Outils de liaison}
        child { node[concept, scale=0.7] {Coordination (mais, or, et)} }
        child { node[concept, scale=0.7] {Subordination (que, si, comme)} }
        child { node[concept, scale=0.7] {Connecteurs logiques} }
    }
    child[concept color=popgold] {
        node[concept] {Contenus}
        child { node[concept, scale=0.7] {Manifestes (Explicites)} }
        child { node[concept, scale=0.7] {Latents (Implicites)} }
    }
    child[concept color=popdark] {
        node[concept] {Production \& Intégration}
        child { node[concept, scale=0.7] {Contraction / Résumé} }
        child { node[concept, scale=0.7] {Exposé oral} }
        child { node[concept, scale=0.7] {Activités intégrées} }
    };
\end{tikzpicture}
\legende{Carte mentale de la leçon 21 : Formulation des idées essentielles d'un texte à résumer}
\end{popfigure}

\begin{bilan}
\courssub{Définitions clés}
\begin{itemize}
    \item \textbf{Texte argumentatif} : Texte visant à convaincre ou persuader un destinataire d'une \cle{thèse} à l'aide d'\cle{arguments} structurés et de \cle{procédés rhétoriques}.
    \item \textbf{Contraction de texte (Résumé)} : Réduction significative (généralement 1/4 à 1/3 de l'original) conservant \emph{uniquement} les \cle{idées essentielles} (thèse + arguments principaux), reformulées en \cle{propres mots}, sans commentaire personnel.
    \item \textbf{Contenu manifeste} : Information explicite, directement énoncée dans le texte (déclarations, faits, arguments formulés).
    \item \textbf{Contenu latent} : Information implicite, sous-entendue, à inférer (valeurs, présupposés, intentions cachées, ironie).
    \item \textbf{Conjonctions de coordination} : Lient deux éléments de même nature/fonction (\emph{mais, or, et, donc, or, ni, car}).
    \item \textbf{Conjonctions de subordination} : Lient une proposition subordonnée à une principale ; marquent la relation logique (cause, conséquence, condition, concession, temps...).
\end{itemize}

\courssub{Méthode express : La contraction en 6 étapes}
\begin{enumerate}
    \item \textbf{Lecture active} : Identifier le \cle{sujet}, la \cle{thèse}, le \cle{ton}, le \cle{plan} (macrostructure). Repérer contenus manifestes/latents.
    \item \textbf{Découpage \& Sélection} : Segmenter en idées nucléaires. \textbf{Biffer} : exemples, détails, répétitions, digressions, adjectifs qualificatifs non essentiels.
    \item \textbf{Hiérarchisation} : Classer les idées retenues (Thèse $\rightarrow$ Arguments principaux $\rightarrow$ Conclusion).
    \item \textbf{Réduction linguistique} : 
    \begin{itemize}
        \item Généraliser (hyponymes $\rightarrow$ hyperonymes : \emph{chien, chat, oiseau} $\rightarrow$ \emph{animaux}).
        \item Fusionner phrases (nominalisation, participes, relatives).
        \item Remplacer périphrases par termes précis.
    \end{itemize}
    \item \textbf{Reformulation \& Liaison} : Réécrire avec son vocabulaire. Assurer la \cle{cohérence} par les \cle{connecteurs logiques} (coord./subord.) et la \cle{progression thématique}.
    \item \textbf{Contrôle final} : Compter les mots (respecter la contrainte $\pm 10\%$). Vérifier : Fidélité au sens ? Aucune idée personnelle ? Correction syntaxique/orthographique ? Fluidité ?
\end{enumerate}

\courssub{Tableau récapitulatif : Conjonctions \& Relations logiques}
\begin{center}
\begin{tabularx}{\linewidth}{|l|X|X|}\hline
\rowcolor{popblueL}
\textbf{Type} & \textbf{Exemples} & \textbf{Relations principales} \\ \hline
Coordination & mais, or, et, donc, or, ni, car & Opposition, addition, conclusion, cause \\ \hline
Subordination (subordonnants) & que, si, comme, lorsque, puisque, bien que, afin que & Complément, condition, cause, temps, concession, but \\ \hline
Adverbes / Locutions (Connecteurs) & cependant, par conséquent, en outre, c'est-à-dire & Opposition, conséquence, ajout, explication \\ \hline
\end{tabularx}
\end{center}

\courssub{Préparation de l'exposé oral (Intégration)}
\begin{itemize}
    \item \textbf{Introduction} : Accroche, sujet, thèse du texte, annonce du plan (2-3 parties).
    \item \textbf{Développement} : Une idée = un paragraphe oral. Appui sur notes (mots-clés), regard sur public, articulation claire (connecteurs).
    \item \textbf{Conclusion} : Synthèse des points forts, ouverture (lien avec actualité/autre texte), formule de politesse.
    \item \textbf{Critères} : Voix/Débit, Posture, Langage (soutenu, précis), Respect du temps, Structure logique.
\end{itemize}
\end{bilan}

\begin{aretenir}[Checklist avant le Bac -- Leçon 21]
\begin{itemize}
    \item $\square$ Je sais identifier la \textbf{thèse} et la structure \textbf{argumentative} d'un texte.
    \item $\square$ Je distingue \textbf{contenus manifestes} (explicites) et \textbf{contenus latents} (implicites).
    \item $\square$ Je maîtrise les \textbf{4 opérations de réduction} : suppression, généralisation, fusion, reformulation.
    \item $\square$ J'utilise correctement les \textbf{conjonctions de coordination} (mais, or, et, donc...) pour lier idées de même rang.
    \item $\square$ J'utilise correctement les \textbf{conjonctions de subordination} (que, si, bien que...) pour marquer les relations logiques complexes.
    \item $\square$ Je reformule les idées essentielles \textbf{avec mes propres mots} (pas de copier-coller).
    \item $\square$ Je respecte la \textbf{contrainte de longueur} imposée (nombre de mots $\pm 10\%$).
    \item $\square$ Mon résumé est \textbf{cohérent, autonome et fidèle} à la pensée de l'auteur.
    \item $\square$ Je sais structurer un \textbf{exposé oral} (Intro / Développement / Conclusion) à partir de mes notes.
    \item $\square$ J'ai révisé les \textbf{lectures méthodiques 4 et 5} de \emph{Le Vieux Nègre et la Médaille} (thèmes, style, argumentation).
    \item $\square$ Je repère les \textbf{procédés rhétoriques} (ironie, répétition, questions rhétoriques...) et leur effet.
    \item $\square$ Je gère mon temps : Lecture (15\%) -- Analyse/Sélection (30\%) -- Rédaction (45\%) -- Relecture (10\%).
\end{itemize}
\end{aretenir}

\findecours

\end{document}
